\documentclass[12pt]{article}
\usepackage[margin=1in]{geometry}
\usepackage{amsmath,amssymb,amsthm,booktabs,longtable,array}
\usepackage{newtxtext,newtxmath}
\usepackage[hidelinks]{hyperref}
\providecommand{\reg}[1]{\ifcsname p4num@#1\endcsname\csname p4num@#1\endcsname\else\textbf{??#1}\PackageWarning{registry}{Unknown key #1}\fi}
\expandafter\def\csname p4num@p4_exp_terms\endcsname{60}
\expandafter\def\csname p4num@p4_e_lo\endcsname{2.718281828459045}
\expandafter\def\csname p4num@p4_e_hi\endcsname{2.718281828459045}
\expandafter\def\csname p4num@u_B\endcsname{1.718282}
\expandafter\def\csname p4num@beta_B\endcsname{0.294655}
\expandafter\def\csname p4num@lambda_B\endcsname{0.336122}
\expandafter\def\csname p4num@b_B\endcsname{0.195615}
\expandafter\def\csname p4num@b0_B\endcsname{0.272648}
\expandafter\def\csname p4num@M_half_B\endcsname{0.731059}
\expandafter\def\csname p4num@m_half_B\endcsname{0.268941}
\expandafter\def\csname p4num@p4_input_h\endcsname{10.000000}
\expandafter\def\csname p4num@p4_input_ell\endcsname{1.000000}
\expandafter\def\csname p4num@p4_input_p\endcsname{0.500000}
\expandafter\def\csname p4num@p4_input_rho\endcsname{0.250000}
\expandafter\def\csname p4num@p4_input_cH\endcsname{6.000000}
\expandafter\def\csname p4num@p4_input_k\endcsname{0.020000}
\expandafter\def\csname p4num@p4_input_pi\endcsname{0.500000}
\expandafter\def\csname p4num@p4_input_r0\endcsname{1.200000}
\expandafter\def\csname p4num@p4_input_r1\endcsname{3.000000}
\expandafter\def\csname p4num@p4_input_r2\endcsname{3.600000}
\expandafter\def\csname p4num@p4_input_epsilon\endcsname{0.100000}
\expandafter\def\csname p4num@p4_input_low_upper\endcsname{0.200000}
\expandafter\def\csname p4num@p4_input_high_lower\endcsname{5.900000}
\expandafter\def\csname p4num@p4_input_high_upper\endcsname{6.100000}
\expandafter\def\csname p4num@p4_input_signal_a\endcsname{0.700000}
\expandafter\def\csname p4num@p4_input_signal_d\endcsname{0.750000}
\expandafter\def\csname p4num@p4_input_s_rho\endcsname{0.850000}
\expandafter\def\csname p4num@p4_input_s_cost\endcsname{7.140000}
\expandafter\def\csname p4num@p4_input_s_fee\endcsname{0.015000}
\expandafter\def\csname p4num@p4_input_s_r0\endcsname{1.100000}
\expandafter\def\csname p4num@p4_input_s_r1\endcsname{2.300000}
\expandafter\def\csname p4num@p4_input_valid_a\endcsname{0.950000}
\expandafter\def\csname p4num@p4_input_valid_d\endcsname{0.505000}
\expandafter\def\csname p4num@p4_input_comp_a\endcsname{0.940000}
\expandafter\def\csname p4num@p4_input_comp_d\endcsname{0.505000}
\expandafter\def\csname p4num@margin_p4_domain_beta_lower\endcsname{0.294655000}
\expandafter\def\csname p4num@margin_p4_domain_beta_upper\endcsname{0.038678333}
\expandafter\def\csname p4num@margin_p4_domain_lambda_lower\endcsname{0.336121769}
\expandafter\def\csname p4num@margin_p4_domain_lambda_upper\endcsname{0.663878231}
\expandafter\def\csname p4num@margin_p4_domain_pi_lower\endcsname{0.500000000}
\expandafter\def\csname p4num@margin_p4_domain_pi_upper\endcsname{0.500000000}
\expandafter\def\csname p4num@margin_p4_domain_rho_lower\endcsname{0.250000000}
\expandafter\def\csname p4num@margin_p4_domain_rho_upper\endcsname{0.750000000}
\expandafter\def\csname p4num@margin_p4_domain_reserve\endcsname{0.500000000}
\expandafter\def\csname p4num@margin_p4_domain_ell_reserve\endcsname{0.500000000}
\expandafter\def\csname p4num@margin_p4_domain_r0_ell\endcsname{0.200000000}
\expandafter\def\csname p4num@margin_p4_domain_r1_r0\endcsname{1.800000000}
\expandafter\def\csname p4num@margin_p4_domain_r2_r1\endcsname{0.600000000}
\expandafter\def\csname p4num@margin_p4_domain_h_r2\endcsname{6.400000000}
\expandafter\def\csname p4num@margin_p4_domain_cost\endcsname{6.000000000}
\expandafter\def\csname p4num@margin_p4_domain_fee\endcsname{0.020000000}
\expandafter\def\csname p4num@margin_p4_domain_epsilon\endcsname{0.100000000}
\expandafter\def\csname p4num@margin_p4_domain_band_width\endcsname{0.200000000}
\expandafter\def\csname p4num@margin_p4_domain_band_lower\endcsname{5.900000000}
\expandafter\def\csname p4num@margin_p4_domain_s_rho_lower\endcsname{0.850000000}
\expandafter\def\csname p4num@margin_p4_domain_s_rho_upper\endcsname{0.150000000}
\expandafter\def\csname p4num@margin_p4_domain_s_r0_ell\endcsname{0.100000000}
\expandafter\def\csname p4num@margin_p4_domain_s_r1_r0\endcsname{1.200000000}
\expandafter\def\csname p4num@margin_p4_domain_s_h_r1\endcsname{7.700000000}
\expandafter\def\csname p4num@margin_p4_domain_s_cost\endcsname{7.140000000}
\expandafter\def\csname p4num@margin_p4_domain_s_fee\endcsname{0.015000000}
\expandafter\def\csname p4num@margin_p4_domain_signal_a_lower\endcsname{0.200000000}
\expandafter\def\csname p4num@margin_p4_domain_signal_a_upper\endcsname{0.300000000}
\expandafter\def\csname p4num@margin_p4_domain_signal_d_lower\endcsname{0.250000000}
\expandafter\def\csname p4num@margin_p4_domain_signal_d_upper\endcsname{0.250000000}
\expandafter\def\csname p4num@margin_p4_domain_signal_d_a\endcsname{0.050000000}
\expandafter\def\csname p4num@margin_p4_domain_valid_a_lower\endcsname{0.450000000}
\expandafter\def\csname p4num@margin_p4_domain_valid_a_upper\endcsname{0.050000000}
\expandafter\def\csname p4num@margin_p4_domain_valid_d_lower\endcsname{0.005000000}
\expandafter\def\csname p4num@margin_p4_domain_valid_d_upper\endcsname{0.495000000}
\expandafter\def\csname p4num@margin_p4_domain_comp_a_lower\endcsname{0.440000000}
\expandafter\def\csname p4num@margin_p4_domain_comp_a_upper\endcsname{0.060000000}
\expandafter\def\csname p4num@margin_p4_domain_comp_d_lower\endcsname{0.005000000}
\expandafter\def\csname p4num@margin_p4_domain_comp_d_upper\endcsname{0.495000000}
\expandafter\def\csname p4num@p4_r0_t0\endcsname{0.291667}
\expandafter\def\csname p4num@p4_r0_th\endcsname{0.704167}
\expandafter\def\csname p4num@p4_r0_tl\endcsname{0.687500}
\expandafter\def\csname p4num@p4_r0_gh\endcsname{9.295833}
\expandafter\def\csname p4num@p4_r0_gl\endcsname{0.312500}
\expandafter\def\csname p4num@p4_r0_D\endcsname{0.016667}
\expandafter\def\csname p4num@p4_r0_wh\endcsname{0.412500}
\expandafter\def\csname p4num@p4_r0_wl\endcsname{0.395833}
\expandafter\def\csname p4num@tau_r0\endcsname{0.633117}
\expandafter\def\csname p4num@margin_p4_gross_gap_r0\endcsname{8.983333333}
\expandafter\def\csname p4num@p4_r1_t0\endcsname{0.416667}
\expandafter\def\csname p4num@p4_r1_th\endcsname{1.541667}
\expandafter\def\csname p4num@p4_r1_tl\endcsname{0.875000}
\expandafter\def\csname p4num@p4_r1_gh\endcsname{8.458333}
\expandafter\def\csname p4num@p4_r1_gl\endcsname{0.125000}
\expandafter\def\csname p4num@p4_r1_D\endcsname{0.666667}
\expandafter\def\csname p4num@p4_r1_wh\endcsname{1.125000}
\expandafter\def\csname p4num@p4_r1_wl\endcsname{0.458333}
\expandafter\def\csname p4num@tau_r1\endcsname{0.705000}
\expandafter\def\csname p4num@margin_p4_gross_gap_r1\endcsname{8.333333333}
\expandafter\def\csname p4num@p4_r2_t0\endcsname{0.430556}
\expandafter\def\csname p4num@p4_r2_th\endcsname{1.834722}
\expandafter\def\csname p4num@p4_r2_tl\endcsname{0.895833}
\expandafter\def\csname p4num@p4_r2_gh\endcsname{8.165278}
\expandafter\def\csname p4num@p4_r2_gl\endcsname{0.104167}
\expandafter\def\csname p4num@p4_r2_D\endcsname{0.938889}
\expandafter\def\csname p4num@p4_r2_wh\endcsname{1.404167}
\expandafter\def\csname p4num@p4_r2_wl\endcsname{0.465278}
\expandafter\def\csname p4num@tau_r2\endcsname{0.731392}
\expandafter\def\csname p4num@margin_p4_gross_gap_r2\endcsname{8.061111111}
\expandafter\def\csname p4num@margin_p4_weak_trade_B\endcsname{0.003333333}
\expandafter\def\csname p4num@margin_p4_strong_trade_B\endcsname{0.024823570}
\expandafter\def\csname p4num@margin_p4_very_strong_trade_B\endcsname{0.043126528}
\expandafter\def\csname p4num@margin_p4_weak_prior_B\endcsname{1.195833333}
\expandafter\def\csname p4num@margin_p4_strong_top_B\endcsname{0.217154822}
\expandafter\def\csname p4num@margin_p4_return_B\endcsname{0.002688902}
\expandafter\def\csname p4num@margin_p4_weak_top_B\endcsname{0.879842898}
\expandafter\def\csname p4num@margin_p4_strong_prior_B\endcsname{1.708333333}
\expandafter\def\csname p4num@margin_p4_atomless_floor_B\endcsname{2.166178511}
\expandafter\def\csname p4num@margin_p4_atomless_weak_B\endcsname{1.095833333}
\expandafter\def\csname p4num@margin_p4_atomless_strong_B\endcsname{0.117154822}
\expandafter\def\csname p4num@margin_p4_atomless_return_B\endcsname{-0.097311098}
\expandafter\def\csname p4num@margin_p4_atomless_floor_r2_B\endcsname{2.072133347}
\expandafter\def\csname p4num@margin_p4_atomless_top_upper_r2_B\endcsname{0.102688902}
\expandafter\def\csname p4num@margin_p4_atomless_top_lower_r2_B\endcsname{0.097311098}
\expandafter\def\csname p4num@margin_p4_ownership_rise_B\endcsname{0.398802607}
\expandafter\def\csname p4num@p4_base_r0_buy_rate\endcsname{0.000000}
\expandafter\def\csname p4num@p4_base_r0_sell_rate\endcsname{0.000000}
\expandafter\def\csname p4num@p4_base_r0_sell_mu\endcsname{0.500000}
\expandafter\def\csname p4num@p4_base_r0_sell_nu\endcsname{0.500000}
\expandafter\def\csname p4num@p4_base_r0_sell_price\endcsname{0.392708}
\expandafter\def\csname p4num@p4_base_r0_sell_eh\endcsname{0.250000}
\expandafter\def\csname p4num@p4_base_r0_sell_el\endcsname{0.250000}
\expandafter\def\csname p4num@p4_base_r0_sell_yp_entry_gap\endcsname{-1.195833333}
\expandafter\def\csname p4num@p4_base_r0_sell_ym_entry_gap\endcsname{-1.195833333}
\expandafter\def\csname p4num@p4_base_r0_sell_tp_payoff\endcsname{-0.022083333}
\expandafter\def\csname p4num@p4_base_r0_sell_tm_payoff\endcsname{-0.017916667}
\expandafter\def\csname p4num@p4_base_r0_zero_mu\endcsname{0.500000}
\expandafter\def\csname p4num@p4_base_r0_zero_nu\endcsname{0.500000}
\expandafter\def\csname p4num@p4_base_r0_zero_price\endcsname{0.392708}
\expandafter\def\csname p4num@p4_base_r0_zero_eh\endcsname{0.250000}
\expandafter\def\csname p4num@p4_base_r0_zero_el\endcsname{0.250000}
\expandafter\def\csname p4num@p4_base_r0_zero_yp_entry_gap\endcsname{-1.195833333}
\expandafter\def\csname p4num@p4_base_r0_zero_ym_entry_gap\endcsname{-1.195833333}
\expandafter\def\csname p4num@p4_base_r0_zero_tp_payoff\endcsname{0.000000000}
\expandafter\def\csname p4num@p4_base_r0_zero_tm_payoff\endcsname{0.000000000}
\expandafter\def\csname p4num@p4_base_r0_buy_mu\endcsname{0.500000}
\expandafter\def\csname p4num@p4_base_r0_buy_nu\endcsname{0.500000}
\expandafter\def\csname p4num@p4_base_r0_buy_price\endcsname{0.392708}
\expandafter\def\csname p4num@p4_base_r0_buy_eh\endcsname{0.250000}
\expandafter\def\csname p4num@p4_base_r0_buy_el\endcsname{0.250000}
\expandafter\def\csname p4num@p4_base_r0_buy_yp_entry_gap\endcsname{-1.195833333}
\expandafter\def\csname p4num@p4_base_r0_buy_ym_entry_gap\endcsname{-1.195833333}
\expandafter\def\csname p4num@p4_base_r0_buy_tp_payoff\endcsname{-0.017916667}
\expandafter\def\csname p4num@p4_base_r0_buy_tm_payoff\endcsname{-0.022083333}
\expandafter\def\csname p4num@p4_base_r0_entry\endcsname{0.250000}
\expandafter\def\csname p4num@p4_base_r0_own\endcsname{0.125000}
\expandafter\def\csname p4num@p4_base_r0_proceeds\endcsname{0.392708}
\expandafter\def\csname p4num@entry_r0\endcsname{0.250000}
\expandafter\def\csname p4num@own_r0\endcsname{0.125000}
\expandafter\def\csname p4num@p4_base_r1_buy_rate\endcsname{1.000000}
\expandafter\def\csname p4num@p4_base_r1_sell_rate\endcsname{1.000000}
\expandafter\def\csname p4num@p4_base_r1_sell_mu\endcsname{0.268941}
\expandafter\def\csname p4num@p4_base_r1_sell_nu\endcsname{0.268941}
\expandafter\def\csname p4num@p4_base_r1_sell_price\endcsname{0.576074}
\expandafter\def\csname p4num@p4_base_r1_sell_eh\endcsname{0.250000}
\expandafter\def\csname p4num@p4_base_r1_sell_el\endcsname{0.250000}
\expandafter\def\csname p4num@p4_base_r1_sell_yp_entry_gap\endcsname{-3.633821489}
\expandafter\def\csname p4num@p4_base_r1_sell_ym_entry_gap\endcsname{-3.633821489}
\expandafter\def\csname p4num@p4_base_r1_sell_tp_payoff\endcsname{-0.141843096}
\expandafter\def\csname p4num@p4_base_r1_sell_tm_payoff\endcsname{0.024823570}
\expandafter\def\csname p4num@p4_base_r1_zero_mu\endcsname{0.500000}
\expandafter\def\csname p4num@p4_base_r1_zero_nu\endcsname{0.500000}
\expandafter\def\csname p4num@p4_base_r1_zero_price\endcsname{0.614583}
\expandafter\def\csname p4num@p4_base_r1_zero_eh\endcsname{0.250000}
\expandafter\def\csname p4num@p4_base_r1_zero_el\endcsname{0.250000}
\expandafter\def\csname p4num@p4_base_r1_zero_yp_entry_gap\endcsname{-1.708333333}
\expandafter\def\csname p4num@p4_base_r1_zero_ym_entry_gap\endcsname{-1.708333333}
\expandafter\def\csname p4num@p4_base_r1_zero_tp_payoff\endcsname{0.000000000}
\expandafter\def\csname p4num@p4_base_r1_zero_tm_payoff\endcsname{0.000000000}
\expandafter\def\csname p4num@p4_base_r1_buy_mu\endcsname{0.731059}
\expandafter\def\csname p4num@p4_base_r1_buy_nu\endcsname{0.731059}
\expandafter\def\csname p4num@p4_base_r1_buy_price\endcsname{1.362372}
\expandafter\def\csname p4num@p4_base_r1_buy_eh\endcsname{1.000000}
\expandafter\def\csname p4num@p4_base_r1_buy_el\endcsname{1.000000}
\expandafter\def\csname p4num@p4_base_r1_buy_yp_entry_gap\endcsname{0.217154822}
\expandafter\def\csname p4num@p4_base_r1_buy_ym_entry_gap\endcsname{0.217154822}
\expandafter\def\csname p4num@p4_base_r1_buy_tp_payoff\endcsname{0.159294281}
\expandafter\def\csname p4num@p4_base_r1_buy_tm_payoff\endcsname{-0.507372386}
\expandafter\def\csname p4num@p4_base_r1_entry\endcsname{0.522757}
\expandafter\def\csname p4num@p4_base_r1_own\endcsname{0.324401}
\expandafter\def\csname p4num@p4_base_r1_proceeds\endcsname{0.872531}
\expandafter\def\csname p4num@entry_r1_B\endcsname{0.522757}
\expandafter\def\csname p4num@own_r1_B\endcsname{0.324401}
\expandafter\def\csname p4num@p4_base_r2_buy_rate\endcsname{1.000000}
\expandafter\def\csname p4num@p4_base_r2_sell_rate\endcsname{1.000000}
\expandafter\def\csname p4num@p4_base_r2_sell_mu\endcsname{0.268941}
\expandafter\def\csname p4num@p4_base_r2_sell_nu\endcsname{0.268941}
\expandafter\def\csname p4num@p4_base_r2_sell_price\endcsname{0.610002}
\expandafter\def\csname p4num@p4_base_r2_sell_eh\endcsname{0.250000}
\expandafter\def\csname p4num@p4_base_r2_sell_el\endcsname{0.250000}
\expandafter\def\csname p4num@p4_base_r2_sell_yp_entry_gap\endcsname{-3.727866653}
\expandafter\def\csname p4num@p4_base_r2_sell_ym_entry_gap\endcsname{-3.727866653}
\expandafter\def\csname p4num@p4_base_r2_sell_tp_payoff\endcsname{-0.191595694}
\expandafter\def\csname p4num@p4_base_r2_sell_tm_payoff\endcsname{0.043126528}
\expandafter\def\csname p4num@p4_base_r2_zero_mu\endcsname{0.500000}
\expandafter\def\csname p4num@p4_base_r2_zero_nu\endcsname{0.500000}
\expandafter\def\csname p4num@p4_base_r2_zero_price\endcsname{0.664236}
\expandafter\def\csname p4num@p4_base_r2_zero_eh\endcsname{0.250000}
\expandafter\def\csname p4num@p4_base_r2_zero_el\endcsname{0.250000}
\expandafter\def\csname p4num@p4_base_r2_zero_yp_entry_gap\endcsname{-1.865277778}
\expandafter\def\csname p4num@p4_base_r2_zero_ym_entry_gap\endcsname{-1.865277778}
\expandafter\def\csname p4num@p4_base_r2_zero_tp_payoff\endcsname{0.000000000}
\expandafter\def\csname p4num@p4_base_r2_zero_tm_payoff\endcsname{0.000000000}
\expandafter\def\csname p4num@p4_base_r2_buy_mu\endcsname{0.731059}
\expandafter\def\csname p4num@p4_base_r2_buy_nu\endcsname{0.731059}
\expandafter\def\csname p4num@p4_base_r2_buy_price\endcsname{0.718471}
\expandafter\def\csname p4num@p4_base_r2_buy_eh\endcsname{0.250000}
\expandafter\def\csname p4num@p4_base_r2_buy_el\endcsname{0.250000}
\expandafter\def\csname p4num@p4_base_r2_buy_yp_entry_gap\endcsname{-0.002688902}
\expandafter\def\csname p4num@p4_base_r2_buy_ym_entry_gap\endcsname{-0.002688902}
\expandafter\def\csname p4num@p4_base_r2_buy_tp_payoff\endcsname{0.043126528}
\expandafter\def\csname p4num@p4_base_r2_buy_tm_payoff\endcsname{-0.191595694}
\expandafter\def\csname p4num@p4_base_r2_entry\endcsname{0.250000}
\expandafter\def\csname p4num@p4_base_r2_own\endcsname{0.125000}
\expandafter\def\csname p4num@p4_base_r2_proceeds\endcsname{0.664236}
\expandafter\def\csname p4num@entry_r2\endcsname{0.250000}
\expandafter\def\csname p4num@own_r2_B\endcsname{0.125000}
\expandafter\def\csname p4num@p4_frozen_r0_buy_rate\endcsname{1.000000}
\expandafter\def\csname p4num@p4_frozen_r0_sell_rate\endcsname{1.000000}
\expandafter\def\csname p4num@p4_frozen_r0_sell_mu\endcsname{0.268941}
\expandafter\def\csname p4num@p4_frozen_r0_sell_nu\endcsname{0.268941}
\expandafter\def\csname p4num@p4_frozen_r0_sell_price\endcsname{0.391746}
\expandafter\def\csname p4num@p4_frozen_r0_sell_eh\endcsname{0.250000}
\expandafter\def\csname p4num@p4_frozen_r0_sell_el\endcsname{0.250000}
\expandafter\def\csname p4num@p4_frozen_r0_sell_yp_entry_gap\endcsname{-3.271509565}
\expandafter\def\csname p4num@p4_frozen_r0_sell_ym_entry_gap\endcsname{-3.271509565}
\expandafter\def\csname p4num@p4_frozen_r0_sell_tp_payoff\endcsname{-0.023046077}
\expandafter\def\csname p4num@p4_frozen_r0_sell_tm_payoff\endcsname{-0.018879411}
\expandafter\def\csname p4num@p4_frozen_r0_zero_mu\endcsname{0.500000}
\expandafter\def\csname p4num@p4_frozen_r0_zero_nu\endcsname{0.500000}
\expandafter\def\csname p4num@p4_frozen_r0_zero_price\endcsname{0.392708}
\expandafter\def\csname p4num@p4_frozen_r0_zero_eh\endcsname{0.250000}
\expandafter\def\csname p4num@p4_frozen_r0_zero_el\endcsname{0.250000}
\expandafter\def\csname p4num@p4_frozen_r0_zero_yp_entry_gap\endcsname{-1.195833333}
\expandafter\def\csname p4num@p4_frozen_r0_zero_ym_entry_gap\endcsname{-1.195833333}
\expandafter\def\csname p4num@p4_frozen_r0_zero_tp_payoff\endcsname{0.000000000}
\expandafter\def\csname p4num@p4_frozen_r0_zero_tm_payoff\endcsname{0.000000000}
\expandafter\def\csname p4num@p4_frozen_r0_buy_mu\endcsname{0.731059}
\expandafter\def\csname p4num@p4_frozen_r0_buy_nu\endcsname{0.731059}
\expandafter\def\csname p4num@p4_frozen_r0_buy_price\endcsname{0.699684}
\expandafter\def\csname p4num@p4_frozen_r0_buy_eh\endcsname{1.000000}
\expandafter\def\csname p4num@p4_frozen_r0_buy_el\endcsname{1.000000}
\expandafter\def\csname p4num@p4_frozen_r0_buy_yp_entry_gap\endcsname{0.879842898}
\expandafter\def\csname p4num@p4_frozen_r0_buy_ym_entry_gap\endcsname{0.879842898}
\expandafter\def\csname p4num@p4_frozen_r0_buy_tp_payoff\endcsname{-0.015517643}
\expandafter\def\csname p4num@p4_frozen_r0_buy_tm_payoff\endcsname{-0.032184310}
\expandafter\def\csname p4num@p4_frozen_r0_entry\endcsname{0.522757}
\expandafter\def\csname p4num@p4_frozen_r0_own\endcsname{0.324401}
\expandafter\def\csname p4num@p4_frozen_r0_proceeds\endcsname{0.503998}
\expandafter\def\csname p4num@p4_frozen_r1_buy_rate\endcsname{1.000000}
\expandafter\def\csname p4num@p4_frozen_r1_sell_rate\endcsname{1.000000}
\expandafter\def\csname p4num@p4_frozen_r1_sell_mu\endcsname{0.268941}
\expandafter\def\csname p4num@p4_frozen_r1_sell_nu\endcsname{0.268941}
\expandafter\def\csname p4num@p4_frozen_r1_sell_price\endcsname{0.576074}
\expandafter\def\csname p4num@p4_frozen_r1_sell_eh\endcsname{0.250000}
\expandafter\def\csname p4num@p4_frozen_r1_sell_el\endcsname{0.250000}
\expandafter\def\csname p4num@p4_frozen_r1_sell_yp_entry_gap\endcsname{-3.633821489}
\expandafter\def\csname p4num@p4_frozen_r1_sell_ym_entry_gap\endcsname{-3.633821489}
\expandafter\def\csname p4num@p4_frozen_r1_sell_tp_payoff\endcsname{-0.141843096}
\expandafter\def\csname p4num@p4_frozen_r1_sell_tm_payoff\endcsname{0.024823570}
\expandafter\def\csname p4num@p4_frozen_r1_zero_mu\endcsname{0.500000}
\expandafter\def\csname p4num@p4_frozen_r1_zero_nu\endcsname{0.500000}
\expandafter\def\csname p4num@p4_frozen_r1_zero_price\endcsname{0.614583}
\expandafter\def\csname p4num@p4_frozen_r1_zero_eh\endcsname{0.250000}
\expandafter\def\csname p4num@p4_frozen_r1_zero_el\endcsname{0.250000}
\expandafter\def\csname p4num@p4_frozen_r1_zero_yp_entry_gap\endcsname{-1.708333333}
\expandafter\def\csname p4num@p4_frozen_r1_zero_ym_entry_gap\endcsname{-1.708333333}
\expandafter\def\csname p4num@p4_frozen_r1_zero_tp_payoff\endcsname{0.000000000}
\expandafter\def\csname p4num@p4_frozen_r1_zero_tm_payoff\endcsname{0.000000000}
\expandafter\def\csname p4num@p4_frozen_r1_buy_mu\endcsname{0.731059}
\expandafter\def\csname p4num@p4_frozen_r1_buy_nu\endcsname{0.731059}
\expandafter\def\csname p4num@p4_frozen_r1_buy_price\endcsname{1.362372}
\expandafter\def\csname p4num@p4_frozen_r1_buy_eh\endcsname{1.000000}
\expandafter\def\csname p4num@p4_frozen_r1_buy_el\endcsname{1.000000}
\expandafter\def\csname p4num@p4_frozen_r1_buy_yp_entry_gap\endcsname{0.217154822}
\expandafter\def\csname p4num@p4_frozen_r1_buy_ym_entry_gap\endcsname{0.217154822}
\expandafter\def\csname p4num@p4_frozen_r1_buy_tp_payoff\endcsname{0.159294281}
\expandafter\def\csname p4num@p4_frozen_r1_buy_tm_payoff\endcsname{-0.507372386}
\expandafter\def\csname p4num@p4_frozen_r1_entry\endcsname{0.522757}
\expandafter\def\csname p4num@p4_frozen_r1_own\endcsname{0.324401}
\expandafter\def\csname p4num@p4_frozen_r1_proceeds\endcsname{0.872531}
\expandafter\def\csname p4num@p4_hidden_r0_buy_rate\endcsname{0.000000}
\expandafter\def\csname p4num@p4_hidden_r0_sell_rate\endcsname{0.000000}
\expandafter\def\csname p4num@p4_hidden_r0_sell_mu\endcsname{0.500000}
\expandafter\def\csname p4num@p4_hidden_r0_sell_nu\endcsname{0.500000}
\expandafter\def\csname p4num@p4_hidden_r0_sell_price\endcsname{0.392708}
\expandafter\def\csname p4num@p4_hidden_r0_sell_eh\endcsname{0.250000}
\expandafter\def\csname p4num@p4_hidden_r0_sell_el\endcsname{0.250000}
\expandafter\def\csname p4num@p4_hidden_r0_sell_yp_entry_gap\endcsname{-1.195833333}
\expandafter\def\csname p4num@p4_hidden_r0_sell_ym_entry_gap\endcsname{-1.195833333}
\expandafter\def\csname p4num@p4_hidden_r0_sell_tp_payoff\endcsname{-0.022083333}
\expandafter\def\csname p4num@p4_hidden_r0_sell_tm_payoff\endcsname{-0.017916667}
\expandafter\def\csname p4num@p4_hidden_r0_zero_mu\endcsname{0.500000}
\expandafter\def\csname p4num@p4_hidden_r0_zero_nu\endcsname{0.500000}
\expandafter\def\csname p4num@p4_hidden_r0_zero_price\endcsname{0.392708}
\expandafter\def\csname p4num@p4_hidden_r0_zero_eh\endcsname{0.250000}
\expandafter\def\csname p4num@p4_hidden_r0_zero_el\endcsname{0.250000}
\expandafter\def\csname p4num@p4_hidden_r0_zero_yp_entry_gap\endcsname{-1.195833333}
\expandafter\def\csname p4num@p4_hidden_r0_zero_ym_entry_gap\endcsname{-1.195833333}
\expandafter\def\csname p4num@p4_hidden_r0_zero_tp_payoff\endcsname{0.000000000}
\expandafter\def\csname p4num@p4_hidden_r0_zero_tm_payoff\endcsname{0.000000000}
\expandafter\def\csname p4num@p4_hidden_r0_buy_mu\endcsname{0.500000}
\expandafter\def\csname p4num@p4_hidden_r0_buy_nu\endcsname{0.500000}
\expandafter\def\csname p4num@p4_hidden_r0_buy_price\endcsname{0.392708}
\expandafter\def\csname p4num@p4_hidden_r0_buy_eh\endcsname{0.250000}
\expandafter\def\csname p4num@p4_hidden_r0_buy_el\endcsname{0.250000}
\expandafter\def\csname p4num@p4_hidden_r0_buy_yp_entry_gap\endcsname{-1.195833333}
\expandafter\def\csname p4num@p4_hidden_r0_buy_ym_entry_gap\endcsname{-1.195833333}
\expandafter\def\csname p4num@p4_hidden_r0_buy_tp_payoff\endcsname{-0.017916667}
\expandafter\def\csname p4num@p4_hidden_r0_buy_tm_payoff\endcsname{-0.022083333}
\expandafter\def\csname p4num@p4_hidden_r0_entry\endcsname{0.250000}
\expandafter\def\csname p4num@p4_hidden_r0_own\endcsname{0.125000}
\expandafter\def\csname p4num@p4_hidden_r0_proceeds\endcsname{0.392708}
\expandafter\def\csname p4num@p4_hidden_r1_buy_rate\endcsname{1.000000}
\expandafter\def\csname p4num@p4_hidden_r1_sell_rate\endcsname{1.000000}
\expandafter\def\csname p4num@p4_hidden_r1_sell_mu\endcsname{0.268941}
\expandafter\def\csname p4num@p4_hidden_r1_sell_nu\endcsname{0.268941}
\expandafter\def\csname p4num@p4_hidden_r1_sell_price\endcsname{0.576074}
\expandafter\def\csname p4num@p4_hidden_r1_sell_eh\endcsname{0.250000}
\expandafter\def\csname p4num@p4_hidden_r1_sell_el\endcsname{0.250000}
\expandafter\def\csname p4num@p4_hidden_r1_sell_yp_entry_gap\endcsname{-1.708333333}
\expandafter\def\csname p4num@p4_hidden_r1_sell_ym_entry_gap\endcsname{-1.708333333}
\expandafter\def\csname p4num@p4_hidden_r1_sell_tp_payoff\endcsname{-0.141843096}
\expandafter\def\csname p4num@p4_hidden_r1_sell_tm_payoff\endcsname{0.024823570}
\expandafter\def\csname p4num@p4_hidden_r1_zero_mu\endcsname{0.500000}
\expandafter\def\csname p4num@p4_hidden_r1_zero_nu\endcsname{0.500000}
\expandafter\def\csname p4num@p4_hidden_r1_zero_price\endcsname{0.614583}
\expandafter\def\csname p4num@p4_hidden_r1_zero_eh\endcsname{0.250000}
\expandafter\def\csname p4num@p4_hidden_r1_zero_el\endcsname{0.250000}
\expandafter\def\csname p4num@p4_hidden_r1_zero_yp_entry_gap\endcsname{-1.708333333}
\expandafter\def\csname p4num@p4_hidden_r1_zero_ym_entry_gap\endcsname{-1.708333333}
\expandafter\def\csname p4num@p4_hidden_r1_zero_tp_payoff\endcsname{0.000000000}
\expandafter\def\csname p4num@p4_hidden_r1_zero_tm_payoff\endcsname{0.000000000}
\expandafter\def\csname p4num@p4_hidden_r1_buy_mu\endcsname{0.731059}
\expandafter\def\csname p4num@p4_hidden_r1_buy_nu\endcsname{0.731059}
\expandafter\def\csname p4num@p4_hidden_r1_buy_price\endcsname{0.653093}
\expandafter\def\csname p4num@p4_hidden_r1_buy_eh\endcsname{0.250000}
\expandafter\def\csname p4num@p4_hidden_r1_buy_el\endcsname{0.250000}
\expandafter\def\csname p4num@p4_hidden_r1_buy_yp_entry_gap\endcsname{-1.708333333}
\expandafter\def\csname p4num@p4_hidden_r1_buy_ym_entry_gap\endcsname{-1.708333333}
\expandafter\def\csname p4num@p4_hidden_r1_buy_tp_payoff\endcsname{0.024823570}
\expandafter\def\csname p4num@p4_hidden_r1_buy_tm_payoff\endcsname{-0.141843096}
\expandafter\def\csname p4num@p4_hidden_r1_entry\endcsname{0.250000}
\expandafter\def\csname p4num@p4_hidden_r1_own\endcsname{0.125000}
\expandafter\def\csname p4num@p4_hidden_r1_proceeds\endcsname{0.614583}
\expandafter\def\csname p4num@p4_control_crossing_r0_r1\endcsname{0.000000}
\expandafter\def\csname p4num@p4_control_crossing_r1_r2\endcsname{0.272757}
\expandafter\def\csname p4num@rC_B\endcsname{3.592659}
\expandafter\def\csname p4num@p4_frozen_ceiling_lower_gap\endcsname{0.0000000000000000000000000000000000000000000000000000001}
\expandafter\def\csname p4num@p4_frozen_ceiling_upper_gap\endcsname{-0.0000000000000000000000000000000000000000000000000000000}
\expandafter\def\csname p4num@p4_frozen_entry_below_ceiling\endcsname{0.522757}
\expandafter\def\csname p4num@p4_frozen_entry_above_ceiling\endcsname{0.250000}
\expandafter\def\csname p4num@p4_atomless_r0_buy_rate\endcsname{0.000000}
\expandafter\def\csname p4num@p4_atomless_r0_sell_rate\endcsname{0.000000}
\expandafter\def\csname p4num@p4_atomless_r0_sell_mu\endcsname{0.500000}
\expandafter\def\csname p4num@p4_atomless_r0_sell_nu\endcsname{0.500000}
\expandafter\def\csname p4num@p4_atomless_r0_sell_price\endcsname{0.392708}
\expandafter\def\csname p4num@p4_atomless_r0_sell_eh\endcsname{0.250000}
\expandafter\def\csname p4num@p4_atomless_r0_sell_el\endcsname{0.250000}
\expandafter\def\csname p4num@p4_atomless_r0_sell_tp_payoff\endcsname{-0.022083333}
\expandafter\def\csname p4num@p4_atomless_r0_sell_tm_payoff\endcsname{-0.017916667}
\expandafter\def\csname p4num@p4_atomless_r0_zero_mu\endcsname{0.500000}
\expandafter\def\csname p4num@p4_atomless_r0_zero_nu\endcsname{0.500000}
\expandafter\def\csname p4num@p4_atomless_r0_zero_price\endcsname{0.392708}
\expandafter\def\csname p4num@p4_atomless_r0_zero_eh\endcsname{0.250000}
\expandafter\def\csname p4num@p4_atomless_r0_zero_el\endcsname{0.250000}
\expandafter\def\csname p4num@p4_atomless_r0_zero_tp_payoff\endcsname{0.000000000}
\expandafter\def\csname p4num@p4_atomless_r0_zero_tm_payoff\endcsname{0.000000000}
\expandafter\def\csname p4num@p4_atomless_r0_buy_mu\endcsname{0.500000}
\expandafter\def\csname p4num@p4_atomless_r0_buy_nu\endcsname{0.500000}
\expandafter\def\csname p4num@p4_atomless_r0_buy_price\endcsname{0.392708}
\expandafter\def\csname p4num@p4_atomless_r0_buy_eh\endcsname{0.250000}
\expandafter\def\csname p4num@p4_atomless_r0_buy_el\endcsname{0.250000}
\expandafter\def\csname p4num@p4_atomless_r0_buy_tp_payoff\endcsname{-0.017916667}
\expandafter\def\csname p4num@p4_atomless_r0_buy_tm_payoff\endcsname{-0.022083333}
\expandafter\def\csname p4num@p4_atomless_r0_entry\endcsname{0.250000}
\expandafter\def\csname p4num@p4_atomless_r0_own\endcsname{0.125000}
\expandafter\def\csname p4num@p4_atomless_r0_proceeds\endcsname{0.392708}
\expandafter\def\csname p4num@p4_atomless_r1_buy_rate\endcsname{1.000000}
\expandafter\def\csname p4num@p4_atomless_r1_sell_rate\endcsname{1.000000}
\expandafter\def\csname p4num@p4_atomless_r1_sell_mu\endcsname{0.268941}
\expandafter\def\csname p4num@p4_atomless_r1_sell_nu\endcsname{0.268941}
\expandafter\def\csname p4num@p4_atomless_r1_sell_price\endcsname{0.576074}
\expandafter\def\csname p4num@p4_atomless_r1_sell_eh\endcsname{0.250000}
\expandafter\def\csname p4num@p4_atomless_r1_sell_el\endcsname{0.250000}
\expandafter\def\csname p4num@p4_atomless_r1_sell_tp_payoff\endcsname{-0.141843096}
\expandafter\def\csname p4num@p4_atomless_r1_sell_tm_payoff\endcsname{0.024823570}
\expandafter\def\csname p4num@p4_atomless_r1_zero_mu\endcsname{0.500000}
\expandafter\def\csname p4num@p4_atomless_r1_zero_nu\endcsname{0.500000}
\expandafter\def\csname p4num@p4_atomless_r1_zero_price\endcsname{0.614583}
\expandafter\def\csname p4num@p4_atomless_r1_zero_eh\endcsname{0.250000}
\expandafter\def\csname p4num@p4_atomless_r1_zero_el\endcsname{0.250000}
\expandafter\def\csname p4num@p4_atomless_r1_zero_tp_payoff\endcsname{0.000000000}
\expandafter\def\csname p4num@p4_atomless_r1_zero_tm_payoff\endcsname{0.000000000}
\expandafter\def\csname p4num@p4_atomless_r1_buy_mu\endcsname{0.731059}
\expandafter\def\csname p4num@p4_atomless_r1_buy_nu\endcsname{0.731059}
\expandafter\def\csname p4num@p4_atomless_r1_buy_price\endcsname{1.362372}
\expandafter\def\csname p4num@p4_atomless_r1_buy_eh\endcsname{1.000000}
\expandafter\def\csname p4num@p4_atomless_r1_buy_el\endcsname{1.000000}
\expandafter\def\csname p4num@p4_atomless_r1_buy_tp_payoff\endcsname{0.159294281}
\expandafter\def\csname p4num@p4_atomless_r1_buy_tm_payoff\endcsname{-0.507372386}
\expandafter\def\csname p4num@p4_atomless_r1_entry\endcsname{0.522757}
\expandafter\def\csname p4num@p4_atomless_r1_own\endcsname{0.324401}
\expandafter\def\csname p4num@p4_atomless_r1_proceeds\endcsname{0.872531}
\expandafter\def\csname p4num@p4_atomless_r2_buy_rate\endcsname{1.000000}
\expandafter\def\csname p4num@p4_atomless_r2_sell_rate\endcsname{1.000000}
\expandafter\def\csname p4num@p4_atomless_r2_sell_mu\endcsname{0.268941}
\expandafter\def\csname p4num@p4_atomless_r2_sell_nu\endcsname{0.268941}
\expandafter\def\csname p4num@p4_atomless_r2_sell_price\endcsname{0.610002}
\expandafter\def\csname p4num@p4_atomless_r2_sell_eh\endcsname{0.250000}
\expandafter\def\csname p4num@p4_atomless_r2_sell_el\endcsname{0.250000}
\expandafter\def\csname p4num@p4_atomless_r2_sell_tp_payoff\endcsname{-0.191595694}
\expandafter\def\csname p4num@p4_atomless_r2_sell_tm_payoff\endcsname{0.043126528}
\expandafter\def\csname p4num@p4_atomless_r2_zero_mu\endcsname{0.500000}
\expandafter\def\csname p4num@p4_atomless_r2_zero_nu\endcsname{0.500000}
\expandafter\def\csname p4num@p4_atomless_r2_zero_price\endcsname{0.664236}
\expandafter\def\csname p4num@p4_atomless_r2_zero_eh\endcsname{0.250000}
\expandafter\def\csname p4num@p4_atomless_r2_zero_el\endcsname{0.250000}
\expandafter\def\csname p4num@p4_atomless_r2_zero_tp_payoff\endcsname{0.000000000}
\expandafter\def\csname p4num@p4_atomless_r2_zero_tm_payoff\endcsname{0.000000000}
\expandafter\def\csname p4num@p4_atomless_r2_buy_mu\endcsname{0.731059}
\expandafter\def\csname p4num@p4_atomless_r2_buy_nu\endcsname{0.731059}
\expandafter\def\csname p4num@p4_atomless_r2_buy_price\endcsname{1.138731}
\expandafter\def\csname p4num@p4_atomless_r2_buy_eh\endcsname{0.614917}
\expandafter\def\csname p4num@p4_atomless_r2_buy_el\endcsname{0.614917}
\expandafter\def\csname p4num@p4_atomless_r2_buy_tp_payoff\endcsname{0.135270204}
\expandafter\def\csname p4num@p4_atomless_r2_buy_tm_payoff\endcsname{-0.442068175}
\expandafter\def\csname p4num@p4_atomless_r2_entry\endcsname{0.382711}
\expandafter\def\csname p4num@p4_atomless_r2_own\endcsname{0.222020}
\expandafter\def\csname p4num@p4_atomless_r2_proceeds\endcsname{0.817075}
\expandafter\def\csname p4num@p4_atomless_max_width_floor_r2\endcsname{0.002688902}
\expandafter\def\csname p4num@p4_atomless_small_width\endcsname{0.001000}
\expandafter\def\csname p4num@margin_p4_atomless_small_width_return\endcsname{0.001688902}
\expandafter\def\csname p4num@margin_p4_signal_s_private\endcsname{0.052500000}
\expandafter\def\csname p4num@margin_p4_signal_s_joint\endcsname{0.045265153}
\expandafter\def\csname p4num@margin_p4_signal_s_weak\endcsname{0.001454545}
\expandafter\def\csname p4num@margin_p4_signal_s_strong\endcsname{0.018594291}
\expandafter\def\csname p4num@margin_p4_signal_s_other_pair\endcsname{4.159475417}
\expandafter\def\csname p4num@margin_p4_signal_b_private\endcsname{-1.050000000}
\expandafter\def\csname p4num@margin_p4_signal_b_joint\endcsname{0.903775226}
\expandafter\def\csname p4num@margin_p4_signal_b_weak\endcsname{-0.048541667}
\expandafter\def\csname p4num@margin_p4_signal_b_strong\endcsname{-0.002070572}
\expandafter\def\csname p4num@margin_p4_signal_b_other_pair\endcsname{3.155195358}
\expandafter\def\csname p4num@margin_p4_signal_valid_b_private\endcsname{1.150916667}
\expandafter\def\csname p4num@margin_p4_signal_valid_b_joint\endcsname{0.058922560}
\expandafter\def\csname p4num@margin_p4_signal_valid_b_weak\endcsname{0.002215625}
\expandafter\def\csname p4num@margin_p4_signal_valid_b_strong\endcsname{0.020341213}
\expandafter\def\csname p4num@margin_p4_signal_valid_b_other_pair\endcsname{0.009997188}
\expandafter\def\csname p4num@margin_p4_signal_comp_b_private\endcsname{1.150916667}
\expandafter\def\csname p4num@margin_p4_signal_comp_b_joint\endcsname{0.020731442}
\expandafter\def\csname p4num@margin_p4_signal_comp_b_weak\endcsname{0.002610833}
\expandafter\def\csname p4num@margin_p4_signal_comp_b_strong\endcsname{0.019444742}
\expandafter\def\csname p4num@margin_p4_signal_comp_b_other_pair\endcsname{0.048821802}
\expandafter\def\csname p4num@margin_p4_s_no_y_weak\endcsname{0.013181818}
\expandafter\def\csname p4num@margin_p4_s_no_y_strong\endcsname{0.018594291}
\expandafter\def\csname p4num@margin_p4_s_no_y_top_r0\endcsname{1.469972860}
\expandafter\def\csname p4num@margin_p4_s_no_y_top_r1\endcsname{1.862796856}
\expandafter\def\csname p4num@margin_p4_s_class_weak\endcsname{0.010454545}
\expandafter\def\csname p4num@margin_p4_s_class_strong\endcsname{0.068985729}
\expandafter\def\csname p4num@margin_p4_s_class_prior_r0\endcsname{2.301363636}
\expandafter\def\csname p4num@margin_p4_s_class_top_r1\endcsname{0.666013879}
\expandafter\def\csname p4num@margin_p4_b_no_y_weak\endcsname{0.019166667}
\expandafter\def\csname p4num@margin_p4_b_no_y_zero_gain\endcsname{0.013333333}
\expandafter\def\csname p4num@margin_p4_b_no_y_full_loss\endcsname{0.002070572}
\expandafter\def\csname p4num@p4_s_signal_r0_buy_rate\endcsname{0.000000}
\expandafter\def\csname p4num@p4_s_signal_r0_sell_rate\endcsname{0.000000}
\expandafter\def\csname p4num@p4_s_signal_r0_sell_mu\endcsname{0.500000}
\expandafter\def\csname p4num@p4_s_signal_r0_sell_nu\endcsname{0.500000}
\expandafter\def\csname p4num@p4_s_signal_r0_sell_price\endcsname{0.603068}
\expandafter\def\csname p4num@p4_s_signal_r0_sell_eh\endcsname{0.850000}
\expandafter\def\csname p4num@p4_s_signal_r0_sell_el\endcsname{0.850000}
\expandafter\def\csname p4num@p4_s_signal_r0_sell_yp_entry_gap\endcsname{-0.052500000}
\expandafter\def\csname p4num@p4_s_signal_r0_sell_ym_entry_gap\endcsname{-4.550227273}
\expandafter\def\csname p4num@p4_s_signal_r0_sell_tp_payoff\endcsname{-0.015772727}
\expandafter\def\csname p4num@p4_s_signal_r0_sell_tm_payoff\endcsname{-0.014227273}
\expandafter\def\csname p4num@p4_s_signal_r0_zero_mu\endcsname{0.500000}
\expandafter\def\csname p4num@p4_s_signal_r0_zero_nu\endcsname{0.500000}
\expandafter\def\csname p4num@p4_s_signal_r0_zero_price\endcsname{0.603068}
\expandafter\def\csname p4num@p4_s_signal_r0_zero_eh\endcsname{0.850000}
\expandafter\def\csname p4num@p4_s_signal_r0_zero_el\endcsname{0.850000}
\expandafter\def\csname p4num@p4_s_signal_r0_zero_yp_entry_gap\endcsname{-0.052500000}
\expandafter\def\csname p4num@p4_s_signal_r0_zero_ym_entry_gap\endcsname{-4.550227273}
\expandafter\def\csname p4num@p4_s_signal_r0_zero_tp_payoff\endcsname{0.000000000}
\expandafter\def\csname p4num@p4_s_signal_r0_zero_tm_payoff\endcsname{0.000000000}
\expandafter\def\csname p4num@p4_s_signal_r0_buy_mu\endcsname{0.500000}
\expandafter\def\csname p4num@p4_s_signal_r0_buy_nu\endcsname{0.500000}
\expandafter\def\csname p4num@p4_s_signal_r0_buy_price\endcsname{0.603068}
\expandafter\def\csname p4num@p4_s_signal_r0_buy_eh\endcsname{0.850000}
\expandafter\def\csname p4num@p4_s_signal_r0_buy_el\endcsname{0.850000}
\expandafter\def\csname p4num@p4_s_signal_r0_buy_yp_entry_gap\endcsname{-0.052500000}
\expandafter\def\csname p4num@p4_s_signal_r0_buy_ym_entry_gap\endcsname{-4.550227273}
\expandafter\def\csname p4num@p4_s_signal_r0_buy_tp_payoff\endcsname{-0.014227273}
\expandafter\def\csname p4num@p4_s_signal_r0_buy_tm_payoff\endcsname{-0.015772727}
\expandafter\def\csname p4num@p4_s_signal_r0_entry\endcsname{0.850000}
\expandafter\def\csname p4num@p4_s_signal_r0_own\endcsname{0.425000}
\expandafter\def\csname p4num@p4_s_signal_r0_proceeds\endcsname{0.603068}
\expandafter\def\csname p4num@p4_s_signal_r1_buy_rate\endcsname{1.000000}
\expandafter\def\csname p4num@p4_s_signal_r1_sell_rate\endcsname{1.000000}
\expandafter\def\csname p4num@p4_s_signal_r1_sell_mu\endcsname{0.407577}
\expandafter\def\csname p4num@p4_s_signal_r1_sell_nu\endcsname{0.268941}
\expandafter\def\csname p4num@p4_s_signal_r1_sell_price\endcsname{0.897388}
\expandafter\def\csname p4num@p4_s_signal_r1_sell_eh\endcsname{0.850000}
\expandafter\def\csname p4num@p4_s_signal_r1_sell_el\endcsname{0.850000}
\expandafter\def\csname p4num@p4_s_signal_r1_sell_yp_entry_gap\endcsname{-1.161828931}
\expandafter\def\csname p4num@p4_s_signal_r1_sell_ym_entry_gap\endcsname{-5.366569501}
\expandafter\def\csname p4num@p4_s_signal_r1_sell_tp_payoff\endcsname{-0.106318752}
\expandafter\def\csname p4num@p4_s_signal_r1_sell_tm_payoff\endcsname{0.018594291}
\expandafter\def\csname p4num@p4_s_signal_r1_zero_mu\endcsname{0.500000}
\expandafter\def\csname p4num@p4_s_signal_r1_zero_nu\endcsname{0.500000}
\expandafter\def\csname p4num@p4_s_signal_r1_zero_price\endcsname{0.926250}
\expandafter\def\csname p4num@p4_s_signal_r1_zero_eh\endcsname{0.850000}
\expandafter\def\csname p4num@p4_s_signal_r1_zero_el\endcsname{0.850000}
\expandafter\def\csname p4num@p4_s_signal_r1_zero_yp_entry_gap\endcsname{-0.502500000}
\expandafter\def\csname p4num@p4_s_signal_r1_zero_ym_entry_gap\endcsname{-4.818804348}
\expandafter\def\csname p4num@p4_s_signal_r1_zero_tp_payoff\endcsname{0.000000000}
\expandafter\def\csname p4num@p4_s_signal_r1_zero_tm_payoff\endcsname{0.000000000}
\expandafter\def\csname p4num@p4_s_signal_r1_buy_mu\endcsname{0.592423}
\expandafter\def\csname p4num@p4_s_signal_r1_buy_nu\endcsname{0.731059}
\expandafter\def\csname p4num@p4_s_signal_r1_buy_price\endcsname{1.016111}
\expandafter\def\csname p4num@p4_s_signal_r1_buy_eh\endcsname{0.962500}
\expandafter\def\csname p4num@p4_s_signal_r1_buy_el\endcsname{0.887500}
\expandafter\def\csname p4num@p4_s_signal_r1_buy_yp_entry_gap\endcsname{0.045265153}
\expandafter\def\csname p4num@p4_s_signal_r1_buy_ym_entry_gap\endcsname{-4.159475417}
\expandafter\def\csname p4num@p4_s_signal_r1_buy_tp_payoff\endcsname{0.026636225}
\expandafter\def\csname p4num@p4_s_signal_r1_buy_tm_payoff\endcsname{-0.128178993}
\expandafter\def\csname p4num@p4_s_signal_r1_entry\endcsname{0.879797}
\expandafter\def\csname p4num@p4_s_signal_r1_own\endcsname{0.449238}
\expandafter\def\csname p4num@p4_s_signal_r1_proceeds\endcsname{0.948434}
\expandafter\def\csname p4num@p4_s_no_y_r0_buy_rate\endcsname{0.000000}
\expandafter\def\csname p4num@p4_s_no_y_r0_sell_rate\endcsname{0.000000}
\expandafter\def\csname p4num@p4_s_no_y_r0_sell_mu\endcsname{0.500000}
\expandafter\def\csname p4num@p4_s_no_y_r0_sell_nu\endcsname{0.500000}
\expandafter\def\csname p4num@p4_s_no_y_r0_sell_price\endcsname{0.603068}
\expandafter\def\csname p4num@p4_s_no_y_r0_sell_eh\endcsname{0.850000}
\expandafter\def\csname p4num@p4_s_no_y_r0_sell_el\endcsname{0.850000}
\expandafter\def\csname p4num@p4_s_no_y_r0_sell_yp_entry_gap\endcsname{-2.301363636}
\expandafter\def\csname p4num@p4_s_no_y_r0_sell_ym_entry_gap\endcsname{-2.301363636}
\expandafter\def\csname p4num@p4_s_no_y_r0_sell_tp_payoff\endcsname{-0.015772727}
\expandafter\def\csname p4num@p4_s_no_y_r0_sell_tm_payoff\endcsname{-0.014227273}
\expandafter\def\csname p4num@p4_s_no_y_r0_zero_mu\endcsname{0.500000}
\expandafter\def\csname p4num@p4_s_no_y_r0_zero_nu\endcsname{0.500000}
\expandafter\def\csname p4num@p4_s_no_y_r0_zero_price\endcsname{0.603068}
\expandafter\def\csname p4num@p4_s_no_y_r0_zero_eh\endcsname{0.850000}
\expandafter\def\csname p4num@p4_s_no_y_r0_zero_el\endcsname{0.850000}
\expandafter\def\csname p4num@p4_s_no_y_r0_zero_yp_entry_gap\endcsname{-2.301363636}
\expandafter\def\csname p4num@p4_s_no_y_r0_zero_ym_entry_gap\endcsname{-2.301363636}
\expandafter\def\csname p4num@p4_s_no_y_r0_zero_tp_payoff\endcsname{0.000000000}
\expandafter\def\csname p4num@p4_s_no_y_r0_zero_tm_payoff\endcsname{0.000000000}
\expandafter\def\csname p4num@p4_s_no_y_r0_buy_mu\endcsname{0.500000}
\expandafter\def\csname p4num@p4_s_no_y_r0_buy_nu\endcsname{0.500000}
\expandafter\def\csname p4num@p4_s_no_y_r0_buy_price\endcsname{0.603068}
\expandafter\def\csname p4num@p4_s_no_y_r0_buy_eh\endcsname{0.850000}
\expandafter\def\csname p4num@p4_s_no_y_r0_buy_el\endcsname{0.850000}
\expandafter\def\csname p4num@p4_s_no_y_r0_buy_yp_entry_gap\endcsname{-2.301363636}
\expandafter\def\csname p4num@p4_s_no_y_r0_buy_ym_entry_gap\endcsname{-2.301363636}
\expandafter\def\csname p4num@p4_s_no_y_r0_buy_tp_payoff\endcsname{-0.014227273}
\expandafter\def\csname p4num@p4_s_no_y_r0_buy_tm_payoff\endcsname{-0.015772727}
\expandafter\def\csname p4num@p4_s_no_y_r0_entry\endcsname{0.850000}
\expandafter\def\csname p4num@p4_s_no_y_r0_own\endcsname{0.425000}
\expandafter\def\csname p4num@p4_s_no_y_r0_proceeds\endcsname{0.603068}
\expandafter\def\csname p4num@p4_s_no_y_r1_buy_rate\endcsname{1.000000}
\expandafter\def\csname p4num@p4_s_no_y_r1_sell_rate\endcsname{1.000000}
\expandafter\def\csname p4num@p4_s_no_y_r1_sell_mu\endcsname{0.407577}
\expandafter\def\csname p4num@p4_s_no_y_r1_sell_nu\endcsname{0.268941}
\expandafter\def\csname p4num@p4_s_no_y_r1_sell_price\endcsname{0.897388}
\expandafter\def\csname p4num@p4_s_no_y_r1_sell_eh\endcsname{0.850000}
\expandafter\def\csname p4num@p4_s_no_y_r1_sell_el\endcsname{0.850000}
\expandafter\def\csname p4num@p4_s_no_y_r1_sell_yp_entry_gap\endcsname{-3.458507492}
\expandafter\def\csname p4num@p4_s_no_y_r1_sell_ym_entry_gap\endcsname{-3.458507492}
\expandafter\def\csname p4num@p4_s_no_y_r1_sell_tp_payoff\endcsname{-0.106318752}
\expandafter\def\csname p4num@p4_s_no_y_r1_sell_tm_payoff\endcsname{0.018594291}
\expandafter\def\csname p4num@p4_s_no_y_r1_zero_mu\endcsname{0.500000}
\expandafter\def\csname p4num@p4_s_no_y_r1_zero_nu\endcsname{0.500000}
\expandafter\def\csname p4num@p4_s_no_y_r1_zero_price\endcsname{0.926250}
\expandafter\def\csname p4num@p4_s_no_y_r1_zero_eh\endcsname{0.850000}
\expandafter\def\csname p4num@p4_s_no_y_r1_zero_el\endcsname{0.850000}
\expandafter\def\csname p4num@p4_s_no_y_r1_zero_yp_entry_gap\endcsname{-2.660652174}
\expandafter\def\csname p4num@p4_s_no_y_r1_zero_ym_entry_gap\endcsname{-2.660652174}
\expandafter\def\csname p4num@p4_s_no_y_r1_zero_tp_payoff\endcsname{0.000000000}
\expandafter\def\csname p4num@p4_s_no_y_r1_zero_tm_payoff\endcsname{0.000000000}
\expandafter\def\csname p4num@p4_s_no_y_r1_buy_mu\endcsname{0.592423}
\expandafter\def\csname p4num@p4_s_no_y_r1_buy_nu\endcsname{0.731059}
\expandafter\def\csname p4num@p4_s_no_y_r1_buy_price\endcsname{0.955112}
\expandafter\def\csname p4num@p4_s_no_y_r1_buy_eh\endcsname{0.850000}
\expandafter\def\csname p4num@p4_s_no_y_r1_buy_el\endcsname{0.850000}
\expandafter\def\csname p4num@p4_s_no_y_r1_buy_yp_entry_gap\endcsname{-1.862796856}
\expandafter\def\csname p4num@p4_s_no_y_r1_buy_ym_entry_gap\endcsname{-1.862796856}
\expandafter\def\csname p4num@p4_s_no_y_r1_buy_tp_payoff\endcsname{0.018594291}
\expandafter\def\csname p4num@p4_s_no_y_r1_buy_tm_payoff\endcsname{-0.106318752}
\expandafter\def\csname p4num@p4_s_no_y_r1_entry\endcsname{0.850000}
\expandafter\def\csname p4num@p4_s_no_y_r1_own\endcsname{0.425000}
\expandafter\def\csname p4num@p4_s_no_y_r1_proceeds\endcsname{0.926250}
\expandafter\def\csname p4num@p4_s_class_r0_buy_rate\endcsname{0.000000}
\expandafter\def\csname p4num@p4_s_class_r0_sell_rate\endcsname{0.000000}
\expandafter\def\csname p4num@p4_s_class_r0_sell_mu\endcsname{0.500000}
\expandafter\def\csname p4num@p4_s_class_r0_sell_nu\endcsname{0.500000}
\expandafter\def\csname p4num@p4_s_class_r0_sell_price\endcsname{0.603068}
\expandafter\def\csname p4num@p4_s_class_r0_sell_eh\endcsname{0.850000}
\expandafter\def\csname p4num@p4_s_class_r0_sell_el\endcsname{0.850000}
\expandafter\def\csname p4num@p4_s_class_r0_sell_yp_entry_gap\endcsname{-2.301363636}
\expandafter\def\csname p4num@p4_s_class_r0_sell_ym_entry_gap\endcsname{-2.301363636}
\expandafter\def\csname p4num@p4_s_class_r0_sell_tp_payoff\endcsname{-0.016931818}
\expandafter\def\csname p4num@p4_s_class_r0_sell_tm_payoff\endcsname{-0.013068182}
\expandafter\def\csname p4num@p4_s_class_r0_zero_mu\endcsname{0.500000}
\expandafter\def\csname p4num@p4_s_class_r0_zero_nu\endcsname{0.500000}
\expandafter\def\csname p4num@p4_s_class_r0_zero_price\endcsname{0.603068}
\expandafter\def\csname p4num@p4_s_class_r0_zero_eh\endcsname{0.850000}
\expandafter\def\csname p4num@p4_s_class_r0_zero_el\endcsname{0.850000}
\expandafter\def\csname p4num@p4_s_class_r0_zero_yp_entry_gap\endcsname{-2.301363636}
\expandafter\def\csname p4num@p4_s_class_r0_zero_ym_entry_gap\endcsname{-2.301363636}
\expandafter\def\csname p4num@p4_s_class_r0_zero_tp_payoff\endcsname{0.000000000}
\expandafter\def\csname p4num@p4_s_class_r0_zero_tm_payoff\endcsname{0.000000000}
\expandafter\def\csname p4num@p4_s_class_r0_buy_mu\endcsname{0.500000}
\expandafter\def\csname p4num@p4_s_class_r0_buy_nu\endcsname{0.500000}
\expandafter\def\csname p4num@p4_s_class_r0_buy_price\endcsname{0.603068}
\expandafter\def\csname p4num@p4_s_class_r0_buy_eh\endcsname{0.850000}
\expandafter\def\csname p4num@p4_s_class_r0_buy_el\endcsname{0.850000}
\expandafter\def\csname p4num@p4_s_class_r0_buy_yp_entry_gap\endcsname{-2.301363636}
\expandafter\def\csname p4num@p4_s_class_r0_buy_ym_entry_gap\endcsname{-2.301363636}
\expandafter\def\csname p4num@p4_s_class_r0_buy_tp_payoff\endcsname{-0.013068182}
\expandafter\def\csname p4num@p4_s_class_r0_buy_tm_payoff\endcsname{-0.016931818}
\expandafter\def\csname p4num@p4_s_class_r0_entry\endcsname{0.850000}
\expandafter\def\csname p4num@p4_s_class_r0_own\endcsname{0.425000}
\expandafter\def\csname p4num@p4_s_class_r0_proceeds\endcsname{0.603068}
\expandafter\def\csname p4num@p4_s_class_r1_buy_rate\endcsname{1.000000}
\expandafter\def\csname p4num@p4_s_class_r1_sell_rate\endcsname{1.000000}
\expandafter\def\csname p4num@p4_s_class_r1_sell_mu\endcsname{0.268941}
\expandafter\def\csname p4num@p4_s_class_r1_sell_nu\endcsname{0.268941}
\expandafter\def\csname p4num@p4_s_class_r1_sell_price\endcsname{0.854094}
\expandafter\def\csname p4num@p4_s_class_r1_sell_eh\endcsname{0.850000}
\expandafter\def\csname p4num@p4_s_class_r1_sell_el\endcsname{0.850000}
\expandafter\def\csname p4num@p4_s_class_r1_sell_yp_entry_gap\endcsname{-4.655290469}
\expandafter\def\csname p4num@p4_s_class_r1_sell_ym_entry_gap\endcsname{-4.655290469}
\expandafter\def\csname p4num@p4_s_class_r1_sell_tp_payoff\endcsname{-0.243296880}
\expandafter\def\csname p4num@p4_s_class_r1_sell_tm_payoff\endcsname{0.068985729}
\expandafter\def\csname p4num@p4_s_class_r1_zero_mu\endcsname{0.500000}
\expandafter\def\csname p4num@p4_s_class_r1_zero_nu\endcsname{0.500000}
\expandafter\def\csname p4num@p4_s_class_r1_zero_price\endcsname{0.926250}
\expandafter\def\csname p4num@p4_s_class_r1_zero_eh\endcsname{0.850000}
\expandafter\def\csname p4num@p4_s_class_r1_zero_el\endcsname{0.850000}
\expandafter\def\csname p4num@p4_s_class_r1_zero_yp_entry_gap\endcsname{-2.660652174}
\expandafter\def\csname p4num@p4_s_class_r1_zero_ym_entry_gap\endcsname{-2.660652174}
\expandafter\def\csname p4num@p4_s_class_r1_zero_tp_payoff\endcsname{0.000000000}
\expandafter\def\csname p4num@p4_s_class_r1_zero_tm_payoff\endcsname{0.000000000}
\expandafter\def\csname p4num@p4_s_class_r1_buy_mu\endcsname{0.731059}
\expandafter\def\csname p4num@p4_s_class_r1_buy_nu\endcsname{0.731059}
\expandafter\def\csname p4num@p4_s_class_r1_buy_price\endcsname{0.998406}
\expandafter\def\csname p4num@p4_s_class_r1_buy_eh\endcsname{0.850000}
\expandafter\def\csname p4num@p4_s_class_r1_buy_el\endcsname{0.850000}
\expandafter\def\csname p4num@p4_s_class_r1_buy_yp_entry_gap\endcsname{-0.666013879}
\expandafter\def\csname p4num@p4_s_class_r1_buy_ym_entry_gap\endcsname{-0.666013879}
\expandafter\def\csname p4num@p4_s_class_r1_buy_tp_payoff\endcsname{0.068985729}
\expandafter\def\csname p4num@p4_s_class_r1_buy_tm_payoff\endcsname{-0.243296880}
\expandafter\def\csname p4num@p4_s_class_r1_entry\endcsname{0.850000}
\expandafter\def\csname p4num@p4_s_class_r1_own\endcsname{0.425000}
\expandafter\def\csname p4num@p4_s_class_r1_proceeds\endcsname{0.926250}
\expandafter\def\csname p4num@p4_valid_b_r0_buy_rate\endcsname{0.000000}
\expandafter\def\csname p4num@p4_valid_b_r0_sell_rate\endcsname{0.000000}
\expandafter\def\csname p4num@p4_valid_b_r0_sell_mu\endcsname{0.500000}
\expandafter\def\csname p4num@p4_valid_b_r0_sell_nu\endcsname{0.500000}
\expandafter\def\csname p4num@p4_valid_b_r0_sell_price\endcsname{0.392708}
\expandafter\def\csname p4num@p4_valid_b_r0_sell_eh\endcsname{0.250000}
\expandafter\def\csname p4num@p4_valid_b_r0_sell_el\endcsname{0.250000}
\expandafter\def\csname p4num@p4_valid_b_r0_sell_yp_entry_gap\endcsname{-1.150916667}
\expandafter\def\csname p4num@p4_valid_b_r0_sell_ym_entry_gap\endcsname{-1.240750000}
\expandafter\def\csname p4num@p4_valid_b_r0_sell_tp_payoff\endcsname{-0.021875000}
\expandafter\def\csname p4num@p4_valid_b_r0_sell_tm_payoff\endcsname{-0.018125000}
\expandafter\def\csname p4num@p4_valid_b_r0_zero_mu\endcsname{0.500000}
\expandafter\def\csname p4num@p4_valid_b_r0_zero_nu\endcsname{0.500000}
\expandafter\def\csname p4num@p4_valid_b_r0_zero_price\endcsname{0.392708}
\expandafter\def\csname p4num@p4_valid_b_r0_zero_eh\endcsname{0.250000}
\expandafter\def\csname p4num@p4_valid_b_r0_zero_el\endcsname{0.250000}
\expandafter\def\csname p4num@p4_valid_b_r0_zero_yp_entry_gap\endcsname{-1.150916667}
\expandafter\def\csname p4num@p4_valid_b_r0_zero_ym_entry_gap\endcsname{-1.240750000}
\expandafter\def\csname p4num@p4_valid_b_r0_zero_tp_payoff\endcsname{0.000000000}
\expandafter\def\csname p4num@p4_valid_b_r0_zero_tm_payoff\endcsname{0.000000000}
\expandafter\def\csname p4num@p4_valid_b_r0_buy_mu\endcsname{0.500000}
\expandafter\def\csname p4num@p4_valid_b_r0_buy_nu\endcsname{0.500000}
\expandafter\def\csname p4num@p4_valid_b_r0_buy_price\endcsname{0.392708}
\expandafter\def\csname p4num@p4_valid_b_r0_buy_eh\endcsname{0.250000}
\expandafter\def\csname p4num@p4_valid_b_r0_buy_el\endcsname{0.250000}
\expandafter\def\csname p4num@p4_valid_b_r0_buy_yp_entry_gap\endcsname{-1.150916667}
\expandafter\def\csname p4num@p4_valid_b_r0_buy_ym_entry_gap\endcsname{-1.240750000}
\expandafter\def\csname p4num@p4_valid_b_r0_buy_tp_payoff\endcsname{-0.018125000}
\expandafter\def\csname p4num@p4_valid_b_r0_buy_tm_payoff\endcsname{-0.021875000}
\expandafter\def\csname p4num@p4_valid_b_r0_entry\endcsname{0.250000}
\expandafter\def\csname p4num@p4_valid_b_r0_own\endcsname{0.125000}
\expandafter\def\csname p4num@p4_valid_b_r0_proceeds\endcsname{0.392708}
\expandafter\def\csname p4num@p4_valid_b_r1_buy_rate\endcsname{1.000000}
\expandafter\def\csname p4num@p4_valid_b_r1_sell_rate\endcsname{1.000000}
\expandafter\def\csname p4num@p4_valid_b_r1_sell_mu\endcsname{0.292047}
\expandafter\def\csname p4num@p4_valid_b_r1_sell_nu\endcsname{0.268941}
\expandafter\def\csname p4num@p4_valid_b_r1_sell_price\endcsname{0.579925}
\expandafter\def\csname p4num@p4_valid_b_r1_sell_eh\endcsname{0.250000}
\expandafter\def\csname p4num@p4_valid_b_r1_sell_el\endcsname{0.250000}
\expandafter\def\csname p4num@p4_valid_b_r1_sell_yp_entry_gap\endcsname{-3.406669479}
\expandafter\def\csname p4num@p4_valid_b_r1_sell_ym_entry_gap\endcsname{-3.475589226}
\expandafter\def\csname p4num@p4_valid_b_r1_sell_tp_payoff\endcsname{-0.129658787}
\expandafter\def\csname p4num@p4_valid_b_r1_sell_tm_payoff\endcsname{0.020341213}
\expandafter\def\csname p4num@p4_valid_b_r1_zero_mu\endcsname{0.500000}
\expandafter\def\csname p4num@p4_valid_b_r1_zero_nu\endcsname{0.500000}
\expandafter\def\csname p4num@p4_valid_b_r1_zero_price\endcsname{0.614583}
\expandafter\def\csname p4num@p4_valid_b_r1_zero_eh\endcsname{0.250000}
\expandafter\def\csname p4num@p4_valid_b_r1_zero_el\endcsname{0.250000}
\expandafter\def\csname p4num@p4_valid_b_r1_zero_yp_entry_gap\endcsname{-1.666666667}
\expandafter\def\csname p4num@p4_valid_b_r1_zero_ym_entry_gap\endcsname{-1.750000000}
\expandafter\def\csname p4num@p4_valid_b_r1_zero_tp_payoff\endcsname{0.000000000}
\expandafter\def\csname p4num@p4_valid_b_r1_zero_tm_payoff\endcsname{0.000000000}
\expandafter\def\csname p4num@p4_valid_b_r1_buy_mu\endcsname{0.707953}
\expandafter\def\csname p4num@p4_valid_b_r1_buy_nu\endcsname{0.731059}
\expandafter\def\csname p4num@p4_valid_b_r1_buy_price\endcsname{1.000590}
\expandafter\def\csname p4num@p4_valid_b_r1_buy_eh\endcsname{0.628750}
\expandafter\def\csname p4num@p4_valid_b_r1_buy_el\endcsname{0.621250}
\expandafter\def\csname p4num@p4_valid_b_r1_buy_yp_entry_gap\endcsname{0.058922560}
\expandafter\def\csname p4num@p4_valid_b_r1_buy_ym_entry_gap\endcsname{-0.009997188}
\expandafter\def\csname p4num@p4_valid_b_r1_buy_tp_payoff\endcsname{0.082290189}
\expandafter\def\csname p4num@p4_valid_b_r1_buy_tm_payoff\endcsname{-0.298053561}
\expandafter\def\csname p4num@p4_valid_b_r1_entry\endcsname{0.386946}
\expandafter\def\csname p4num@p4_valid_b_r1_own\endcsname{0.222515}
\expandafter\def\csname p4num@p4_valid_b_r1_proceeds\endcsname{0.742360}
\expandafter\def\csname p4num@p4_comp_b_r0_buy_rate\endcsname{0.000000}
\expandafter\def\csname p4num@p4_comp_b_r0_sell_rate\endcsname{0.000000}
\expandafter\def\csname p4num@p4_comp_b_r0_sell_mu\endcsname{0.500000}
\expandafter\def\csname p4num@p4_comp_b_r0_sell_nu\endcsname{0.500000}
\expandafter\def\csname p4num@p4_comp_b_r0_sell_price\endcsname{0.392708}
\expandafter\def\csname p4num@p4_comp_b_r0_sell_eh\endcsname{0.250000}
\expandafter\def\csname p4num@p4_comp_b_r0_sell_el\endcsname{0.250000}
\expandafter\def\csname p4num@p4_comp_b_r0_sell_yp_entry_gap\endcsname{-1.150916667}
\expandafter\def\csname p4num@p4_comp_b_r0_sell_ym_entry_gap\endcsname{-1.240750000}
\expandafter\def\csname p4num@p4_comp_b_r0_sell_tp_payoff\endcsname{-0.021833333}
\expandafter\def\csname p4num@p4_comp_b_r0_sell_tm_payoff\endcsname{-0.018166667}
\expandafter\def\csname p4num@p4_comp_b_r0_zero_mu\endcsname{0.500000}
\expandafter\def\csname p4num@p4_comp_b_r0_zero_nu\endcsname{0.500000}
\expandafter\def\csname p4num@p4_comp_b_r0_zero_price\endcsname{0.392708}
\expandafter\def\csname p4num@p4_comp_b_r0_zero_eh\endcsname{0.250000}
\expandafter\def\csname p4num@p4_comp_b_r0_zero_el\endcsname{0.250000}
\expandafter\def\csname p4num@p4_comp_b_r0_zero_yp_entry_gap\endcsname{-1.150916667}
\expandafter\def\csname p4num@p4_comp_b_r0_zero_ym_entry_gap\endcsname{-1.240750000}
\expandafter\def\csname p4num@p4_comp_b_r0_zero_tp_payoff\endcsname{0.000000000}
\expandafter\def\csname p4num@p4_comp_b_r0_zero_tm_payoff\endcsname{0.000000000}
\expandafter\def\csname p4num@p4_comp_b_r0_buy_mu\endcsname{0.500000}
\expandafter\def\csname p4num@p4_comp_b_r0_buy_nu\endcsname{0.500000}
\expandafter\def\csname p4num@p4_comp_b_r0_buy_price\endcsname{0.392708}
\expandafter\def\csname p4num@p4_comp_b_r0_buy_eh\endcsname{0.250000}
\expandafter\def\csname p4num@p4_comp_b_r0_buy_el\endcsname{0.250000}
\expandafter\def\csname p4num@p4_comp_b_r0_buy_yp_entry_gap\endcsname{-1.150916667}
\expandafter\def\csname p4num@p4_comp_b_r0_buy_ym_entry_gap\endcsname{-1.240750000}
\expandafter\def\csname p4num@p4_comp_b_r0_buy_tp_payoff\endcsname{-0.018166667}
\expandafter\def\csname p4num@p4_comp_b_r0_buy_tm_payoff\endcsname{-0.021833333}
\expandafter\def\csname p4num@p4_comp_b_r0_entry\endcsname{0.250000}
\expandafter\def\csname p4num@p4_comp_b_r0_own\endcsname{0.125000}
\expandafter\def\csname p4num@p4_comp_b_r0_proceeds\endcsname{0.392708}
\expandafter\def\csname p4num@p4_comp_b_r1_buy_rate\endcsname{1.000000}
\expandafter\def\csname p4num@p4_comp_b_r1_sell_rate\endcsname{1.000000}
\expandafter\def\csname p4num@p4_comp_b_r1_sell_mu\endcsname{0.296668}
\expandafter\def\csname p4num@p4_comp_b_r1_sell_nu\endcsname{0.268941}
\expandafter\def\csname p4num@p4_comp_b_r1_sell_price\endcsname{0.580695}
\expandafter\def\csname p4num@p4_comp_b_r1_sell_eh\endcsname{0.250000}
\expandafter\def\csname p4num@p4_comp_b_r1_sell_el\endcsname{0.250000}
\expandafter\def\csname p4num@p4_comp_b_r1_sell_yp_entry_gap\endcsname{-3.367844864}
\expandafter\def\csname p4num@p4_comp_b_r1_sell_ym_entry_gap\endcsname{-3.437398108}
\expandafter\def\csname p4num@p4_comp_b_r1_sell_tp_payoff\endcsname{-0.127221925}
\expandafter\def\csname p4num@p4_comp_b_r1_sell_tm_payoff\endcsname{0.019444742}
\expandafter\def\csname p4num@p4_comp_b_r1_zero_mu\endcsname{0.500000}
\expandafter\def\csname p4num@p4_comp_b_r1_zero_nu\endcsname{0.500000}
\expandafter\def\csname p4num@p4_comp_b_r1_zero_price\endcsname{0.614583}
\expandafter\def\csname p4num@p4_comp_b_r1_zero_eh\endcsname{0.250000}
\expandafter\def\csname p4num@p4_comp_b_r1_zero_el\endcsname{0.250000}
\expandafter\def\csname p4num@p4_comp_b_r1_zero_yp_entry_gap\endcsname{-1.666666667}
\expandafter\def\csname p4num@p4_comp_b_r1_zero_ym_entry_gap\endcsname{-1.750000000}
\expandafter\def\csname p4num@p4_comp_b_r1_zero_tp_payoff\endcsname{0.000000000}
\expandafter\def\csname p4num@p4_comp_b_r1_zero_tm_payoff\endcsname{0.000000000}
\expandafter\def\csname p4num@p4_comp_b_r1_buy_mu\endcsname{0.703332}
\expandafter\def\csname p4num@p4_comp_b_r1_buy_nu\endcsname{0.731059}
\expandafter\def\csname p4num@p4_comp_b_r1_buy_price\endcsname{0.998637}
\expandafter\def\csname p4num@p4_comp_b_r1_buy_eh\endcsname{0.628750}
\expandafter\def\csname p4num@p4_comp_b_r1_buy_el\endcsname{0.621250}
\expandafter\def\csname p4num@p4_comp_b_r1_buy_yp_entry_gap\endcsname{0.020731442}
\expandafter\def\csname p4num@p4_comp_b_r1_buy_ym_entry_gap\endcsname{-0.048821802}
\expandafter\def\csname p4num@p4_comp_b_r1_buy_tp_payoff\endcsname{0.080017073}
\expandafter\def\csname p4num@p4_comp_b_r1_buy_tm_payoff\endcsname{-0.291874593}
\expandafter\def\csname p4num@p4_comp_b_r1_entry\endcsname{0.386933}
\expandafter\def\csname p4num@p4_comp_b_r1_own\endcsname{0.221878}
\expandafter\def\csname p4num@p4_comp_b_r1_proceeds\endcsname{0.741930}
\expandafter\def\csname p4num@p4_comp_b_no_y_r0_buy_rate\endcsname{0.000000}
\expandafter\def\csname p4num@p4_comp_b_no_y_r0_sell_rate\endcsname{0.000000}
\expandafter\def\csname p4num@p4_comp_b_no_y_r0_sell_mu\endcsname{0.500000}
\expandafter\def\csname p4num@p4_comp_b_no_y_r0_sell_nu\endcsname{0.500000}
\expandafter\def\csname p4num@p4_comp_b_no_y_r0_sell_price\endcsname{0.392708}
\expandafter\def\csname p4num@p4_comp_b_no_y_r0_sell_eh\endcsname{0.250000}
\expandafter\def\csname p4num@p4_comp_b_no_y_r0_sell_el\endcsname{0.250000}
\expandafter\def\csname p4num@p4_comp_b_no_y_r0_sell_yp_entry_gap\endcsname{-1.195833333}
\expandafter\def\csname p4num@p4_comp_b_no_y_r0_sell_ym_entry_gap\endcsname{-1.195833333}
\expandafter\def\csname p4num@p4_comp_b_no_y_r0_sell_tp_payoff\endcsname{-0.021833333}
\expandafter\def\csname p4num@p4_comp_b_no_y_r0_sell_tm_payoff\endcsname{-0.018166667}
\expandafter\def\csname p4num@p4_comp_b_no_y_r0_zero_mu\endcsname{0.500000}
\expandafter\def\csname p4num@p4_comp_b_no_y_r0_zero_nu\endcsname{0.500000}
\expandafter\def\csname p4num@p4_comp_b_no_y_r0_zero_price\endcsname{0.392708}
\expandafter\def\csname p4num@p4_comp_b_no_y_r0_zero_eh\endcsname{0.250000}
\expandafter\def\csname p4num@p4_comp_b_no_y_r0_zero_el\endcsname{0.250000}
\expandafter\def\csname p4num@p4_comp_b_no_y_r0_zero_yp_entry_gap\endcsname{-1.195833333}
\expandafter\def\csname p4num@p4_comp_b_no_y_r0_zero_ym_entry_gap\endcsname{-1.195833333}
\expandafter\def\csname p4num@p4_comp_b_no_y_r0_zero_tp_payoff\endcsname{0.000000000}
\expandafter\def\csname p4num@p4_comp_b_no_y_r0_zero_tm_payoff\endcsname{0.000000000}
\expandafter\def\csname p4num@p4_comp_b_no_y_r0_buy_mu\endcsname{0.500000}
\expandafter\def\csname p4num@p4_comp_b_no_y_r0_buy_nu\endcsname{0.500000}
\expandafter\def\csname p4num@p4_comp_b_no_y_r0_buy_price\endcsname{0.392708}
\expandafter\def\csname p4num@p4_comp_b_no_y_r0_buy_eh\endcsname{0.250000}
\expandafter\def\csname p4num@p4_comp_b_no_y_r0_buy_el\endcsname{0.250000}
\expandafter\def\csname p4num@p4_comp_b_no_y_r0_buy_yp_entry_gap\endcsname{-1.195833333}
\expandafter\def\csname p4num@p4_comp_b_no_y_r0_buy_ym_entry_gap\endcsname{-1.195833333}
\expandafter\def\csname p4num@p4_comp_b_no_y_r0_buy_tp_payoff\endcsname{-0.018166667}
\expandafter\def\csname p4num@p4_comp_b_no_y_r0_buy_tm_payoff\endcsname{-0.021833333}
\expandafter\def\csname p4num@p4_comp_b_no_y_r0_entry\endcsname{0.250000}
\expandafter\def\csname p4num@p4_comp_b_no_y_r0_own\endcsname{0.125000}
\expandafter\def\csname p4num@p4_comp_b_no_y_r0_proceeds\endcsname{0.392708}
\expandafter\def\csname p4num@p4_comp_b_no_y_r1_buy_rate\endcsname{1.000000}
\expandafter\def\csname p4num@p4_comp_b_no_y_r1_sell_rate\endcsname{1.000000}
\expandafter\def\csname p4num@p4_comp_b_no_y_r1_sell_mu\endcsname{0.296668}
\expandafter\def\csname p4num@p4_comp_b_no_y_r1_sell_nu\endcsname{0.268941}
\expandafter\def\csname p4num@p4_comp_b_no_y_r1_sell_price\endcsname{0.580695}
\expandafter\def\csname p4num@p4_comp_b_no_y_r1_sell_eh\endcsname{0.250000}
\expandafter\def\csname p4num@p4_comp_b_no_y_r1_sell_el\endcsname{0.250000}
\expandafter\def\csname p4num@p4_comp_b_no_y_r1_sell_yp_entry_gap\endcsname{-3.402762910}
\expandafter\def\csname p4num@p4_comp_b_no_y_r1_sell_ym_entry_gap\endcsname{-3.402762910}
\expandafter\def\csname p4num@p4_comp_b_no_y_r1_sell_tp_payoff\endcsname{-0.127221925}
\expandafter\def\csname p4num@p4_comp_b_no_y_r1_sell_tm_payoff\endcsname{0.019444742}
\expandafter\def\csname p4num@p4_comp_b_no_y_r1_zero_mu\endcsname{0.500000}
\expandafter\def\csname p4num@p4_comp_b_no_y_r1_zero_nu\endcsname{0.500000}
\expandafter\def\csname p4num@p4_comp_b_no_y_r1_zero_price\endcsname{0.614583}
\expandafter\def\csname p4num@p4_comp_b_no_y_r1_zero_eh\endcsname{0.250000}
\expandafter\def\csname p4num@p4_comp_b_no_y_r1_zero_el\endcsname{0.250000}
\expandafter\def\csname p4num@p4_comp_b_no_y_r1_zero_yp_entry_gap\endcsname{-1.708333333}
\expandafter\def\csname p4num@p4_comp_b_no_y_r1_zero_ym_entry_gap\endcsname{-1.708333333}
\expandafter\def\csname p4num@p4_comp_b_no_y_r1_zero_tp_payoff\endcsname{0.000000000}
\expandafter\def\csname p4num@p4_comp_b_no_y_r1_zero_tm_payoff\endcsname{0.000000000}
\expandafter\def\csname p4num@p4_comp_b_no_y_r1_buy_mu\endcsname{0.703332}
\expandafter\def\csname p4num@p4_comp_b_no_y_r1_buy_nu\endcsname{0.731059}
\expandafter\def\csname p4num@p4_comp_b_no_y_r1_buy_price\endcsname{0.648472}
\expandafter\def\csname p4num@p4_comp_b_no_y_r1_buy_eh\endcsname{0.250000}
\expandafter\def\csname p4num@p4_comp_b_no_y_r1_buy_el\endcsname{0.250000}
\expandafter\def\csname p4num@p4_comp_b_no_y_r1_buy_yp_entry_gap\endcsname{-0.013903757}
\expandafter\def\csname p4num@p4_comp_b_no_y_r1_buy_ym_entry_gap\endcsname{-0.013903757}
\expandafter\def\csname p4num@p4_comp_b_no_y_r1_buy_tp_payoff\endcsname{0.019444742}
\expandafter\def\csname p4num@p4_comp_b_no_y_r1_buy_tm_payoff\endcsname{-0.127221925}
\expandafter\def\csname p4num@p4_comp_b_no_y_r1_entry\endcsname{0.250000}
\expandafter\def\csname p4num@p4_comp_b_no_y_r1_own\endcsname{0.125000}
\expandafter\def\csname p4num@p4_comp_b_no_y_r1_proceeds\endcsname{0.614583}
\expandafter\def\csname p4num@p4_comp_b_class_r0_buy_rate\endcsname{0.000000}
\expandafter\def\csname p4num@p4_comp_b_class_r0_sell_rate\endcsname{0.000000}
\expandafter\def\csname p4num@p4_comp_b_class_r0_sell_mu\endcsname{0.500000}
\expandafter\def\csname p4num@p4_comp_b_class_r0_sell_nu\endcsname{0.500000}
\expandafter\def\csname p4num@p4_comp_b_class_r0_sell_price\endcsname{0.392708}
\expandafter\def\csname p4num@p4_comp_b_class_r0_sell_eh\endcsname{0.250000}
\expandafter\def\csname p4num@p4_comp_b_class_r0_sell_el\endcsname{0.250000}
\expandafter\def\csname p4num@p4_comp_b_class_r0_sell_yp_entry_gap\endcsname{-1.195833333}
\expandafter\def\csname p4num@p4_comp_b_class_r0_sell_ym_entry_gap\endcsname{-1.195833333}
\expandafter\def\csname p4num@p4_comp_b_class_r0_sell_tp_payoff\endcsname{-0.022083333}
\expandafter\def\csname p4num@p4_comp_b_class_r0_sell_tm_payoff\endcsname{-0.017916667}
\expandafter\def\csname p4num@p4_comp_b_class_r0_zero_mu\endcsname{0.500000}
\expandafter\def\csname p4num@p4_comp_b_class_r0_zero_nu\endcsname{0.500000}
\expandafter\def\csname p4num@p4_comp_b_class_r0_zero_price\endcsname{0.392708}
\expandafter\def\csname p4num@p4_comp_b_class_r0_zero_eh\endcsname{0.250000}
\expandafter\def\csname p4num@p4_comp_b_class_r0_zero_el\endcsname{0.250000}
\expandafter\def\csname p4num@p4_comp_b_class_r0_zero_yp_entry_gap\endcsname{-1.195833333}
\expandafter\def\csname p4num@p4_comp_b_class_r0_zero_ym_entry_gap\endcsname{-1.195833333}
\expandafter\def\csname p4num@p4_comp_b_class_r0_zero_tp_payoff\endcsname{0.000000000}
\expandafter\def\csname p4num@p4_comp_b_class_r0_zero_tm_payoff\endcsname{0.000000000}
\expandafter\def\csname p4num@p4_comp_b_class_r0_buy_mu\endcsname{0.500000}
\expandafter\def\csname p4num@p4_comp_b_class_r0_buy_nu\endcsname{0.500000}
\expandafter\def\csname p4num@p4_comp_b_class_r0_buy_price\endcsname{0.392708}
\expandafter\def\csname p4num@p4_comp_b_class_r0_buy_eh\endcsname{0.250000}
\expandafter\def\csname p4num@p4_comp_b_class_r0_buy_el\endcsname{0.250000}
\expandafter\def\csname p4num@p4_comp_b_class_r0_buy_yp_entry_gap\endcsname{-1.195833333}
\expandafter\def\csname p4num@p4_comp_b_class_r0_buy_ym_entry_gap\endcsname{-1.195833333}
\expandafter\def\csname p4num@p4_comp_b_class_r0_buy_tp_payoff\endcsname{-0.017916667}
\expandafter\def\csname p4num@p4_comp_b_class_r0_buy_tm_payoff\endcsname{-0.022083333}
\expandafter\def\csname p4num@p4_comp_b_class_r0_entry\endcsname{0.250000}
\expandafter\def\csname p4num@p4_comp_b_class_r0_own\endcsname{0.125000}
\expandafter\def\csname p4num@p4_comp_b_class_r0_proceeds\endcsname{0.392708}
\expandafter\def\csname p4num@p4_comp_b_class_r1_buy_rate\endcsname{1.000000}
\expandafter\def\csname p4num@p4_comp_b_class_r1_sell_rate\endcsname{1.000000}
\expandafter\def\csname p4num@p4_comp_b_class_r1_sell_mu\endcsname{0.268941}
\expandafter\def\csname p4num@p4_comp_b_class_r1_sell_nu\endcsname{0.268941}
\expandafter\def\csname p4num@p4_comp_b_class_r1_sell_price\endcsname{0.576074}
\expandafter\def\csname p4num@p4_comp_b_class_r1_sell_eh\endcsname{0.250000}
\expandafter\def\csname p4num@p4_comp_b_class_r1_sell_el\endcsname{0.250000}
\expandafter\def\csname p4num@p4_comp_b_class_r1_sell_yp_entry_gap\endcsname{-3.633821489}
\expandafter\def\csname p4num@p4_comp_b_class_r1_sell_ym_entry_gap\endcsname{-3.633821489}
\expandafter\def\csname p4num@p4_comp_b_class_r1_sell_tp_payoff\endcsname{-0.141843096}
\expandafter\def\csname p4num@p4_comp_b_class_r1_sell_tm_payoff\endcsname{0.024823570}
\expandafter\def\csname p4num@p4_comp_b_class_r1_zero_mu\endcsname{0.500000}
\expandafter\def\csname p4num@p4_comp_b_class_r1_zero_nu\endcsname{0.500000}
\expandafter\def\csname p4num@p4_comp_b_class_r1_zero_price\endcsname{0.614583}
\expandafter\def\csname p4num@p4_comp_b_class_r1_zero_eh\endcsname{0.250000}
\expandafter\def\csname p4num@p4_comp_b_class_r1_zero_el\endcsname{0.250000}
\expandafter\def\csname p4num@p4_comp_b_class_r1_zero_yp_entry_gap\endcsname{-1.708333333}
\expandafter\def\csname p4num@p4_comp_b_class_r1_zero_ym_entry_gap\endcsname{-1.708333333}
\expandafter\def\csname p4num@p4_comp_b_class_r1_zero_tp_payoff\endcsname{0.000000000}
\expandafter\def\csname p4num@p4_comp_b_class_r1_zero_tm_payoff\endcsname{0.000000000}
\expandafter\def\csname p4num@p4_comp_b_class_r1_buy_mu\endcsname{0.731059}
\expandafter\def\csname p4num@p4_comp_b_class_r1_buy_nu\endcsname{0.731059}
\expandafter\def\csname p4num@p4_comp_b_class_r1_buy_price\endcsname{1.362372}
\expandafter\def\csname p4num@p4_comp_b_class_r1_buy_eh\endcsname{1.000000}
\expandafter\def\csname p4num@p4_comp_b_class_r1_buy_el\endcsname{1.000000}
\expandafter\def\csname p4num@p4_comp_b_class_r1_buy_yp_entry_gap\endcsname{0.217154822}
\expandafter\def\csname p4num@p4_comp_b_class_r1_buy_ym_entry_gap\endcsname{0.217154822}
\expandafter\def\csname p4num@p4_comp_b_class_r1_buy_tp_payoff\endcsname{0.159294281}
\expandafter\def\csname p4num@p4_comp_b_class_r1_buy_tm_payoff\endcsname{-0.507372386}
\expandafter\def\csname p4num@p4_comp_b_class_r1_entry\endcsname{0.522757}
\expandafter\def\csname p4num@p4_comp_b_class_r1_own\endcsname{0.324401}
\expandafter\def\csname p4num@p4_comp_b_class_r1_proceeds\endcsname{0.872531}
\expandafter\def\csname p4num@p4_valid_b_no_y_r0_buy_rate\endcsname{0.000000}
\expandafter\def\csname p4num@p4_valid_b_no_y_r0_sell_rate\endcsname{0.000000}
\expandafter\def\csname p4num@p4_valid_b_no_y_r0_sell_mu\endcsname{0.500000}
\expandafter\def\csname p4num@p4_valid_b_no_y_r0_sell_nu\endcsname{0.500000}
\expandafter\def\csname p4num@p4_valid_b_no_y_r0_sell_price\endcsname{0.392708}
\expandafter\def\csname p4num@p4_valid_b_no_y_r0_sell_eh\endcsname{0.250000}
\expandafter\def\csname p4num@p4_valid_b_no_y_r0_sell_el\endcsname{0.250000}
\expandafter\def\csname p4num@p4_valid_b_no_y_r0_sell_yp_entry_gap\endcsname{-1.195833333}
\expandafter\def\csname p4num@p4_valid_b_no_y_r0_sell_ym_entry_gap\endcsname{-1.195833333}
\expandafter\def\csname p4num@p4_valid_b_no_y_r0_sell_tp_payoff\endcsname{-0.021875000}
\expandafter\def\csname p4num@p4_valid_b_no_y_r0_sell_tm_payoff\endcsname{-0.018125000}
\expandafter\def\csname p4num@p4_valid_b_no_y_r0_zero_mu\endcsname{0.500000}
\expandafter\def\csname p4num@p4_valid_b_no_y_r0_zero_nu\endcsname{0.500000}
\expandafter\def\csname p4num@p4_valid_b_no_y_r0_zero_price\endcsname{0.392708}
\expandafter\def\csname p4num@p4_valid_b_no_y_r0_zero_eh\endcsname{0.250000}
\expandafter\def\csname p4num@p4_valid_b_no_y_r0_zero_el\endcsname{0.250000}
\expandafter\def\csname p4num@p4_valid_b_no_y_r0_zero_yp_entry_gap\endcsname{-1.195833333}
\expandafter\def\csname p4num@p4_valid_b_no_y_r0_zero_ym_entry_gap\endcsname{-1.195833333}
\expandafter\def\csname p4num@p4_valid_b_no_y_r0_zero_tp_payoff\endcsname{0.000000000}
\expandafter\def\csname p4num@p4_valid_b_no_y_r0_zero_tm_payoff\endcsname{0.000000000}
\expandafter\def\csname p4num@p4_valid_b_no_y_r0_buy_mu\endcsname{0.500000}
\expandafter\def\csname p4num@p4_valid_b_no_y_r0_buy_nu\endcsname{0.500000}
\expandafter\def\csname p4num@p4_valid_b_no_y_r0_buy_price\endcsname{0.392708}
\expandafter\def\csname p4num@p4_valid_b_no_y_r0_buy_eh\endcsname{0.250000}
\expandafter\def\csname p4num@p4_valid_b_no_y_r0_buy_el\endcsname{0.250000}
\expandafter\def\csname p4num@p4_valid_b_no_y_r0_buy_yp_entry_gap\endcsname{-1.195833333}
\expandafter\def\csname p4num@p4_valid_b_no_y_r0_buy_ym_entry_gap\endcsname{-1.195833333}
\expandafter\def\csname p4num@p4_valid_b_no_y_r0_buy_tp_payoff\endcsname{-0.018125000}
\expandafter\def\csname p4num@p4_valid_b_no_y_r0_buy_tm_payoff\endcsname{-0.021875000}
\expandafter\def\csname p4num@p4_valid_b_no_y_r0_entry\endcsname{0.250000}
\expandafter\def\csname p4num@p4_valid_b_no_y_r0_own\endcsname{0.125000}
\expandafter\def\csname p4num@p4_valid_b_no_y_r0_proceeds\endcsname{0.392708}
\expandafter\def\csname p4num@p4_valid_b_no_y_r1_buy_rate\endcsname{1.000000}
\expandafter\def\csname p4num@p4_valid_b_no_y_r1_sell_rate\endcsname{1.000000}
\expandafter\def\csname p4num@p4_valid_b_no_y_r1_sell_mu\endcsname{0.292047}
\expandafter\def\csname p4num@p4_valid_b_no_y_r1_sell_nu\endcsname{0.268941}
\expandafter\def\csname p4num@p4_valid_b_no_y_r1_sell_price\endcsname{0.579925}
\expandafter\def\csname p4num@p4_valid_b_no_y_r1_sell_eh\endcsname{0.250000}
\expandafter\def\csname p4num@p4_valid_b_no_y_r1_sell_el\endcsname{0.250000}
\expandafter\def\csname p4num@p4_valid_b_no_y_r1_sell_yp_entry_gap\endcsname{-3.441272673}
\expandafter\def\csname p4num@p4_valid_b_no_y_r1_sell_ym_entry_gap\endcsname{-3.441272673}
\expandafter\def\csname p4num@p4_valid_b_no_y_r1_sell_tp_payoff\endcsname{-0.129658787}
\expandafter\def\csname p4num@p4_valid_b_no_y_r1_sell_tm_payoff\endcsname{0.020341213}
\expandafter\def\csname p4num@p4_valid_b_no_y_r1_zero_mu\endcsname{0.500000}
\expandafter\def\csname p4num@p4_valid_b_no_y_r1_zero_nu\endcsname{0.500000}
\expandafter\def\csname p4num@p4_valid_b_no_y_r1_zero_price\endcsname{0.614583}
\expandafter\def\csname p4num@p4_valid_b_no_y_r1_zero_eh\endcsname{0.250000}
\expandafter\def\csname p4num@p4_valid_b_no_y_r1_zero_el\endcsname{0.250000}
\expandafter\def\csname p4num@p4_valid_b_no_y_r1_zero_yp_entry_gap\endcsname{-1.708333333}
\expandafter\def\csname p4num@p4_valid_b_no_y_r1_zero_ym_entry_gap\endcsname{-1.708333333}
\expandafter\def\csname p4num@p4_valid_b_no_y_r1_zero_tp_payoff\endcsname{0.000000000}
\expandafter\def\csname p4num@p4_valid_b_no_y_r1_zero_tm_payoff\endcsname{0.000000000}
\expandafter\def\csname p4num@p4_valid_b_no_y_r1_buy_mu\endcsname{0.707953}
\expandafter\def\csname p4num@p4_valid_b_no_y_r1_buy_nu\endcsname{0.731059}
\expandafter\def\csname p4num@p4_valid_b_no_y_r1_buy_price\endcsname{1.346968}
\expandafter\def\csname p4num@p4_valid_b_no_y_r1_buy_eh\endcsname{1.000000}
\expandafter\def\csname p4num@p4_valid_b_no_y_r1_buy_el\endcsname{1.000000}
\expandafter\def\csname p4num@p4_valid_b_no_y_r1_buy_yp_entry_gap\endcsname{0.024606006}
\expandafter\def\csname p4num@p4_valid_b_no_y_r1_buy_ym_entry_gap\endcsname{0.024606006}
\expandafter\def\csname p4num@p4_valid_b_no_y_r1_buy_tp_payoff\endcsname{0.141364853}
\expandafter\def\csname p4num@p4_valid_b_no_y_r1_buy_tm_payoff\endcsname{-0.458635147}
\expandafter\def\csname p4num@p4_valid_b_no_y_r1_entry\endcsname{0.522757}
\expandafter\def\csname p4num@p4_valid_b_no_y_r1_own\endcsname{0.318099}
\expandafter\def\csname p4num@p4_valid_b_no_y_r1_proceeds\endcsname{0.868330}
\expandafter\def\csname p4num@margin_p4_comp_b_no_y_weak\endcsname{0.005333333}
\expandafter\def\csname p4num@margin_p4_comp_b_no_y_strong\endcsname{0.019444742}
\expandafter\def\csname p4num@margin_p4_comp_b_no_y_prior_r0\endcsname{1.195833333}
\expandafter\def\csname p4num@margin_p4_comp_b_no_y_top_r1\endcsname{0.013903757}
\expandafter\def\csname p4num@margin_p4_comp_b_class_weak\endcsname{0.003333333}
\expandafter\def\csname p4num@margin_p4_comp_b_class_strong\endcsname{0.024823570}
\expandafter\def\csname p4num@margin_p4_comp_b_class_prior_r0\endcsname{1.195833333}
\expandafter\def\csname p4num@margin_p4_comp_b_class_prior_r1\endcsname{1.708333333}
\expandafter\def\csname p4num@margin_p4_comp_b_class_top_r1\endcsname{0.217154822}
\expandafter\def\csname p4num@margin_p4_valid_b_no_y_weak\endcsname{0.005000000}
\expandafter\def\csname p4num@margin_p4_valid_b_no_y_strong\endcsname{0.020341213}
\expandafter\def\csname p4num@margin_p4_valid_b_no_y_prior_r0\endcsname{1.195833333}
\expandafter\def\csname p4num@margin_p4_valid_b_no_y_prior_r1\endcsname{1.708333333}
\expandafter\def\csname p4num@margin_p4_valid_b_no_y_top_r1\endcsname{0.024606006}
\expandafter\def\csname p4num@margin_p4_valid_b_no_y_top_exclusion\endcsname{-0.024606006}
\expandafter\def\csname p4num@p4_b_signal_r0_buy_rate\endcsname{0.721409}
\expandafter\def\csname p4num@p4_b_signal_r0_sell_rate\endcsname{0.721409}
\expandafter\def\csname p4num@p4_b_signal_r0_sell_mu\endcsname{0.423473}
\expandafter\def\csname p4num@p4_b_signal_r0_sell_nu\endcsname{0.308682}
\expandafter\def\csname p4num@p4_b_signal_r0_sell_price\endcsname{0.533438}
\expandafter\def\csname p4num@p4_b_signal_r0_sell_eh\endcsname{0.812500}
\expandafter\def\csname p4num@p4_b_signal_r0_sell_el\endcsname{0.437500}
\expandafter\def\csname p4num@p4_b_signal_r0_sell_yp_entry_gap\endcsname{0.491669568}
\expandafter\def\csname p4num@p4_b_signal_r0_sell_ym_entry_gap\endcsname{-3.920616786}
\expandafter\def\csname p4num@p4_b_signal_r0_sell_tp_payoff\endcsname{-0.064791667}
\expandafter\def\csname p4num@p4_b_signal_r0_sell_tm_payoff\endcsname{0.000000000}
\expandafter\def\csname p4num@p4_b_signal_r0_zero_mu\endcsname{0.500000}
\expandafter\def\csname p4num@p4_b_signal_r0_zero_nu\endcsname{0.500000}
\expandafter\def\csname p4num@p4_b_signal_r0_zero_price\endcsname{0.545833}
\expandafter\def\csname p4num@p4_b_signal_r0_zero_eh\endcsname{0.812500}
\expandafter\def\csname p4num@p4_b_signal_r0_zero_el\endcsname{0.437500}
\expandafter\def\csname p4num@p4_b_signal_r0_zero_yp_entry_gap\endcsname{1.050000000}
\expandafter\def\csname p4num@p4_b_signal_r0_zero_ym_entry_gap\endcsname{-3.441666667}
\expandafter\def\csname p4num@p4_b_signal_r0_zero_tp_payoff\endcsname{0.000000000}
\expandafter\def\csname p4num@p4_b_signal_r0_zero_tm_payoff\endcsname{0.000000000}
\expandafter\def\csname p4num@p4_b_signal_r0_buy_mu\endcsname{0.576527}
\expandafter\def\csname p4num@p4_b_signal_r0_buy_nu\endcsname{0.691318}
\expandafter\def\csname p4num@p4_b_signal_r0_buy_price\endcsname{0.558229}
\expandafter\def\csname p4num@p4_b_signal_r0_buy_eh\endcsname{0.812500}
\expandafter\def\csname p4num@p4_b_signal_r0_buy_el\endcsname{0.437500}
\expandafter\def\csname p4num@p4_b_signal_r0_buy_yp_entry_gap\endcsname{1.528950119}
\expandafter\def\csname p4num@p4_b_signal_r0_buy_ym_entry_gap\endcsname{-2.883336235}
\expandafter\def\csname p4num@p4_b_signal_r0_buy_tp_payoff\endcsname{-0.000000000}
\expandafter\def\csname p4num@p4_b_signal_r0_buy_tm_payoff\endcsname{-0.064791667}
\expandafter\def\csname p4num@p4_b_signal_r0_entry\endcsname{0.625000}
\expandafter\def\csname p4num@p4_b_signal_r0_own\endcsname{0.406250}
\expandafter\def\csname p4num@p4_b_signal_r0_proceeds\endcsname{0.545833}
\expandafter\def\csname p4num@p4_b_signal_r1_buy_rate\endcsname{1.000000}
\expandafter\def\csname p4num@p4_b_signal_r1_sell_rate\endcsname{0.775969}
\expandafter\def\csname p4num@p4_b_signal_r1_sell_mu\endcsname{0.420000}
\expandafter\def\csname p4num@p4_b_signal_r1_sell_nu\endcsname{0.300000}
\expandafter\def\csname p4num@p4_b_signal_r1_sell_price\endcsname{0.601250}
\expandafter\def\csname p4num@p4_b_signal_r1_sell_eh\endcsname{0.250000}
\expandafter\def\csname p4num@p4_b_signal_r1_sell_el\endcsname{0.250000}
\expandafter\def\csname p4num@p4_b_signal_r1_sell_yp_entry_gap\endcsname{-0.168478261}
\expandafter\def\csname p4num@p4_b_signal_r1_sell_ym_entry_gap\endcsname{-4.254629630}
\expandafter\def\csname p4num@p4_b_signal_r1_sell_tp_payoff\endcsname{-0.066666667}
\expandafter\def\csname p4num@p4_b_signal_r1_sell_tm_payoff\endcsname{0.000000000}
\expandafter\def\csname p4num@p4_b_signal_r1_zero_mu\endcsname{0.475733}
\expandafter\def\csname p4num@p4_b_signal_r1_zero_nu\endcsname{0.439331}
\expandafter\def\csname p4num@p4_b_signal_r1_zero_price\endcsname{0.956642}
\expandafter\def\csname p4num@p4_b_signal_r1_zero_eh\endcsname{0.812500}
\expandafter\def\csname p4num@p4_b_signal_r1_zero_el\endcsname{0.437500}
\expandafter\def\csname p4num@p4_b_signal_r1_zero_yp_entry_gap\endcsname{0.219556089}
\expandafter\def\csname p4num@p4_b_signal_r1_zero_ym_entry_gap\endcsname{-3.939744864}
\expandafter\def\csname p4num@p4_b_signal_r1_zero_tp_payoff\endcsname{0.000000000}
\expandafter\def\csname p4num@p4_b_signal_r1_zero_tm_payoff\endcsname{0.000000000}
\expandafter\def\csname p4num@p4_b_signal_r1_buy_mu\endcsname{0.592423}
\expandafter\def\csname p4num@p4_b_signal_r1_buy_nu\endcsname{0.731059}
\expandafter\def\csname p4num@p4_b_signal_r1_buy_price\endcsname{1.039906}
\expandafter\def\csname p4num@p4_b_signal_r1_buy_eh\endcsname{0.812500}
\expandafter\def\csname p4num@p4_b_signal_r1_buy_el\endcsname{0.437500}
\expandafter\def\csname p4num@p4_b_signal_r1_buy_yp_entry_gap\endcsname{0.903775226}
\expandafter\def\csname p4num@p4_b_signal_r1_buy_ym_entry_gap\endcsname{-3.155195358}
\expandafter\def\csname p4num@p4_b_signal_r1_buy_tp_payoff\endcsname{0.056760364}
\expandafter\def\csname p4num@p4_b_signal_r1_buy_tm_payoff\endcsname{-0.228656303}
\expandafter\def\csname p4num@p4_b_signal_r1_entry\endcsname{0.512521}
\expandafter\def\csname p4num@p4_b_signal_r1_own\endcsname{0.329227}
\expandafter\def\csname p4num@p4_b_signal_r1_proceeds\endcsname{0.871057}
\expandafter\def\csname p4num@p4_b_no_y_r0_buy_rate\endcsname{0.000000}
\expandafter\def\csname p4num@p4_b_no_y_r0_sell_rate\endcsname{0.000000}
\expandafter\def\csname p4num@p4_b_no_y_r0_sell_mu\endcsname{0.500000}
\expandafter\def\csname p4num@p4_b_no_y_r0_sell_nu\endcsname{0.500000}
\expandafter\def\csname p4num@p4_b_no_y_r0_sell_price\endcsname{0.392708}
\expandafter\def\csname p4num@p4_b_no_y_r0_sell_eh\endcsname{0.250000}
\expandafter\def\csname p4num@p4_b_no_y_r0_sell_el\endcsname{0.250000}
\expandafter\def\csname p4num@p4_b_no_y_r0_sell_yp_entry_gap\endcsname{-1.195833333}
\expandafter\def\csname p4num@p4_b_no_y_r0_sell_ym_entry_gap\endcsname{-1.195833333}
\expandafter\def\csname p4num@p4_b_no_y_r0_sell_tp_payoff\endcsname{-0.020833333}
\expandafter\def\csname p4num@p4_b_no_y_r0_sell_tm_payoff\endcsname{-0.019166667}
\expandafter\def\csname p4num@p4_b_no_y_r0_zero_mu\endcsname{0.500000}
\expandafter\def\csname p4num@p4_b_no_y_r0_zero_nu\endcsname{0.500000}
\expandafter\def\csname p4num@p4_b_no_y_r0_zero_price\endcsname{0.392708}
\expandafter\def\csname p4num@p4_b_no_y_r0_zero_eh\endcsname{0.250000}
\expandafter\def\csname p4num@p4_b_no_y_r0_zero_el\endcsname{0.250000}
\expandafter\def\csname p4num@p4_b_no_y_r0_zero_yp_entry_gap\endcsname{-1.195833333}
\expandafter\def\csname p4num@p4_b_no_y_r0_zero_ym_entry_gap\endcsname{-1.195833333}
\expandafter\def\csname p4num@p4_b_no_y_r0_zero_tp_payoff\endcsname{0.000000000}
\expandafter\def\csname p4num@p4_b_no_y_r0_zero_tm_payoff\endcsname{0.000000000}
\expandafter\def\csname p4num@p4_b_no_y_r0_buy_mu\endcsname{0.500000}
\expandafter\def\csname p4num@p4_b_no_y_r0_buy_nu\endcsname{0.500000}
\expandafter\def\csname p4num@p4_b_no_y_r0_buy_price\endcsname{0.392708}
\expandafter\def\csname p4num@p4_b_no_y_r0_buy_eh\endcsname{0.250000}
\expandafter\def\csname p4num@p4_b_no_y_r0_buy_el\endcsname{0.250000}
\expandafter\def\csname p4num@p4_b_no_y_r0_buy_yp_entry_gap\endcsname{-1.195833333}
\expandafter\def\csname p4num@p4_b_no_y_r0_buy_ym_entry_gap\endcsname{-1.195833333}
\expandafter\def\csname p4num@p4_b_no_y_r0_buy_tp_payoff\endcsname{-0.019166667}
\expandafter\def\csname p4num@p4_b_no_y_r0_buy_tm_payoff\endcsname{-0.020833333}
\expandafter\def\csname p4num@p4_b_no_y_r0_entry\endcsname{0.250000}
\expandafter\def\csname p4num@p4_b_no_y_r0_own\endcsname{0.125000}
\expandafter\def\csname p4num@p4_b_no_y_r0_proceeds\endcsname{0.392708}
\expandafter\def\csname p4num@p4_b_no_y_r1_buy_rate\endcsname{0.775969}
\expandafter\def\csname p4num@p4_b_no_y_r1_sell_rate\endcsname{0.775969}
\expandafter\def\csname p4num@p4_b_no_y_r1_sell_mu\endcsname{0.420000}
\expandafter\def\csname p4num@p4_b_no_y_r1_sell_nu\endcsname{0.300000}
\expandafter\def\csname p4num@p4_b_no_y_r1_sell_price\endcsname{0.601250}
\expandafter\def\csname p4num@p4_b_no_y_r1_sell_eh\endcsname{0.250000}
\expandafter\def\csname p4num@p4_b_no_y_r1_sell_el\endcsname{0.250000}
\expandafter\def\csname p4num@p4_b_no_y_r1_sell_yp_entry_gap\endcsname{-2.375000000}
\expandafter\def\csname p4num@p4_b_no_y_r1_sell_ym_entry_gap\endcsname{-2.375000000}
\expandafter\def\csname p4num@p4_b_no_y_r1_sell_tp_payoff\endcsname{-0.066666667}
\expandafter\def\csname p4num@p4_b_no_y_r1_sell_tm_payoff\endcsname{0.000000000}
\expandafter\def\csname p4num@p4_b_no_y_r1_zero_mu\endcsname{0.500000}
\expandafter\def\csname p4num@p4_b_no_y_r1_zero_nu\endcsname{0.500000}
\expandafter\def\csname p4num@p4_b_no_y_r1_zero_price\endcsname{0.614583}
\expandafter\def\csname p4num@p4_b_no_y_r1_zero_eh\endcsname{0.250000}
\expandafter\def\csname p4num@p4_b_no_y_r1_zero_el\endcsname{0.250000}
\expandafter\def\csname p4num@p4_b_no_y_r1_zero_yp_entry_gap\endcsname{-1.708333333}
\expandafter\def\csname p4num@p4_b_no_y_r1_zero_ym_entry_gap\endcsname{-1.708333333}
\expandafter\def\csname p4num@p4_b_no_y_r1_zero_tp_payoff\endcsname{0.000000000}
\expandafter\def\csname p4num@p4_b_no_y_r1_zero_tm_payoff\endcsname{0.000000000}
\expandafter\def\csname p4num@p4_b_no_y_r1_buy_mu\endcsname{0.580000}
\expandafter\def\csname p4num@p4_b_no_y_r1_buy_nu\endcsname{0.700000}
\expandafter\def\csname p4num@p4_b_no_y_r1_buy_price\endcsname{0.627917}
\expandafter\def\csname p4num@p4_b_no_y_r1_buy_eh\endcsname{0.250000}
\expandafter\def\csname p4num@p4_b_no_y_r1_buy_el\endcsname{0.250000}
\expandafter\def\csname p4num@p4_b_no_y_r1_buy_yp_entry_gap\endcsname{-1.041666667}
\expandafter\def\csname p4num@p4_b_no_y_r1_buy_ym_entry_gap\endcsname{-1.041666667}
\expandafter\def\csname p4num@p4_b_no_y_r1_buy_tp_payoff\endcsname{-0.000000000}
\expandafter\def\csname p4num@p4_b_no_y_r1_buy_tm_payoff\endcsname{-0.066666667}
\expandafter\def\csname p4num@p4_b_no_y_r1_entry\endcsname{0.250000}
\expandafter\def\csname p4num@p4_b_no_y_r1_own\endcsname{0.125000}
\expandafter\def\csname p4num@p4_b_no_y_r1_proceeds\endcsname{0.614583}
\expandafter\def\csname p4num@p4_b_no_y_r1_zero_gross\endcsname{0.033333333}
\expandafter\def\csname p4num@p4_b_no_y_r1_full_gross\endcsname{0.017929428}
\expandafter\def\csname p4num@p4_b_signal_entry_change\endcsname{-0.112479}
\expandafter\def\csname p4num@p4_b_signal_own_change\endcsname{-0.077023}
\expandafter\def\csname p4num@margin_p4_b_weak_all_yplus\endcsname{0.363883929}
\expandafter\def\csname p4num@margin_p4_b_weak_no_yminus\endcsname{2.755550596}
\expandafter\def\csname p4num@margin_p4_b_weak_zero_return\endcsname{0.012395833}
\expandafter\def\csname p4num@margin_p4_b_weak_full_loss\endcsname{0.002574837}
\expandafter\def\csname p4num@margin_p4_b_weak_rate_room\endcsname{0.478698495}
\expandafter\def\csname p4num@margin_p4_b_strong_buy\endcsname{0.056760364}
\expandafter\def\csname p4num@margin_p4_b_strong_prior_yplus\endcsname{0.375000000}
\expandafter\def\csname p4num@margin_p4_b_strong_no_yminus\endcsname{3.155195358}
\expandafter\def\csname p4num@margin_p4_b_strong_threshold_floor_return\endcsname{0.003899371}
\expandafter\def\csname p4num@margin_p4_b_strong_rate_room\endcsname{0.384948495}
\expandafter\def\csname p4num@margin_p4_b_strong_sale_public_gap\endcsname{0.023396226}
\expandafter\def\csname p4num@margin_p4_b_strong_zero_yplus\endcsname{0.219556089}
\expandafter\def\csname p4num@margin_p4_b_no_y_top_exclusion_r0\endcsname{0.365562841}
\expandafter\def\csname p4num@margin_p4_b_no_y_top_exclusion_r1\endcsname{0.938138071}
\expandafter\def\csname p4num@p4_accuracy_68\endcsname{0.68}
\expandafter\def\csname p4num@p4_accuracy_69\endcsname{0.69}
\expandafter\def\csname p4num@p4_accuracy_70\endcsname{0.70}
\expandafter\def\csname p4num@p4_accuracy_71\endcsname{0.71}
\expandafter\def\csname p4num@p4_accuracy_72\endcsname{0.72}
\expandafter\def\csname p4num@p4_accuracy_73\endcsname{0.73}
\expandafter\def\csname p4num@p4_accuracy_74\endcsname{0.74}
\expandafter\def\csname p4num@p4_accuracy_75\endcsname{0.75}
\expandafter\def\csname p4num@p4_accuracy_76\endcsname{0.76}
\expandafter\def\csname p4num@p4_accuracy_77\endcsname{0.77}
\expandafter\def\csname p4num@margin_p4_grid_68_73_private\endcsname{0.232409091}
\expandafter\def\csname p4num@margin_p4_grid_68_73_joint\endcsname{-0.149269935}
\expandafter\def\csname p4num@margin_p4_grid_68_73_weak\endcsname{0.003653455}
\expandafter\def\csname p4num@margin_p4_grid_68_73_strong\endcsname{0.015234862}
\expandafter\def\csname p4num@margin_p4_grid_68_73_other_pair\endcsname{4.033112080}
\expandafter\def\csname p4num@p4_grid_68_73_sufficient\endcsname{0}
\expandafter\def\csname p4num@margin_p4_grid_68_74_private\endcsname{0.142454545}
\expandafter\def\csname p4num@margin_p4_grid_68_74_joint\endcsname{-0.077065620}
\expandafter\def\csname p4num@margin_p4_grid_68_74_weak\endcsname{0.003231273}
\expandafter\def\csname p4num@margin_p4_grid_68_74_strong\endcsname{0.015234862}
\expandafter\def\csname p4num@margin_p4_grid_68_74_other_pair\endcsname{4.131892686}
\expandafter\def\csname p4num@p4_grid_68_74_sufficient\endcsname{0}
\expandafter\def\csname p4num@margin_p4_grid_68_75_private\endcsname{0.052500000}
\expandafter\def\csname p4num@margin_p4_grid_68_75_joint\endcsname{-0.005304891}
\expandafter\def\csname p4num@margin_p4_grid_68_75_weak\endcsname{0.002809091}
\expandafter\def\csname p4num@margin_p4_grid_68_75_strong\endcsname{0.015234862}
\expandafter\def\csname p4num@margin_p4_grid_68_75_other_pair\endcsname{4.231390265}
\expandafter\def\csname p4num@p4_grid_68_75_sufficient\endcsname{0}
\expandafter\def\csname p4num@margin_p4_grid_68_76_private\endcsname{-0.037454545}
\expandafter\def\csname p4num@margin_p4_grid_68_76_joint\endcsname{0.066016329}
\expandafter\def\csname p4num@margin_p4_grid_68_76_weak\endcsname{0.002386909}
\expandafter\def\csname p4num@margin_p4_grid_68_76_strong\endcsname{0.015234862}
\expandafter\def\csname p4num@margin_p4_grid_68_76_other_pair\endcsname{4.331612652}
\expandafter\def\csname p4num@p4_grid_68_76_sufficient\endcsname{0}
\expandafter\def\csname p4num@margin_p4_grid_68_77_private\endcsname{-0.127409091}
\expandafter\def\csname p4num@margin_p4_grid_68_77_joint\endcsname{0.136902065}
\expandafter\def\csname p4num@margin_p4_grid_68_77_weak\endcsname{0.001964727}
\expandafter\def\csname p4num@margin_p4_grid_68_77_strong\endcsname{0.015234862}
\expandafter\def\csname p4num@margin_p4_grid_68_77_other_pair\endcsname{4.432567795}
\expandafter\def\csname p4num@p4_grid_68_77_sufficient\endcsname{0}
\expandafter\def\csname p4num@margin_p4_grid_69_73_private\endcsname{0.232409091}
\expandafter\def\csname p4num@margin_p4_grid_69_73_joint\endcsname{-0.122237854}
\expandafter\def\csname p4num@margin_p4_grid_69_73_weak\endcsname{0.003023091}
\expandafter\def\csname p4num@margin_p4_grid_69_73_strong\endcsname{0.016914577}
\expandafter\def\csname p4num@margin_p4_grid_69_73_other_pair\endcsname{3.996061415}
\expandafter\def\csname p4num@p4_grid_69_73_sufficient\endcsname{0}
\expandafter\def\csname p4num@margin_p4_grid_69_74_private\endcsname{0.142454545}
\expandafter\def\csname p4num@margin_p4_grid_69_74_joint\endcsname{-0.050844653}
\expandafter\def\csname p4num@margin_p4_grid_69_74_weak\endcsname{0.002577455}
\expandafter\def\csname p4num@margin_p4_grid_69_74_strong\endcsname{0.016914577}
\expandafter\def\csname p4num@margin_p4_grid_69_74_other_pair\endcsname{4.095455527}
\expandafter\def\csname p4num@p4_grid_69_74_sufficient\endcsname{0}
\expandafter\def\csname p4num@margin_p4_grid_69_75_private\endcsname{0.052500000}
\expandafter\def\csname p4num@margin_p4_grid_69_75_joint\endcsname{0.020087547}
\expandafter\def\csname p4num@margin_p4_grid_69_75_weak\endcsname{0.002131818}
\expandafter\def\csname p4num@margin_p4_grid_69_75_strong\endcsname{0.016914577}
\expandafter\def\csname p4num@margin_p4_grid_69_75_other_pair\endcsname{4.195615001}
\expandafter\def\csname p4num@p4_grid_69_75_sufficient\endcsname{1}
\expandafter\def\csname p4num@p4_grid_69_75_entry\endcsname{0.879671}
\expandafter\def\csname p4num@p4_grid_69_75_own\endcsname{0.449049}
\expandafter\def\csname p4num@margin_p4_grid_69_76_private\endcsname{-0.037454545}
\expandafter\def\csname p4num@margin_p4_grid_69_76_joint\endcsname{0.090563195}
\expandafter\def\csname p4num@margin_p4_grid_69_76_weak\endcsname{0.001686182}
\expandafter\def\csname p4num@margin_p4_grid_69_76_strong\endcsname{0.016914577}
\expandafter\def\csname p4num@margin_p4_grid_69_76_other_pair\endcsname{4.296548712}
\expandafter\def\csname p4num@p4_grid_69_76_sufficient\endcsname{0}
\expandafter\def\csname p4num@margin_p4_grid_69_77_private\endcsname{-0.127409091}
\expandafter\def\csname p4num@margin_p4_grid_69_77_joint\endcsname{0.160586687}
\expandafter\def\csname p4num@margin_p4_grid_69_77_weak\endcsname{0.001240545}
\expandafter\def\csname p4num@margin_p4_grid_69_77_strong\endcsname{0.016914577}
\expandafter\def\csname p4num@margin_p4_grid_69_77_other_pair\endcsname{4.398265673}
\expandafter\def\csname p4num@p4_grid_69_77_sufficient\endcsname{0}
\expandafter\def\csname p4num@margin_p4_grid_70_73_private\endcsname{0.232409091}
\expandafter\def\csname p4num@margin_p4_grid_70_73_joint\endcsname{-0.095417614}
\expandafter\def\csname p4num@margin_p4_grid_70_73_weak\endcsname{0.002392727}
\expandafter\def\csname p4num@margin_p4_grid_70_73_strong\endcsname{0.018594291}
\expandafter\def\csname p4num@margin_p4_grid_70_73_other_pair\endcsname{3.958666433}
\expandafter\def\csname p4num@p4_grid_70_73_sufficient\endcsname{0}
\expandafter\def\csname p4num@margin_p4_grid_70_74_private\endcsname{0.142454545}
\expandafter\def\csname p4num@margin_p4_grid_70_74_joint\endcsname{-0.024837375}
\expandafter\def\csname p4num@margin_p4_grid_70_74_weak\endcsname{0.001923636}
\expandafter\def\csname p4num@margin_p4_grid_70_74_strong\endcsname{0.018594291}
\expandafter\def\csname p4num@margin_p4_grid_70_74_other_pair\endcsname{4.058663595}
\expandafter\def\csname p4num@p4_grid_70_74_sufficient\endcsname{0}
\expandafter\def\csname p4num@margin_p4_grid_70_75_private\endcsname{0.052500000}
\expandafter\def\csname p4num@margin_p4_grid_70_75_joint\endcsname{0.045265153}
\expandafter\def\csname p4num@margin_p4_grid_70_75_weak\endcsname{0.001454545}
\expandafter\def\csname p4num@margin_p4_grid_70_75_strong\endcsname{0.018594291}
\expandafter\def\csname p4num@margin_p4_grid_70_75_other_pair\endcsname{4.159475417}
\expandafter\def\csname p4num@p4_grid_70_75_sufficient\endcsname{1}
\expandafter\def\csname p4num@p4_grid_70_75_entry\endcsname{0.879797}
\expandafter\def\csname p4num@p4_grid_70_75_own\endcsname{0.449238}
\expandafter\def\csname p4num@margin_p4_grid_70_76_private\endcsname{-0.037454545}
\expandafter\def\csname p4num@margin_p4_grid_70_76_joint\endcsname{0.114894806}
\expandafter\def\csname p4num@margin_p4_grid_70_76_weak\endcsname{0.000985455}
\expandafter\def\csname p4num@margin_p4_grid_70_76_strong\endcsname{0.018594291}
\expandafter\def\csname p4num@margin_p4_grid_70_76_other_pair\endcsname{4.261111895}
\expandafter\def\csname p4num@p4_grid_70_76_sufficient\endcsname{0}
\expandafter\def\csname p4num@margin_p4_grid_70_77_private\endcsname{-0.127409091}
\expandafter\def\csname p4num@margin_p4_grid_70_77_joint\endcsname{0.184056351}
\expandafter\def\csname p4num@margin_p4_grid_70_77_weak\endcsname{0.000516364}
\expandafter\def\csname p4num@margin_p4_grid_70_77_strong\endcsname{0.018594291}
\expandafter\def\csname p4num@margin_p4_grid_70_77_other_pair\endcsname{4.363583190}
\expandafter\def\csname p4num@p4_grid_70_77_sufficient\endcsname{0}
\expandafter\def\csname p4num@margin_p4_grid_71_73_private\endcsname{0.232409091}
\expandafter\def\csname p4num@margin_p4_grid_71_73_joint\endcsname{-0.068806733}
\expandafter\def\csname p4num@margin_p4_grid_71_73_weak\endcsname{0.001762364}
\expandafter\def\csname p4num@margin_p4_grid_71_73_strong\endcsname{0.020274006}
\expandafter\def\csname p4num@margin_p4_grid_71_73_other_pair\endcsname{3.920922312}
\expandafter\def\csname p4num@p4_grid_71_73_sufficient\endcsname{0}
\expandafter\def\csname p4num@margin_p4_grid_71_74_private\endcsname{0.142454545}
\expandafter\def\csname p4num@margin_p4_grid_71_74_joint\endcsname{0.000958814}
\expandafter\def\csname p4num@margin_p4_grid_71_74_weak\endcsname{0.001269818}
\expandafter\def\csname p4num@margin_p4_grid_71_74_strong\endcsname{0.020274006}
\expandafter\def\csname p4num@margin_p4_grid_71_74_other_pair\endcsname{4.021511685}
\expandafter\def\csname p4num@p4_grid_71_74_sufficient\endcsname{1}
\expandafter\def\csname p4num@p4_grid_71_74_entry\endcsname{0.879817}
\expandafter\def\csname p4num@p4_grid_71_74_own\endcsname{0.449102}
\expandafter\def\csname p4num@margin_p4_grid_71_75_private\endcsname{0.052500000}
\expandafter\def\csname p4num@margin_p4_grid_71_75_joint\endcsname{0.070230645}
\expandafter\def\csname p4num@margin_p4_grid_71_75_weak\endcsname{0.000777273}
\expandafter\def\csname p4num@margin_p4_grid_71_75_strong\endcsname{0.020274006}
\expandafter\def\csname p4num@margin_p4_grid_71_75_other_pair\endcsname{4.122965921}
\expandafter\def\csname p4num@p4_grid_71_75_sufficient\endcsname{1}
\expandafter\def\csname p4num@p4_grid_71_75_entry\endcsname{0.879923}
\expandafter\def\csname p4num@p4_grid_71_75_own\endcsname{0.449427}
\expandafter\def\csname p4num@margin_p4_grid_71_76_private\endcsname{-0.037454545}
\expandafter\def\csname p4num@margin_p4_grid_71_76_joint\endcsname{0.139013981}
\expandafter\def\csname p4num@margin_p4_grid_71_76_weak\endcsname{0.000284727}
\expandafter\def\csname p4num@margin_p4_grid_71_76_strong\endcsname{0.020274006}
\expandafter\def\csname p4num@margin_p4_grid_71_76_other_pair\endcsname{4.225296222}
\expandafter\def\csname p4num@p4_grid_71_76_sufficient\endcsname{0}
\expandafter\def\csname p4num@margin_p4_grid_71_77_private\endcsname{-0.127409091}
\expandafter\def\csname p4num@margin_p4_grid_71_77_joint\endcsname{0.207313971}
\expandafter\def\csname p4num@margin_p4_grid_71_77_weak\endcsname{-0.000207818}
\expandafter\def\csname p4num@margin_p4_grid_71_77_strong\endcsname{0.020274006}
\expandafter\def\csname p4num@margin_p4_grid_71_77_other_pair\endcsname{4.328513984}
\expandafter\def\csname p4num@p4_grid_71_77_sufficient\endcsname{0}
\expandafter\def\csname p4num@margin_p4_grid_72_73_private\endcsname{0.232409091}
\expandafter\def\csname p4num@margin_p4_grid_72_73_joint\endcsname{-0.042402770}
\expandafter\def\csname p4num@margin_p4_grid_72_73_weak\endcsname{0.001132000}
\expandafter\def\csname p4num@margin_p4_grid_72_73_strong\endcsname{0.021953721}
\expandafter\def\csname p4num@margin_p4_grid_72_73_other_pair\endcsname{3.882824139}
\expandafter\def\csname p4num@p4_grid_72_73_sufficient\endcsname{0}
\expandafter\def\csname p4num@margin_p4_grid_72_74_private\endcsname{0.142454545}
\expandafter\def\csname p4num@margin_p4_grid_72_74_joint\endcsname{0.026546476}
\expandafter\def\csname p4num@margin_p4_grid_72_74_weak\endcsname{0.000616000}
\expandafter\def\csname p4num@margin_p4_grid_72_74_strong\endcsname{0.021953721}
\expandafter\def\csname p4num@margin_p4_grid_72_74_other_pair\endcsname{3.983994488}
\expandafter\def\csname p4num@p4_grid_72_74_sufficient\endcsname{1}
\expandafter\def\csname p4num@p4_grid_72_74_entry\endcsname{0.879938}
\expandafter\def\csname p4num@p4_grid_72_74_own\endcsname{0.449288}
\expandafter\def\csname p4num@margin_p4_grid_72_75_private\endcsname{0.052500000}
\expandafter\def\csname p4num@margin_p4_grid_72_75_joint\endcsname{0.094986690}
\expandafter\def\csname p4num@margin_p4_grid_72_75_weak\endcsname{0.000100000}
\expandafter\def\csname p4num@margin_p4_grid_72_75_strong\endcsname{0.021953721}
\expandafter\def\csname p4num@margin_p4_grid_72_75_other_pair\endcsname{4.086080804}
\expandafter\def\csname p4num@p4_grid_72_75_sufficient\endcsname{1}
\expandafter\def\csname p4num@p4_grid_72_75_entry\endcsname{0.880049}
\expandafter\def\csname p4num@p4_grid_72_75_own\endcsname{0.449616}
\expandafter\def\csname p4num@margin_p4_grid_72_76_private\endcsname{-0.037454545}
\expandafter\def\csname p4num@margin_p4_grid_72_76_joint\endcsname{0.162923490}
\expandafter\def\csname p4num@margin_p4_grid_72_76_weak\endcsname{-0.000416000}
\expandafter\def\csname p4num@margin_p4_grid_72_76_strong\endcsname{0.021953721}
\expandafter\def\csname p4num@margin_p4_grid_72_76_other_pair\endcsname{4.189095583}
\expandafter\def\csname p4num@p4_grid_72_76_sufficient\endcsname{0}
\expandafter\def\csname p4num@margin_p4_grid_72_77_private\endcsname{-0.127409091}
\expandafter\def\csname p4num@margin_p4_grid_72_77_joint\endcsname{0.230362408}
\expandafter\def\csname p4num@margin_p4_grid_72_77_weak\endcsname{-0.000932000}
\expandafter\def\csname p4num@margin_p4_grid_72_77_strong\endcsname{0.021953721}
\expandafter\def\csname p4num@margin_p4_grid_72_77_other_pair\endcsname{4.293051551}
\expandafter\def\csname p4num@p4_grid_72_77_sufficient\endcsname{0}
\expandafter\def\csname p4num@p4_grid_met_count\endcsname{6}
\expandafter\def\csname p4num@p4_b_signal_z_min\endcsname{0.446193819}
\expandafter\def\csname p4num@margin_p4_b_d_ge_a_impossible\endcsname{0.049029700}
\expandafter\def\csname p4num@margin_p4_min_domain_b\endcsname{0.020000000}
\expandafter\def\csname p4num@margin_p4_min_domain_s\endcsname{0.015000000}
\expandafter\def\csname p4num@margin_p4_min_domain_valid\endcsname{0.005000000}
\expandafter\def\csname p4num@margin_p4_min_domain_comp\endcsname{0.005000000}
\expandafter\def\csname p4num@margin_p4_min_signal_comp_b\endcsname{0.002610833}
\expandafter\def\csname p4num@margin_p4_min_comp_b_no_y\endcsname{0.005333333}
\expandafter\def\csname p4num@margin_p4_min_comp_b_class\endcsname{0.003333333}
\expandafter\def\csname p4num@margin_p4_min_valid_b_no_y\endcsname{0.005000000}
\expandafter\def\csname p4num@margin_p4_min_auction\endcsname{0.016666667}
\expandafter\def\csname p4num@margin_p4_min_actual_weak\endcsname{0.002574837}
\expandafter\def\csname p4num@margin_p4_min_actual_strong\endcsname{0.003899371}
\expandafter\def\csname p4num@margin_p4_min_s_no_y\endcsname{0.013181818}
\expandafter\def\csname p4num@margin_p4_min_s_class\endcsname{0.010454545}
\expandafter\def\csname p4num@margin_p4_min_b_no_y\endcsname{0.002070572}
\expandafter\def\csname p4num@margin_p4_min_grid\endcsname{0.000100000}

\newtheorem{proposition}{Proposition}
\newtheorem{lemma}{Lemma}
\newcommand{\Ent}{\mathsf E}
\newcommand{\OH}{\mathsf O_H}
\newcommand{\RT}{\mathcal R_T}
\newcommand{\ind}{\mathbf 1}
\newcommand{\D}{\Delta}
\newcommand{\supp}{\operatorname{supp}}
\newcommand{\calQ}{\mathcal Q}
\title{Information, preparation costs, and private signals}
\author{}
\date{23 October 2026}
\begin{document}
\maketitle

\section{The objects and the equilibrium conditions}

The order slot has the menu $\calQ=\{-1,0,1\}$.
The informed investor arrives with probability $\lambda$.
Otherwise, the liquidity probabilities are $(\beta,1-2\beta,\beta)$.
Write
\[
 b=(1-\lambda)\beta,\qquad b_0=(1-\lambda)(1-2\beta),\qquad u=\lambda/b.
\]
The quantity $b_0$ denotes the liquidity mass at zero; $C$ denotes the preparation cost.

Write $\Delta=\Delta_T$ and $\Delta_i=\Delta_T(F_i)$, as in the core proof.
The parameter $a_T$ denotes investor signal accuracy.
This notation reserves $a$ for the research cost parameter.

\subsection{Calibration B}

The primitive declaration is
\[
\begin{gathered}
(h,\ell,p)=(\reg{p4_input_h},\reg{p4_input_ell},\reg{p4_input_p}),\\
(\rho,c_H,k,\pi)=(\reg{p4_input_rho},\reg{p4_input_cH},\reg{p4_input_k},\reg{p4_input_pi}),\\
(r_0,r_1,r_2)=(\reg{p4_input_r0},\reg{p4_input_r1},\reg{p4_input_r2}),\\
c_L=0,\qquad 1+u=e,\qquad \beta=\reg{beta_B}\ \text{exactly},\qquad
\lambda=\frac{u\beta}{1+u\beta}.
\end{gathered}
\]
The derived presence probability is $\reg{lambda_B}$ to six decimals.
Every computation retains its unrounded definition.
The point-mass classes and the uniform incumbent define this benchmark.

\subsection{Auction objects}

Let the incumbent distribution be $F$ and the challenger class be $J\in\{H,L\}$.
Its prior is $\pi\in(0,1)$.
The class distributions have finite means and supports above $p$.
The distribution $G_H$ first-order dominates $G_L$.
All information here depends on the acquisition value only through its class.

A prepared challenger with value $\theta$ earns
\begin{equation}\label{eq:rent}
 g_F(\theta)=\mathbb E[(\theta-\max\{p,R\})_+]
 =\int_p^\theta F(v)\,dv.
\end{equation}
Indeed, the positive part equals $\int_p^\theta\ind\{R<v\}\,dv$ except on a null set of integration points.
Integration over the class gives
\begin{equation}\label{eq:classprofits}
 g_j(F)=\int_p^\infty[1-G_j(v)]F(v)\,dv,
 \qquad t_j(F)=\mathbb E[\theta\mid j]-g_j(F).
\end{equation}
With no challenger, $t_0(F)=p\Pr(R\ge p)$.
Define
\[
 w_j=t_j-t_0,\quad \D=t_H-t_L,\quad
 B_F(\mu)=g_L+\mu(g_H-g_L),\quad
 W_H(F)=\Pr(\theta\ge R\mid H).
\]
The bidder tie convention has no effect for the continuous incumbent laws used below.

A stronger incumbent has $F_1(v)\le F_0(v)$.
Equation~\eqref{eq:classprofits} gives $g_j(F_1)\le g_j(F_0)$ for both classes.
It also gives
\[
 \D(F_1)-\D(F_0)
 =\int_p^\infty[G_L(v)-G_H(v)][F_0(v)-F_1(v)]\,dv\ge0.
\]
Thus gross preparation profit falls at every fixed belief, while the target spread rises weakly.

For $R\sim U[0,r]$ and $p<\ell<r<h$, direct integration gives
\begin{align}\label{eq:auction}
 t_0&=p(1-p/r),&
 t_H&=r/2+p^2/(2r),&
 t_L&=\ell-(\ell^2-p^2)/(2r),\\
 g_H&=h-r/2-p^2/(2r),&
 g_L&=(\ell^2-p^2)/(2r),&
 \D&=(r-\ell)^2/(2r).\notag
\end{align}
For example, $t_H=p^2/r+\int_p^r v\,dv/r$.
For the low class, $t_L=p^2/r+\int_p^\ell v\,dv/r+\ell(r-\ell)/r$.
The rent identities then follow from $t_j+g_j=\theta_j$.
These formulas agree with every shared object in the team rules and equation (4) of the pinned source.

\subsection{The tuple without a private challenger signal}

An equilibrium is the tuple
\begin{equation}\label{eq:tuple}
 \mathcal E_F=(\sigma_H,\sigma_L,P,\mu,\zeta,\mathrm{bids}).
\end{equation}
The distributions $\sigma_j$ specify orders conditional on informed arrival.
Write $P_x=P(x)$ and $\mu_x=\mu(x)$; the derived challenger belief is $\bar\mu(P)$.
The cost-specific preparation rule is $\zeta(P,C)$.
For atomic costs, the core rule $\zeta(P)$ equals $\zeta(P,c_H)$.
Prepared bidders use $b_I(R)=R$ and $b_C(\theta)=\theta$.

The equilibrium conditions are:
\begin{enumerate}
\item Bayes consistency at every order.
\item Challenger beliefs: $\bar\mu(P)=\Pr(H\mid P)$ on the price image.
\item Optimal preparation given the price, the cost, and the available private information.
\item Competitive prices: $P_x=\mathbb E[V_T\mid X=x]$.
\item Optimal informed orders against the stated prices and preparation rule.
\item Truthful bids in the second-price auction with reserve $p$.
\end{enumerate}

A unilateral order keeps the candidate prices and preparation rule fixed.
The payoff conditions on informed arrival; it has no additional factor $\lambda$.
The tuple specifies preparation on the attainable price image.
Every admissible order reaches that image.

The investor rule IM permits mixtures between payoff maximizers.
The main statements use IM and challenger rule C1.
The alternative I0 selects zero when zero ties for the maximum.
The challenger rule C1 prepares at equality.
The alternative CM permits any preparation probability at equality.
The zero-cost challenger prepares at zero gross profit under both rules, as in the core definition.
Strict order and entry comparisons give the same result under these rules.
An atomless cost law makes the challenger tie rule irrelevant almost surely.

Put $\eta_{-1}=\eta_1=\beta$ and $\eta_0=1-2\beta$.
For arbitrary investor mixtures,
\begin{equation}\label{eq:bayes}
 f_j(x)=(1-\lambda)\eta_x+\lambda\sigma_j(x),\qquad
 \mu_x=\frac{\pi f_H(x)}{\pi f_H(x)+(1-\pi)f_L(x)}.
\end{equation}
Every order has positive mass in both classes.
For $0<\beta\le1/3$, its liquidity mass is at least $b$.
Consequently,
\begin{equation}\label{eq:bounds}
 m(\pi)=\frac{\pi}{\pi+(1-\pi)(1+u)}
 \le\mu_x\le
 \frac{\pi(1+u)}{\pi(1+u)+1-\pi}=M(\pi).
\end{equation}
To derive the bounds, divide the largest possible likelihood $b_x+\lambda$ by the smallest possible likelihood $b_x$.
The resulting ratio is at most $1+u$.
Bayes' odds formula gives both endpoints in~\eqref{eq:bounds}.

The price posterior, an average of order posteriors, obeys the same bounds.

Let $H_C$ be the independent cost CDF.
Set $e(\mu)=H_C(B_F(\mu))$ and assume $e\ge\rho>0$ on the feasible beliefs.
Competitive prices in any candidate satisfy
\begin{equation}\label{eq:inverse}
 P_x=t_0+e(P_x)(w_L+\D\mu_x),\qquad
 \mu_x=\frac{P_x-t_0-e(P_x)w_L}{e(P_x)\D}.
\end{equation}
Here $e(P)$ is the candidate's entry probability, before any sufficiency claim.
When $\D>0$, this identity makes $\mu_x$ a function of the price.
The conditional-expectation identity gives $\bar\mu(P_x)=\mu_x$.
This construction supplies the price and entry parts of~\eqref{eq:tuple}:
\begin{equation}\label{eq:construction}
 \zeta(P_x,C)=\ind\{C\le B_F(\mu_x)\},\quad
 P_x=t_0+e(\mu_x)(w_L+\D\mu_x),\quad \bar\mu(P_x)=\mu_x.
\end{equation}
For $\mu_2>\mu_1$, its price increment equals
\[
 e(\mu_1)\D(\mu_2-\mu_1)
 +[e(\mu_2)-e(\mu_1)](w_L+\D\mu_2)>0.
\]
This proves that the construction reveals the posterior when $w_L\ge0$ and $g_H>g_L$.

Against this fixed tuple, the four nonzero-order payoffs are
\begin{align}\label{eq:payoffs}
 U_H(1)&=e(\mu_1)\D(1-\mu_1)-k,&
 U_H(-1)&=-e(\mu_{-1})\D(1-\mu_{-1})-k,\\
 U_L(-1)&=e(\mu_{-1})\D\mu_{-1}-k,&
 U_L(1)&=-e(\mu_1)\D\mu_1-k.\notag
\end{align}
Both zero-order payoffs equal zero.
Wrong-signed orders lose strictly.
Correctly signed gross returns are at most $\D$ and at least $\rho\kappa(\pi)\D$, where
\[
 \kappa(\pi)=\min\{m(\pi),1-M(\pi)\}.
\]
These bounds hold for all pure or mixed candidate profiles and any price pools.

Truthful bids satisfy the final condition directly.
A challenger wins precisely when its bid exceeds the larger of the reserve and the incumbent bid.
A true-value bid wins all positive-surplus events and rejects all negative-surplus events.
The incumbent has the same argument with the two bidders interchanged.
Thus the bid component is optimal for every tuple constructed below.

\section{R5: The fixed-information control}

\begin{proposition}[Fixed information; analytical]\label{prop:control}
Fix the joint law of $(J,S)$, the conditional value laws $G_j$, and the independent cost law.
Let the value depend on $S$ only through $J$.
If the incumbent strengthens from $F_0$ to $F_1$, entry does not rise under C1.
Its decrease is exactly
\begin{equation}\label{eq:crossing}
 \Ent_0-\Ent_1
 =\Pr\{B_{F_1}(\mu(S))<C\le B_{F_0}(\mu(S))\},
 \qquad \mu(S)=\Pr(H\mid S).
\end{equation}
The decrease is strict exactly when this probability is positive.
\end{proposition}

\begin{proof}
Cost independence makes $\mu(S)$ the complete class belief at the preparation decision.
The gross profit is $B_{F_i}(\mu(S))$ in economy $i$.
Equation~\eqref{eq:classprofits} gives $B_{F_1}(\mu(S))\le B_{F_0}(\mu(S))$ pointwise.
Therefore
\[
 \ind\{C\le B_{F_0}(\mu(S))\}-\ind\{C\le B_{F_1}(\mu(S))\}
 =\ind\{B_{F_1}(\mu(S))<C\le B_{F_0}(\mu(S))\}.
\]
Its expectation proves both statements.
\end{proof}

The control fixes the information experiment, not the numerical price level.
In the slot version, it fixes $\lambda,\beta,\pi,G_H,G_L,H_C$ and both conditional order distributions.
The challenger then reoptimizes preparation, and market makers recompute the competitive prices.
Equation~\eqref{eq:inverse} ensures that these prices reveal the same order posterior at each strength.
The acquisition payoffs and price levels may change.

The proposition also holds with any fixed tie probability at equality, used in both economies.
Under CM alone, an unchanged gross profit at a cost atom can support different tie probabilities across economies.
An atomless law, or strict entry gaps at all realized posteriors, removes that issue.
Every displayed control has strict entry gaps.

\subsection{The frozen full-order tuple}

Impose $(\sigma_H,\sigma_L)=(\delta_1,\delta_{-1})$.
The complete conditional laws and posteriors are
\begin{equation}\label{eq:full}
\begin{array}{c|ccc}
 x&-1&0&1\\ \hline
 f_H(x)&b&b_0&b+\lambda\\
 f_L(x)&b+\lambda&b_0&b\\
 \mu_x&m(\pi)&\pi&M(\pi).
\end{array}
\end{equation}
For atomic costs, set $e_x=\rho+(1-\rho)\ind\{B_F(\mu_x)\ge c_H\}$.
Use~\eqref{eq:construction} for the prices, challenger beliefs, and entry rule.
Use the truthful bids from~\eqref{eq:tuple}.
Bayes' rule follows directly from each column of~\eqref{eq:full}.
The price and entry equations verify competitive prices and optimal preparation at all orders.

The four investor payoffs remain~\eqref{eq:payoffs}.
At the weak benchmark, the buy and sale payoffs are
\[
 \Delta_0(1-M)-k=\reg{p4_frozen_r0_buy_tp_payoff},\qquad
 \rho\Delta_0m-k=\reg{p4_frozen_r0_sell_tm_payoff}.
\]
Both are negative, so this weak tuple is an imposed control.
At the strong benchmark, both correct orders pay strictly more than zero.
The lower bound $\rho m\Delta_1-k=\reg{margin_p4_strong_trade_B}>0$ forces those orders in every candidate.
Thus the strong frozen tuple also equals the unique equilibrium.

At the benchmark prior, the atomic entry and ownership formulas are
\begin{equation}\label{eq:frozenlevels}
 \Ent=\rho+(1-\rho)(b+\lambda/2),\qquad
 \OH=\frac{\rho+(1-\rho)(b+\lambda)}2,
\end{equation}
provided $B_r(1/2)<c_H<B_r(M)$.
The target proceeds satisfy
\begin{equation}\label{eq:revenue}
 \RT=t_0+\frac{\bar e_Hw_H+\bar e_Lw_L}{2},\qquad
 \bar e_H=\rho+(1-\rho)(b+\lambda),\quad
 \bar e_L=\rho+(1-\rho)b.
\end{equation}
These follow from the class probabilities of the buy order.

For the declared auction inputs,
\[
 \tau(r)=\frac{48r-3}{4(20r-r^2-1)},\qquad
 \tau'(r)=\frac{48r^2-6r+12}{4(20r-r^2-1)^2}>0.
\]
The derivative numerator equals $48(r-1/16)^2+189/16$.
At $r>1$, $\tau(r)>\tau(1)=5/8>1/2$.
Let $r_C$ solve $B_{r_C}(M)=c_H$.
The unique solution is $r_C=\reg{rC_B}$.
The frozen entry path under C1 is
\begin{equation}\label{eq:frozenpath}
 \Ent(r)=
 \begin{cases}
 \rho+(1-\rho)(b+\lambda/2)=\reg{p4_frozen_entry_below_ceiling},&1<r\le r_C,\\
 \rho=\reg{p4_frozen_entry_above_ceiling},&r>r_C.
 \end{cases}
\end{equation}
For $1<r<h$, the strict threshold increase proves both branches.
For $r\ge h$, first-order dominance keeps $B_r(M)<c_H$, so the lower branch persists.
Rule C1 includes the top atom at $r_C$.
Rule CM permits any entry level between these two levels at $r_C$; it changes no other point.
At both $r_0$ and $r_1$, only the top atom recruits the high-cost type.
The strict-crossing probability is therefore $\reg{p4_control_crossing_r0_r1}$.
Between $r_1$ and $r_2$, it is $\reg{p4_control_crossing_r1_r2}$.

\subsection{The price-hidden equilibrium tuples}

The hidden-price control gives the challenger only its prior and its cost.
Market makers still observe the order and set competitive prices.
Under $B_{F_i}(\pi)<c_H$, the entry rule is $\zeta(C)=\ind\{C=0\}$ and $e=\rho$.
At $F_0$, take $\sigma_H=\sigma_L=\delta_0$.
All public beliefs are $\pi$, and every price equals $t_0+\rho(w_L+\D\pi)$.

At $F_1$, take the full orders and public beliefs in~\eqref{eq:full}.
The price is $P_x=t_0+\rho(w_L+\D\mu_x)$ at each order.
The challenger's belief stays $\pi$, because it does not receive the price or the order.
The entry rule, beliefs, and price formula directly verify the first four equilibrium conditions.
The payoffs in~\eqref{eq:payoffs} now use $e=\rho$ at all orders.

If $\Delta_0<k<\rho\kappa(\pi)\Delta_1$, the weak correct orders lose and the strong correct orders gain.
The wrong orders lose at both strengths.
These bounds hold for every profile and prove uniqueness of the two order tuples.
The bids satisfy the preceding auction argument.
Entry is $\rho$ at both strengths, while $\OH=\pi\rho W_H(F_i)$.

\begin{table}[htbp]\centering\small
\caption{Every information-control row at calibration B. Values are computer-assisted.}
\begin{tabular}{@{}lrrrrr@{}}\toprule
Control & $q_H$ & $q_L$ & $\Ent$ & $\OH$ & $\RT$\\\midrule
Frozen, $r_0$ & $1$ & $-1$ & \reg{p4_frozen_r0_entry}&\reg{p4_frozen_r0_own}&\reg{p4_frozen_r0_proceeds}\\
Frozen, $r_1$ & $1$ & $-1$ & \reg{p4_frozen_r1_entry}&\reg{p4_frozen_r1_own}&\reg{p4_frozen_r1_proceeds}\\
Price hidden, $r_0$ & $0$ & $0$ & \reg{p4_hidden_r0_entry}&\reg{p4_hidden_r0_own}&\reg{p4_hidden_r0_proceeds}\\
Price hidden, $r_1$ & $1$ & $-1$ & \reg{p4_hidden_r1_entry}&\reg{p4_hidden_r1_own}&\reg{p4_hidden_r1_proceeds}\\\bottomrule
\end{tabular}
\end{table}

\section{R7: Atomless preparation costs}

Atomless preparation costs are the leading application of the next result.
The support conditions also cover cost laws with atoms.

\begin{proposition}[The cost-support reversal; analytical]\label{prop:atomless}
Let $F_1$ strengthen $F_0$.
Use the class and auction assumptions of Section 1.
Use $0<\pi,\rho,\lambda<1$, $0<\beta\le1/3$, and an independent cost law
\[
 H_C=\rho H^0+(1-\rho)H^1,\quad
 \supp H^0\subset[0,L],\quad \supp H^1\subset[c_-,c_+].
\]
Both component laws are arbitrary probability laws on their stated supports.
Suppose
\begin{align}\label{eq:costconditions}
 &0<L<B_{F_1}(m(\pi)),\qquad
 B_{F_0}(\pi)<c_-<c_+<B_{F_1}(M(\pi)),\\
 &\Delta_0<k<\rho\kappa(\pi)\Delta_1.\notag
\end{align}
Then the unique orders are zero at $F_0$ and full at $F_1$.
These conclusions hold under IM or I0 and C1 or CM.
Entry and class entry at $F_1$ are exactly
\begin{equation}\label{eq:costentry}
 \Ent_1=\rho+(1-\rho)(b+\pi\lambda),\quad
 \bar e_H=\rho+(1-\rho)(b+\lambda),\quad
 \bar e_L=\rho+(1-\rho)b.
\end{equation}
At $F_0$, entry is $\rho$.
High-class ownership rises exactly when
\begin{equation}\label{eq:owncondition}
 [\rho+(1-\rho)(b+\lambda)]W_H(F_1)>\rho W_H(F_0).
\end{equation}
In particular, it rises if $W_H(F_0)=W_H(F_1)=1$.
\end{proposition}

\begin{proof}
\emph{Consequences of the assumptions.}
Class dominance gives $\Delta_i\ge0$, and auction payoffs give $w_L(F_i)\ge0$.
Since $k>0$, the strong-order inequality implies $\Delta_1>0$.
The support gaps give
\[
 B_{F_1}(M)>B_{F_0}(\pi)\ge B_{F_1}(\pi),\qquad
 B_{F_0}(M)\ge B_{F_1}(M)>B_{F_0}(\pi).
\]
Since $M>\pi$, both class profit gaps are positive.
No assumption requires $\Delta_0>0$.

\emph{Construction.}
At $F_0$, set $(\sigma_H,\sigma_L)=(\delta_0,\delta_0)$, $\mu_x=\pi$, and $e_x=\rho$ for every order.
Set $P_x=t_0+\rho(w_L+\Delta_0\pi)$ and $\zeta(P_x,C)=\ind\{C\le B_{F_0}(\pi)\}$.
At $F_1$, use~\eqref{eq:full}, set $e_1=1$ and $e_0=e_{-1}=\rho$, and use~\eqref{eq:construction}.
Both tuples use the truthful bids.

\emph{Bayes consistency and prices.}
The zero-order tuple has identical class likelihoods at every order.
The full-order tuple has the three likelihood pairs in~\eqref{eq:full}.
Thus its posteriors are $(m,\pi,M)$.
The stated prices equal the conditional means of $t_0+e_x(t_j-t_0)$.
At the weak tuple, every order has the same posterior and price.
At the strong tuple, the price increment after~\eqref{eq:construction} is positive.
Thus each price supplies its stated posterior, including when $\Delta_0=0$.

\emph{Preparation.}
At every feasible posterior in either economy,
\[
 B_{F_i}(\mu)\ge B_{F_1}(m)>L.
\]
Hence every low-component cost prepares.
At the weak prior, $B_{F_0}(\pi)<c_-$ excludes the entire high component.
At the strong top posterior, $B_{F_1}(M)>c_+$ includes the entire high component.
At the other strong posteriors,
\[
 B_{F_1}(m)<B_{F_1}(\pi)\le B_{F_0}(\pi)<c_-.
\]
This verifies each entry value against its stated tuple.

\emph{Investor incentives.}
At the weak tuple, the two correct gross returns are $\rho\Delta_0(1-\pi)$ and $\rho\Delta_0\pi$.
Both are below $k$ by~\eqref{eq:costconditions}.
At the strong tuple, a high-class buy nets $\Delta_1(1-M)-k>0$.
A low-class sale nets $\rho\Delta_1m-k>0$.
The remaining two nonzero actions have the negative payoffs in~\eqref{eq:payoffs}.
Zero pays zero, so the stated orders maximize payoffs.
The preceding bid argument verifies truthful bids.

\emph{Uniqueness.}
For any candidate profile, the price posterior obeys~\eqref{eq:bounds}, so the support restriction gives $e(P)\ge\rho$.
Competitive prices give the payoffs in~\eqref{eq:payoffs} with $e(P_x)$, before price sufficiency.
Every correct weak order earns at most $\Delta_0-k<0$.
If $\Delta_0=0$, every nonzero order pays $-k$.
Thus the weak orders are zero, and their common prior fixes entry and prices.
At the strong strength, $\Delta_1>0$ and~\eqref{eq:inverse} gives price sufficiency under every profile.
Every correct strong order earns at least $\rho\kappa\Delta_1-k>0$.
Wrong orders lose at both strengths.
Thus every candidate must use the constructed orders, which determine the prices and entry probabilities.
Every realized cost gap is strict, so C1 and CM give the same continuation.

\emph{Outcomes.}
A buy has probabilities $b+\lambda$ in class $H$ and $b$ in class $L$.
Only this order admits high-cost entry.
The formulas in~\eqref{eq:costentry} follow by total probability.
Conditional on class, the value remains independent of preparation and the order.
Thus ownership equals $\pi\bar e_HW_H(F_1)$ and $\pi\rho W_H(F_0)$ in the two economies.
Their strict comparison is~\eqref{eq:owncondition}.
\end{proof}

Neither atomlessness nor a density assumption enters this proof.
If both component laws are atomless, the entire cost law is atomless.
The atomic benchmark also satisfies the result with $H^0=\delta_0$ and $H^1=\delta_{c_H}$.
The common support bound $L$ remains a strictly positive upper envelope for that mass.

\subsection{Calibration B and the third part of Proposition 2}

Use the declared continuous law
\begin{equation}\label{eq:costlaw}
 H_C=\frac14 U_{[0,1/5]}+\frac34 U_{[59/10,61/10]}.
\end{equation}
Here $U_{[v,w]}$ denotes a uniform CDF.
Every strict inequality in~\eqref{eq:costconditions} holds.
Strong entry is $\reg{p4_atomless_r1_entry}$ and ownership is $\reg{p4_atomless_r1_own}$.
Both values equal their atomic-cost counterparts exactly.
The strict support gaps put every realized posterior outside both cost bands at the first two strengths.

The third part has an additional support condition.
For a stronger $F_2$, sufficient strict conditions are
\begin{equation}\label{eq:returnconditions}
 L<B_{F_2}(m),\qquad B_{F_2}(M)<c_-,\qquad
 k<\rho\kappa\Delta_2.
\end{equation}
The full-order tuple then has $e_x=\rho$ at every order.
Equation~\eqref{eq:payoffs} verifies the orders, and the same bound proves their uniqueness under every profile.
The support inequalities exclude the entire high component and include the entire low component.
Therefore entry returns to $\rho$.

For~\eqref{eq:costlaw}, the second margin in~\eqref{eq:returnconditions} is
\[
 c_--B_{r_2}(M)=\reg{margin_p4_atomless_return_B}<0.
\]
The top profit lies strictly inside the high-cost band.
The full-order equilibrium still exists uniquely, but its entry is
\begin{equation}\label{eq:costramp}
 \Ent(r)=\rho+(1-\rho)(b+\pi\lambda)
 H^1(B_r(M)),
\end{equation}
whenever the low component enters everywhere and $B_r(\pi)<c_-$.
The same expression replaces $b+\pi\lambda$ by $b+\lambda$ for high-class entry.
At $r_2$, it gives $\Ent=\reg{p4_atomless_r2_entry}$ and $\OH=\reg{p4_atomless_r2_own}$.
This follows by direct evaluation of the uniform CDF, not by a continuity approximation.

For a symmetric half-width $\varepsilon$, the return condition requires
\[
 0<\varepsilon<c_H-B_{r_2}(M)
 =\reg{p4_atomless_max_width_floor_r2}.
\]
The choice $\varepsilon=\reg{p4_atomless_small_width}$ satisfies all three parts.
Its return margin is $\reg{margin_p4_atomless_small_width_return}>0$.
For any atomless $H^1$,~\eqref{eq:costramp} is continuous and nonincreasing in strength.
The uniform law makes it affine in $B_r(M)$ within its transition band.

\section{R8: Complementary private signals}

\subsection{Information, prices, and the tuple}

The prior is $\pi=1/2$ and values are the point masses $h$ and $\ell$.
The investor sees $T\in\{+,-\}$ and the challenger sees $Y\in\{+,-\}$.
Conditional on the class, the signals are independent and have accuracies
\[
 \begin{aligned}
 \Pr(T=+\mid H)&=\Pr(T=-\mid L)=a_T,\\
 \Pr(Y=+\mid H)&=\Pr(Y=-\mid L)=d,
 \qquad 1/2<a_T,d<1.
 \end{aligned}
\]
The identity draw, the liquidity order, the cost, and $R$ are independent of these signals conditional on class.
Preparation reveals the acquisition value.

Replace $\sigma_H,\sigma_L$ in~\eqref{eq:tuple} by $\sigma_+,\sigma_-$.
Add the beliefs $\nu_x=\Pr(T=+\mid X=x)$ and $\mu_{P,y}=\Pr(H\mid P,Y=y)$.
Replace $\zeta(P,C)$ by $\zeta(P,Y,C)$.
The six equilibrium conditions remain the same, with the investor choice conditional on $T$.

Write $A_t(x)=(1-\lambda)\eta_x+\lambda\sigma_t(x)$.
Conditional independence and equal signal priors give
\begin{align}\label{eq:signalsbelief}
 \nu_x&=\frac{A_+(x)}{A_+(x)+A_-(x)},&
 \mu_x&=(1-a_T)+(2a_T-1)\nu_x,\\
 \phi_+(\mu)&=\frac{d\mu}{d\mu+(1-d)(1-\mu)},&
 \phi_-(\mu)&=\frac{(1-d)\mu}{(1-d)\mu+d(1-\mu)}.\notag
\end{align}
The identity $J\perp X\mid T$ gives the affine formula.
The identity $Y\perp X\mid J$ gives $\mu_{P,y}=\phi_y(\bar\mu(P))$.
Both identities remain true under arbitrary signal-contingent mixtures.
The likelihood bound gives $m\le\nu_x\le M$, where $m=m(1/2)$ and $M=M(1/2)$.
Thus
\[
 \mu_-=(1-a_T)+(2a_T-1)m\le\mu_x\le
 (1-a_T)+(2a_T-1)M=\mu_+.
\]

The zero-cost challenger always prepares because $g_L>0$.
At a public posterior $\mu$, put
\begin{align}\label{eq:signalentry}
 I_y(\mu)&=\ind\{B_r(\phi_y(\mu))\ge c_H\},\\
 e_H(\mu)&=\rho+(1-\rho)[dI_+(\mu)+(1-d)I_-(\mu)],\notag\\
 e_L(\mu)&=\rho+(1-\rho)[(1-d)I_+(\mu)+dI_-(\mu)],\notag\\
 \Delta_P(\mu)&=e_H(\mu)w_H-e_L(\mu)w_L.\notag
\end{align}
Since $\phi_+>\phi_-$, optimality gives $I_+\ge I_-$, also under CM.
Consequently,
\begin{equation}\label{eq:signalbounds}
 \rho\D\le\Delta_P
 =e_L\D+(e_H-e_L)w_H
 \le\D+(1-\rho)(2d-1)w_H=:D^+.
\end{equation}
Both price coefficients are class entry rates at that price.
They are not the unconditional rates across prices.

In any candidate, competitive prices satisfy
\begin{equation}\label{eq:signalinverse}
 P_x=t_0+e_L(P_x)w_L+\mu_x\Delta_P(P_x),\qquad
 \mu_x=\frac{P_x-t_0-e_L(P_x)w_L}{\Delta_P(P_x)}.
\end{equation}
The positive coefficient gives $\bar\mu(P)=\mu_x$ before any order restriction.

Construct the entry rates from~\eqref{eq:signalentry} and the prices from~\eqref{eq:signalinverse}.
The cost rule is
\[
 \zeta(P_x,y,C)=\ind\{C\le B_r(\phi_y(\mu_x))\}.
\]
For $\mu_2>\mu_1$, the price change is
\begin{align*}
 P(\mu_2)-P(\mu_1)={}&(\mu_2-\mu_1)\Delta_P(\mu_1)\\
 &+\mu_2[e_H(\mu_2)-e_H(\mu_1)]w_H\\
 &+(1-\mu_2)[e_L(\mu_2)-e_L(\mu_1)]w_L>0.
\end{align*}
Each bracket is nonnegative, and the first term is positive.
This verifies that the constructed prices supply the stated beliefs.

A unilateral order reveals no further class information to the investor.
Its class probabilities remain $a_T$ after $T=+$ and $1-a_T$ after $T=-$.
Subtract the fixed price in~\eqref{eq:signalinverse} from the corresponding terminal value.
The four nonzero-order payoffs are
\begin{align}\label{eq:signalpayoffs}
 U_+(1)&=(2a_T-1)(1-\nu_1)\Delta_P(\mu_1)-k,\\
 U_+(-1)&=-(2a_T-1)(1-\nu_{-1})\Delta_P(\mu_{-1})-k,\notag\\
 U_-(-1)&=(2a_T-1)\nu_{-1}\Delta_P(\mu_{-1})-k,\notag\\
 U_-(1)&=-(2a_T-1)\nu_1\Delta_P(\mu_1)-k.\notag
\end{align}
Zero gives zero after either signal.
Wrong-signed orders lose strictly.
Every correct gross return lies between $(2a_T-1)m\rho\D$ and $(2a_T-1)D^+$.

\subsection{The general result}

\begin{proposition}[Complementary private signals; analytical]\label{prop:signals}
Choose $0<p<\ell<r_0<r_1<h$, $0<\rho,\lambda<1$, and $0<\beta\le1/3$.
Let $1/2<a_T,d<1$ and suppose
\begin{align}\label{eq:signalconditions}
 B_{r_0}(d)&<c_H<B_{r_1}(\phi_+(\mu_+)),\\
 (2a_T-1)D^+_0&<k<(2a_T-1)m\rho\Delta_1.\notag
\end{align}
The unique weak orders are zero, with entry $\rho$.
The unique strong orders are $q(+)=1$ and $q(-)=-1$, with entry strictly above $\rho$.
These order conclusions hold under IM or I0 and C1 or CM.
Under C1, the tuples below determine every continuation choice.
If the strong realized entry gaps are strict, the tuples are also unique under CM.
The hypotheses permit $d>a_T$.
\end{proposition}

\begin{proof}
\emph{Weak tuple.}
Take $\sigma_+=\sigma_-=\delta_0$.
At every order set $\nu_x=\mu_x=1/2$, $\mu_{P,+}=d$, and $\mu_{P,-}=1-d$.
Set $\zeta(P_x,y,C)=\ind\{C=0\}$ and
\[
 P_x=t_0+\rho(w_H+w_L)/2.
\]
These assignments, with truthful bids, specify the whole tuple.
The identical conditional order laws verify Bayes' rule.
The cost inequality excludes high-cost entry after either private signal.
Low-cost entry is strict because $g_L>0$.
The price is the conditional target mean.

At this fixed tuple, each correct order earns $(2a_T-1)\rho\Delta_0/2-k$.
It is negative because $\rho\Delta_0/2\le D^+_0$ and~\eqref{eq:signalconditions} holds.
The two wrong orders earn its negative gross component minus $k$.
Thus zero is optimal after both signals.

\emph{Strong tuple.}
Take $\sigma_+=\delta_1$ and $\sigma_-=\delta_{-1}$.
Specify the order laws by
\begin{equation}\label{eq:signalfull}
\begin{array}{c|ccc}
 x&-1&0&1\\ \hline
 A_+(x)&b&b_0&b+\lambda\\
 A_-(x)&b+\lambda&b_0&b\\
 \nu_x&m&1/2&M\\
 \mu_x&\mu_-&1/2&\mu_+.
\end{array}
\end{equation}
Use~\eqref{eq:signalsbelief} for each private posterior.
Use~\eqref{eq:signalentry} for each class entry rate and the stated cost rule for each private realization.
Use~\eqref{eq:signalinverse} for all three prices, and use truthful bids.
This explicit tuple covers every order and both private signals.

Every column of~\eqref{eq:signalfull} satisfies Bayes' rule.
The price formula averages the class target values with those beliefs.
The cost rule maximizes gross profit minus the realized cost.
Against these fixed prices and entry rules,
\[
 U_+(1)\ge(2a_T-1)m\rho\Delta_1-k>0,\qquad
 U_-(-1)\ge(2a_T-1)m\rho\Delta_1-k>0.
\]
The wrong-order payoffs in~\eqref{eq:signalpayoffs} are negative and zero pays zero.
Thus the two full orders are optimal.
The bid argument in Section 1 verifies the auction choices.

\emph{Uniqueness and additional entry.}
For every candidate,~\eqref{eq:signalbounds} and the likelihood bounds hold.
They force zero orders in the weak economy and full orders in the strong economy.
This proves uniqueness without a restriction to pure profiles.
At the strong buy price, $B_{r_1}(\phi_+(\mu_+))>c_H$.
The event $(X,Y)=(1,+)$ has positive probability in each class.
It adds strictly positive high-cost preparation and proves $\Ent_1>\rho$.
\end{proof}

The class laws of the strong order are
\[
 \begin{aligned}
 f_H(1)&=b+a_T\lambda,& f_L(1)&=b+(1-a_T)\lambda,\\
 f_H(-1)&=b+(1-a_T)\lambda,& f_L(-1)&=b+a_T\lambda,\\
 f_H(0)&=b_0,&f_L(0)&=b_0.
 \end{aligned}
\]
The general outcome formulas use all six entry indicators:
\begin{equation}\label{eq:signaloutcomes}
 \bar e_j=\rho+(1-\rho)\sum_xf_j(x)
 \big[\Pr(Y=+\mid j)I_+(\mu_x)+\Pr(Y=-\mid j)I_-(\mu_x)\big].
\end{equation}
Then $\Ent=(\bar e_H+\bar e_L)/2$, $\OH=\bar e_H/2$, and~\eqref{eq:revenue} gives proceeds.
The ownership identity uses $h>r_1$.

The extra strict inequality
\begin{equation}\label{eq:onlypair}
 B_{r_1}(\phi_-(\mu_+))<c_H
\end{equation}
leaves exactly one high-cost entry pair: $(X,Y)=(1,+)$.
Indeed, $B_{r_1}(d)\le B_{r_0}(d)<c_H$ excludes $Y=+$ at zero and negative orders.
Monotonicity of $\phi_-$ and~\eqref{eq:onlypair} exclude $Y=-$ everywhere.
Thus
\begin{equation}\label{eq:onepair}
 \bar e_H=\rho+(1-\rho)d(b+a_T\lambda),\qquad
 \bar e_L=\rho+(1-\rho)(1-d)(b+(1-a_T)\lambda).
\end{equation}
Every entry choice is strict on this parameter set.

To require both sources of information, add
\begin{equation}\label{eq:signalcomplement}
 B_{r_1}(\mu_+)<c_H.
\end{equation}
Without $Y$, the weak upper return bound still forces zero orders.
Then $B_{r_0}(1/2)<B_{r_0}(d)<c_H$ gives entry $\rho$.
At the strong strength, the lower return bound still forces full orders.
Condition~\eqref{eq:signalcomplement} excludes high-cost preparation under every profile, so strong entry is also $\rho$.
With $Y$, the joint-entry inequality instead gives entry above $\rho$.
The strict inequalities define an open region for this complementarity result.

\subsection{The signal example with B slot parameters}

Denote the source's signal example by S(B).
It retains $h,\ell,p,\pi,\beta,\lambda$ from B but declares
\[
 \begin{gathered}
 \rho=\reg{p4_input_s_rho},\quad c_H=\reg{p4_input_s_cost},\quad
 k=\reg{p4_input_s_fee},\\
 r_0=\reg{p4_input_s_r0},\quad r_1=\reg{p4_input_s_r1},\\
 (a_T,d)=(\reg{p4_input_signal_a},\reg{p4_input_signal_d}).
 \end{gathered}
\]
Equations~\eqref{eq:signalconditions} and~\eqref{eq:onlypair} hold strictly.
Entry rises from $\reg{p4_s_signal_r0_entry}$ to $\reg{p4_s_signal_r1_entry}$.
Ownership rises from $\reg{p4_s_signal_r0_own}$ to $\reg{p4_s_signal_r1_own}$.
These are the earlier headline displays, recomputed with the exact decimal $\beta$ in the present brief.
They describe a strength comparison in S(B), not a comparison at the unchanged B primitives.

For clarity, the table also removes the challenger's private signal while it preserves the investor signal.
That control has $d=1/2$ in the posterior formulas.
Its maximum public belief gives $B_{r_i}(\mu_+)<c_H$ at both strengths, so entry is always $\rho$.
At $r_0$, the upper return bound excludes every nonzero order.
At $r_1$, the lower return bound forces both full orders.
The tuple formulas above then verify the complete control equilibria.

\begin{table}[htbp]\centering\small
\caption{Information comparisons in S(B); computer-assisted.}
\begin{tabular}{@{}lrrrr@{}}\toprule
Information & $\Ent_0$ & $\Ent_1$ & $\mathsf O_{H,0}$ & $\mathsf O_{H,1}$\\\midrule
Investor sees $J$, no $Y$ & \reg{p4_s_class_r0_entry}&\reg{p4_s_class_r1_entry}&\reg{p4_s_class_r0_own}&\reg{p4_s_class_r1_own}\\
Investor sees $T$, no $Y$ & \reg{p4_s_no_y_r0_entry}&\reg{p4_s_no_y_r1_entry}&\reg{p4_s_no_y_r0_own}&\reg{p4_s_no_y_r1_own}\\
Investor sees $T$, challenger sees $Y$ & \reg{p4_s_signal_r0_entry}&\reg{p4_s_signal_r1_entry}&\reg{p4_s_signal_r0_own}&\reg{p4_s_signal_r1_own}\\\bottomrule
\end{tabular}
\end{table}

For the class-informed row, the weak spread is below $k$ and the strong lower return exceeds $k$.
Its strong top belief also leaves gross profit below $c_H$.
The zero and full tuples of Section 1 therefore prove that row.

The original accuracy grid uses
\[
\begin{aligned}
a_T&\in\{\reg{p4_accuracy_68},\reg{p4_accuracy_69},\reg{p4_accuracy_70},\reg{p4_accuracy_71},\reg{p4_accuracy_72}\},\\
d&\in\{\reg{p4_accuracy_73},\reg{p4_accuracy_74},\reg{p4_accuracy_75},\reg{p4_accuracy_76},\reg{p4_accuracy_77}\}.
\end{aligned}
\]
Exactly the six pairs in Table~\ref{tab:pairs} satisfy the four sufficient inequalities.
Each also satisfies~\eqref{eq:onlypair}, and each has $d>a_T$.
The table states sufficient pairs on this grid; equations~\eqref{eq:signalconditions} define the full open sufficient region.

\begin{table}[htbp]\centering\small
\caption{The sufficient signal pairs in S(B); computer-assisted.}\label{tab:pairs}
\begin{tabular}{@{}rrrr@{}}\toprule
$a_T$ & $d$ & $\Ent_1$ & $\mathsf O_{H,1}$\\\midrule
\reg{p4_accuracy_69}&\reg{p4_accuracy_75}&\reg{p4_grid_69_75_entry}&\reg{p4_grid_69_75_own}\\
\reg{p4_accuracy_70}&\reg{p4_accuracy_75}&\reg{p4_grid_70_75_entry}&\reg{p4_grid_70_75_own}\\
\reg{p4_accuracy_71}&\reg{p4_accuracy_74}&\reg{p4_grid_71_74_entry}&\reg{p4_grid_71_74_own}\\
\reg{p4_accuracy_71}&\reg{p4_accuracy_75}&\reg{p4_grid_71_75_entry}&\reg{p4_grid_71_75_own}\\
\reg{p4_accuracy_72}&\reg{p4_accuracy_74}&\reg{p4_grid_72_74_entry}&\reg{p4_grid_72_74_own}\\
\reg{p4_accuracy_72}&\reg{p4_accuracy_75}&\reg{p4_grid_72_75_entry}&\reg{p4_grid_72_75_own}\\\bottomrule
\end{tabular}
\end{table}

\subsection{The unchanged B primitives}

The earlier signal pair has a different result at unchanged B.
Its private-only margin is $\reg{margin_p4_signal_b_private}$.
The weak-order margin is $\reg{margin_p4_signal_b_weak}$, and the strong-order margin is $\reg{margin_p4_signal_b_strong}$.
Thus the earlier sufficient theorem does not apply to that pair at B.
The next calculation derives its actual equilibria under IM.

\paragraph{Exact B continuation (analytical).}
Equations~\eqref{eq:signalbounds} and~\eqref{eq:signalpayoffs} exclude wrong orders under every candidate profile.
Define the effective order rates
\[
 s=\lambda\Pr(q=1\mid T=+),\qquad t=\lambda\Pr(q=-1\mid T=-).
\]
The associated order distributions are
\[
 \sigma_+=\frac{s}{\lambda}\delta_1+
 \left(1-\frac{s}{\lambda}\right)\delta_0,\qquad
 \sigma_-=\frac{t}{\lambda}\delta_{-1}+
 \left(1-\frac{t}{\lambda}\right)\delta_0.
\]
Their signal likelihoods are
\begin{equation}\label{eq:mixedlaws}
\begin{array}{c|ccc}
 x&-1&0&1\\ \hline
 A_+(x)&b&b_0+\lambda-s&b+s\\
 A_-(x)&b+t&b_0+\lambda-t&b.
\end{array}
\end{equation}
Equations~\eqref{eq:signalsbelief},~\eqref{eq:signalentry}, and~\eqref{eq:signalinverse} specify every remaining component of the tuple.
The bids remain truthful.

The buy column depends only on $s$, and the sale column depends only on $t$.
Specifically,
\[
 \nu_1(s)=\frac{b+s}{2b+s},\qquad \nu_{-1}(t)=\frac{b}{2b+t},\qquad
 \mu_x=(1-a_T)+(2a_T-1)\nu_x.
\]
Under C1, each posterior fixes its entry rates and $\Delta_P$.
Thus the two correct gross returns are separate functions:
\[
 \begin{aligned}
 K_+(s)&=(2a_T-1)\frac{b}{2b+s}\Delta_P(\mu_1(s)),\\
 K_-(t)&=(2a_T-1)\frac{b}{2b+t}\Delta_P(\mu_{-1}(t)).
 \end{aligned}
\]
At a cost threshold, CM replaces the relevant coefficient by its admissible interval.
That interval still depends only on the corresponding order rate.
Zero pays zero after either signal, although its posterior can depend on both rates.
Therefore the buy and sale problems separate, including their endpoint and threshold comparisons.

\paragraph{Weak strength.}
At $r_0$, the worst favorable private posterior satisfies $B_{r_0}(\phi_+(\mu_-))>6$.
The best unfavorable private posterior satisfies $B_{r_0}(\phi_-(\mu_+))<6$.
These inequalities hold under every profile, so $I_+=1$ and $I_-=0$ at every price.
Therefore
\[
 e_H=13/16,\qquad e_L=7/16,\qquad \Delta_P=311/1920.
\]

The two correct gross returns are
\[
 K_+(s)=\frac{2}{5}\frac{b}{2b+s}\frac{311}{1920},\qquad
 K_-(t)=\frac{2}{5}\frac{b}{2b+t}\frac{311}{1920}.
\]
Each function decreases strictly in its argument.
At zero it exceeds $k$; at $\lambda$ it lies below $k$.
The unique zero-profit rates are
\begin{equation}\label{eq:weakmixed}
 s=t=\frac{119}{96}b,\qquad
 \Pr(q=1\mid +)=\Pr(q=-1\mid -)=\frac{119}{96u}.
\end{equation}

At this tuple, each used order and zero yield zero.
Each wrong order loses strictly by~\eqref{eq:signalpayoffs}.
Bayes' rule, prices, entry, and truthful bids follow from the tuple formulas and the strict entry gaps.
The endpoint signs and strict decrease prove uniqueness among all mixtures.
Since entry is class-dependent but price-independent, $\Ent_0=5/8$ and $\mathsf O_{H,0}=13/32$.

\paragraph{Strong strength.}
At $r_1$, $Y=-$ never recruits under any profile.
A nonnegative signal posterior after a buy has $\mu_1\ge1/2$ once wrong orders are excluded.
Since $B_{r_1}(d)>6$, every buy price has $I_+=1$.
There $\Delta_P=137/192$.
The minimum buy payoff is $\reg{margin_p4_b_strong_buy}>0$, so $s=\lambda$ is necessary.

At a sale price, the favorable signal recruits exactly when
\[
 \mu_{-1}\ge\frac{47}{106}
 \quad\Longleftrightarrow\quad
 \nu_{-1}\ge\frac{19}{53}
 \quad\Longleftrightarrow\quad t\le\frac{15}{19}b.
\]
On that interval, the gross sale return exceeds the floor return at the endpoint.
Even the floor return there exceeds $k$, because
\[
 \frac{19}{795}-\frac1{50}=\reg{margin_p4_b_strong_threshold_floor_return}>0.
\]
This also excludes a CM threshold mixture.

Beyond the threshold, $e_H=e_L=1/4$, so the sale return is
\[
 K_-(t)=\frac1{15}\frac{b}{2b+t}.
\]
It decreases strictly, exceeds $k$ just above the threshold, and falls below $k$ at $\lambda$.
Its unique zero-profit rate is
\begin{equation}\label{eq:strongmixed}
 s=\lambda,\qquad t=\frac43 b,\qquad
 \Pr(q=-1\mid -)=\frac4{3u}.
\end{equation}
At this rate, $\nu_{-1}=3/10$ and $\mu_{-1}=21/50<47/106$.
At zero flow, direct substitution in~\eqref{eq:mixedlaws} gives $B_{r_1}(\phi_+(\mu_0))>6$.
Thus $Y=+$ recruits at $X=0,1$ and does not recruit at $X=-1$.
All six entry comparisons are strict.

The tuple defined by~\eqref{eq:strongmixed} and~\eqref{eq:mixedlaws} has positive buy profit and zero sale profit.
The used sale and zero are optimal after $T=-$; wrong orders lose.
Its entry, beliefs, prices, and bids satisfy the other five equilibrium conditions directly.
The buy bound and the sale argument prove uniqueness under IM and either challenger tie rule.
Both mixed equilibria depend on IM; I0 admits no equilibrium at these two parameter points.
Their unique necessary positive interior rates are unavailable under I0.

Only a sale blocks high-cost entry after $Y=+$ at $r_1$.
Its class probabilities are $b+(1-a_T)t$ and $b+a_Tt$.
Hence
\begin{equation}\label{eq:bactual}
 \Ent_1=\frac58-\frac{23}{40}b=\reg{p4_b_signal_r1_entry},\qquad
 \mathsf O_{H,1}=\frac{13}{32}-\frac{63}{160}b=\reg{p4_b_signal_r1_own}.
\end{equation}
Both fall relative to the weak economy.

If the challenger has no private signal, its highest public profit remains below $6$ at both strengths.
Thus $e_H=e_L=\rho$ under every profile.
At $r_0$, $\rho(2a_T-1)\Delta_0/2<k$ forces zero orders.
At $r_1$, the unique zero-profit rates are $s=t=4b/3$.
The same strict-decrease argument verifies that tuple and its uniqueness under IM.
Its entry and ownership stay at $\rho$ and $\rho/2$.
The entry gaps are strict, so C1 and CM give the same result.
At $r_1=3$, this control requires IM; I0 admits no equilibrium.
At zero orders, the gross return is $\reg{p4_b_no_y_r1_zero_gross}>k$.
At full orders, it is $\reg{p4_b_no_y_r1_full_gross}<k$.
Each side therefore needs its unique positive interior rate, which I0 excludes.

\begin{table}[htbp]\centering\small
\caption{Information comparisons at unchanged B, earlier pair $(0.70,0.75)$; computer-assisted.}
\begin{tabular}{@{}lrrrr@{}}\toprule
Information & $\Ent_0$ & $\Ent_1$ & $\mathsf O_{H,0}$ & $\mathsf O_{H,1}$\\\midrule
Investor sees $J$, no $Y$ & \reg{p4_base_r0_entry}&\reg{p4_base_r1_entry}&\reg{p4_base_r0_own}&\reg{p4_base_r1_own}\\
Investor sees $T$, no $Y$ & \reg{p4_b_no_y_r0_entry}&\reg{p4_b_no_y_r1_entry}&\reg{p4_b_no_y_r0_own}&\reg{p4_b_no_y_r1_own}\\
Investor sees $T$, challenger sees $Y$ & \reg{p4_b_signal_r0_entry}&\reg{p4_b_signal_r1_entry}&\reg{p4_b_signal_r0_own}&\reg{p4_b_signal_r1_own}\\\bottomrule
\end{tabular}
\end{table}

The manuscript uses this complementary pair at unchanged B:
\[
 (a_T,d)=(\reg{p4_input_comp_a},\reg{p4_input_comp_d}).
\]
Its four theorem margins, the one-pair margin, and the no-$Y$ exclusion margin are strictly positive.
It has zero weak orders and full strong orders.
Equations~\eqref{eq:onepair} give entry $\reg{p4_comp_b_r0_entry}\to\reg{p4_comp_b_r1_entry}$ and
ownership $\reg{p4_comp_b_r0_own}\to\reg{p4_comp_b_r1_own}$.
Its smallest theorem margin is $\reg{margin_p4_min_signal_comp_b}$.
Condition~\eqref{eq:signalcomplement} has margin $\reg{margin_p4_comp_b_no_y_top_r1}$.
Without $Y$, entry stays at $\rho$ at both strengths.

For this control, use $d=1/2$ and the zero and full tuples above.
The weak bound $k-(2a_T-1)\Delta_0>0$ excludes every nonzero order under every profile.
The weak prior gap then excludes high-cost preparation.
At the strong strength, $(2a_T-1)\rho m\Delta_1-k>0$ forces full orders under every profile.
The top public profit remains below $c_H$, so every class entry rate equals $\rho$.
The tuple formulas verify Bayes' rule, prices, preparation, and bids.

The class-observed control also removes $Y$ and sets $a_T=1$.
It equals the benchmark: zero weak orders, full strong orders, and high-cost entry only at a strong buy price.
The bounds $\Delta_0<k<\rho m\Delta_1$ force these orders under every profile.
The weak prior gap, strong prior gap, and strong top gap verify every entry choice.
Its prices and bids are the tuples of Section 1.
All three rows have strict order and entry comparisons, so all four tie-rule pairs give the same unique equilibria.

\begin{table}[htbp]\centering\small
\caption{The selected pair at unchanged B and both controls; computer-assisted.}\label{tab:compb}
\begin{tabular}{@{}lrrrrr@{}}\toprule
Information & $\Ent_0$ & $\Ent_1$ & $\mathsf O_{H,0}$ & $\mathsf O_{H,1}$ & Min. margin\\\midrule
Both private signals &\reg{p4_comp_b_r0_entry}&\reg{p4_comp_b_r1_entry}&\reg{p4_comp_b_r0_own}&\reg{p4_comp_b_r1_own}&\reg{margin_p4_min_signal_comp_b}\\
Investor sees $T$, no $Y$ &\reg{p4_comp_b_no_y_r0_entry}&\reg{p4_comp_b_no_y_r1_entry}&\reg{p4_comp_b_no_y_r0_own}&\reg{p4_comp_b_no_y_r1_own}&\reg{margin_p4_min_comp_b_no_y}\\
Investor sees $J$, no $Y$ &\reg{p4_comp_b_class_r0_entry}&\reg{p4_comp_b_class_r1_entry}&\reg{p4_comp_b_class_r0_own}&\reg{p4_comp_b_class_r1_own}&\reg{margin_p4_min_comp_b_class}\\\bottomrule
\end{tabular}
\end{table}

The earlier alternative $(a_T,d)=(\reg{p4_input_valid_a},\reg{p4_input_valid_d})$ also satisfies the reversal theorem.
Its entry remains $\reg{p4_valid_b_r0_entry}\to\reg{p4_valid_b_r1_entry}$ and ownership remains
$\reg{p4_valid_b_r0_own}\to\reg{p4_valid_b_r1_own}$.
It does not require the challenger signal.
Without $Y$, its entry is $\reg{p4_valid_b_no_y_r0_entry}\to\reg{p4_valid_b_no_y_r1_entry}$ and ownership is
$\reg{p4_valid_b_no_y_r0_own}\to\reg{p4_valid_b_no_y_r1_own}$.
The public exclusion margin is negative: $c_H-B_{r_1}(\mu_+)=\reg{margin_p4_valid_b_no_y_top_exclusion}$.
The strict control margins in the registry force zero weak orders and full strong orders.
Thus this pair shows a reversal, while the selected pair establishes complementarity.

The four sufficient conditions cannot hold with $d\ge a_T$ at unchanged B.
To prove this, put $z=2a_T-1$.
The strong bound requires
\[
 z>z_*:=\frac{k}{m\rho\Delta_1}=\reg{p4_b_signal_z_min}.
\]
If $d\ge a_T$, then $2d-1\ge z$.
The weak gross-return bound consequently exceeds
\[
 z_*[\Delta_0+(1-\rho)z_*w_H(r_0)]
 =k+\reg{margin_p4_b_d_ge_a_impossible}>k.
\]
This contradicts the weak inequality in~\eqref{eq:signalconditions}.
It is a restriction on this sufficient region; the exact mixed comparison above has its own proof.

\section{Parameter sets, strict margins, and the certificate}

\subsection{Open parameter sets}

The equalities that define a signal law, a cost support, or a uniform incumbent are model definitions.
The table states the strict inequalities that define the economic parameter regions.
At $\beta<1/3$, every region is open in its finite primitive coordinates.
The proofs also admit the structural boundary $\beta=1/3$.
For general class laws, the same statement applies to perturbations that preserve the displayed strict moment inequalities.

\begin{table}[htbp]\centering\small
\caption{Admissible parameter regions; analytical.}
\begin{tabular}{@{}p{0.23\linewidth}p{0.70\linewidth}@{}}\toprule
Result & Strict economic inequalities\\\midrule
Atomic controls & $\Delta_0<k<\rho\kappa\Delta_1$, $B_0(\pi)<c_H$, and $B_1(\pi)<c_H< B_1(M)$; the frozen weak row also uses $c_H<B_0(M)$.\\
Cost reversal & $L<B_1(m)$, $B_0(\pi)<c_-<c_+<B_1(M)$, and $\Delta_0<k<\rho\kappa\Delta_1$.\\
Cost return & $L<B_2(m)$, $B_2(M)<c_-$, and $k<\rho\kappa\Delta_2$.\\
Signal reversal & The four inequalities in~\eqref{eq:signalconditions}; the one-pair formula also uses~\eqref{eq:onlypair}.\\
Signal complementarity & The signal-reversal conditions and $B_{r_1}(\mu_+)<c_H$.\\
Ownership & Equation~\eqref{eq:owncondition}; $h>r_1$ makes it automatic for the point masses.\\\bottomrule
\end{tabular}
\end{table}

The common primitive domain is $0<\pi,\rho,\lambda<1$, $0<\beta<1/3$, $k>0$, $c_H>0$, and $0<p<\ell<r_0<r_1<h$.
The cost example also has $L>0$ and $c_+>c_->0$.
The signal domain adds $1/2<a_T,d<1$.
Table~\ref{tab:primitive} records all boundary distances used by these declarations.

\begin{longtable}{@{}p{0.60\linewidth}r@{}}
\caption{Primitive boundary distances; exact or computer-assisted.}\label{tab:primitive}\\
\toprule Quantity & Margin\\\midrule\endfirsthead
\toprule Quantity & Margin\\\midrule\endhead
$\beta$; $1/3-\beta$ & \reg{margin_p4_domain_beta_lower}; \reg{margin_p4_domain_beta_upper}\\
$\lambda$; $1-\lambda$ & \reg{margin_p4_domain_lambda_lower}; \reg{margin_p4_domain_lambda_upper}\\
$\pi$; $1-\pi$ & \reg{margin_p4_domain_pi_lower}; \reg{margin_p4_domain_pi_upper}\\
$\rho$; $1-\rho$ at B & \reg{margin_p4_domain_rho_lower}; \reg{margin_p4_domain_rho_upper}\\
$p$; $\ell-p$ & \reg{margin_p4_domain_reserve}; \reg{margin_p4_domain_ell_reserve}\\
$r_0-\ell$; $r_1-r_0$ at B & \reg{margin_p4_domain_r0_ell}; \reg{margin_p4_domain_r1_r0}\\
$r_2-r_1$; $h-r_2$ at B & \reg{margin_p4_domain_r2_r1}; \reg{margin_p4_domain_h_r2}\\
$c_H$; $k$ at B & \reg{margin_p4_domain_cost}; \reg{margin_p4_domain_fee}\\
$\varepsilon$; $c_+-c_-$ & \reg{margin_p4_domain_epsilon}; \reg{margin_p4_domain_band_width}\\
$c_-$ & \reg{margin_p4_domain_band_lower}\\
$\rho$; $1-\rho$ in S(B) & \reg{margin_p4_domain_s_rho_lower}; \reg{margin_p4_domain_s_rho_upper}\\
$r_0-\ell$; $r_1-r_0$ in S(B) & \reg{margin_p4_domain_s_r0_ell}; \reg{margin_p4_domain_s_r1_r0}\\
$h-r_1$ in S(B) & \reg{margin_p4_domain_s_h_r1}\\
$c_H$; $k$ in S(B) & \reg{margin_p4_domain_s_cost}; \reg{margin_p4_domain_s_fee}\\
$a_T-1/2$; $1-a_T$ at the earlier pair & \reg{margin_p4_domain_signal_a_lower}; \reg{margin_p4_domain_signal_a_upper}\\
$d-1/2$; $1-d$ at the earlier pair & \reg{margin_p4_domain_signal_d_lower}; \reg{margin_p4_domain_signal_d_upper}\\
$d-a_T$ at the earlier pair & \reg{margin_p4_domain_signal_d_a}\\
$a_T-1/2$; $1-a_T$ at the earlier alternative & \reg{margin_p4_domain_valid_a_lower}; \reg{margin_p4_domain_valid_a_upper}\\
$d-1/2$; $1-d$ at the earlier alternative & \reg{margin_p4_domain_valid_d_lower}; \reg{margin_p4_domain_valid_d_upper}\\
$a_T-1/2$; $1-a_T$ at the selected B pair & \reg{margin_p4_domain_comp_a_lower}; \reg{margin_p4_domain_comp_a_upper}\\
$d-1/2$; $1-d$ at the selected B pair & \reg{margin_p4_domain_comp_d_lower}; \reg{margin_p4_domain_comp_d_upper}\\
\bottomrule
\end{longtable}

\begin{longtable}{@{}p{0.65\linewidth}r@{}}
\caption{R5 and R7: strict hypothesis margins at B; computer-assisted where they depend on $e$.}\\
\toprule Quantity & Margin\\\midrule\endfirsthead
\toprule Quantity & Margin\\\midrule\endhead
$k-\Delta_0$ & \reg{margin_p4_weak_trade_B}\\
$\rho m\Delta_1-k$ & \reg{margin_p4_strong_trade_B}\\
$\rho m\Delta_2-k$ & \reg{margin_p4_very_strong_trade_B}\\
$c_H-B_0(1/2)$ & \reg{margin_p4_weak_prior_B}\\
$c_H-B_1(1/2)$ & \reg{margin_p4_strong_prior_B}\\
$B_0(M)-c_H$ & \reg{margin_p4_weak_top_B}\\
$B_1(M)-c_H$ & \reg{margin_p4_strong_top_B}\\
$c_H-B_2(M)$ & \reg{margin_p4_return_B}\\
$B_1(m)-L$, with $L=1/5$ & \reg{margin_p4_atomless_floor_B}\\
$c_--B_0(1/2)$ & \reg{margin_p4_atomless_weak_B}\\
$B_1(M)-c_+$ & \reg{margin_p4_atomless_strong_B}\\
$B_2(m)-L$ & \reg{margin_p4_atomless_floor_r2_B}\\
$c_--B_2(M)$: the declared width fails this condition & \reg{margin_p4_atomless_return_B}\\
$c_+-B_2(M)$ & \reg{margin_p4_atomless_top_upper_r2_B}\\
$B_2(M)-c_-$ & \reg{margin_p4_atomless_top_lower_r2_B}\\
$c_H-0.001-B_2(M)$ & \reg{margin_p4_atomless_small_width_return}\\
$\bar e_H-\rho$, with $W_H(F_0)=W_H(F_1)=1$ & \reg{margin_p4_ownership_rise_B}\\
$\Delta_0$; $\Delta_1$; $\Delta_2$ & \reg{p4_r0_D}; \reg{p4_r1_D}; \reg{p4_r2_D}\\
$g_H(F_0)-g_L(F_0)$ & \reg{margin_p4_gross_gap_r0}\\
$g_H(F_1)-g_L(F_1)$ & \reg{margin_p4_gross_gap_r1}\\
$g_H(F_2)-g_L(F_2)$ & \reg{margin_p4_gross_gap_r2}\\
$w_L(F_0)$; $w_L(F_1)$; $w_L(F_2)$ & \reg{p4_r0_wl}; \reg{p4_r1_wl}; \reg{p4_r2_wl}\\
\bottomrule
\end{longtable}

\begin{table}[htbp]\centering\small
\caption{R8 margins in three declarations; computer-assisted.}
\begin{tabular}{@{}lrrr@{}}\toprule
Margin & S(B), earlier pair & B, earlier pair & B, selected pair\\\midrule
$c_H-B_0(d)$ &\reg{margin_p4_signal_s_private}&\reg{margin_p4_signal_b_private}&\reg{margin_p4_signal_comp_b_private}\\
$B_1(\phi_+(\mu_+))-c_H$ &\reg{margin_p4_signal_s_joint}&\reg{margin_p4_signal_b_joint}&\reg{margin_p4_signal_comp_b_joint}\\
$k-(2a_T-1)D^+_0$ &\reg{margin_p4_signal_s_weak}&\reg{margin_p4_signal_b_weak}&\reg{margin_p4_signal_comp_b_weak}\\
$(2a_T-1)m\rho\Delta_1-k$ &\reg{margin_p4_signal_s_strong}&\reg{margin_p4_signal_b_strong}&\reg{margin_p4_signal_comp_b_strong}\\
$c_H-B_1(\phi_-(\mu_+))$ &\reg{margin_p4_signal_s_other_pair}&\reg{margin_p4_signal_b_other_pair}&\reg{margin_p4_signal_comp_b_other_pair}\\\bottomrule
\end{tabular}
\end{table}

\begin{longtable}{@{}p{0.65\linewidth}r@{}}
\caption{Strict margins for the actual B signal equilibria; computer-assisted.}\\
\toprule Quantity & Margin\\\midrule\endfirsthead
\toprule Quantity & Margin\\\midrule\endhead
$B_0(\phi_+(\mu_-))-c_H$ & \reg{margin_p4_b_weak_all_yplus}\\
$c_H-B_0(\phi_-(\mu_+))$ & \reg{margin_p4_b_weak_no_yminus}\\
$(2a_T-1)(311/1920)/2-k$ & \reg{margin_p4_b_weak_zero_return}\\
$k-(2a_T-1)m(311/1920)$ & \reg{margin_p4_b_weak_full_loss}\\
$u-119/96$ & \reg{margin_p4_b_weak_rate_room}\\
$(2a_T-1)m(137/192)-k$ & \reg{margin_p4_b_strong_buy}\\
$B_1(d)-c_H$ & \reg{margin_p4_b_strong_prior_yplus}\\
$c_H-B_1(\phi_-(\mu_+))$ & \reg{margin_p4_b_strong_no_yminus}\\
$19/795-k$ & \reg{margin_p4_b_strong_threshold_floor_return}\\
$u-4/3$ & \reg{margin_p4_b_strong_rate_room}\\
$47/106-21/50$ & \reg{margin_p4_b_strong_sale_public_gap}\\
$B_1(\phi_+(\mu_0))-c_H$ at the mixed tuple & \reg{margin_p4_b_strong_zero_yplus}\\
$c_H-B_0(\mu_+)$ without $Y$ & \reg{margin_p4_b_no_y_top_exclusion_r0}\\
$c_H-B_1(\mu_+)$ without $Y$ & \reg{margin_p4_b_no_y_top_exclusion_r1}\\
$k-\rho(2a_T-1)\Delta_0/2$ without $Y$ & \reg{margin_p4_b_no_y_weak}\\
$\rho(2a_T-1)\Delta_1/2-k$ without $Y$ & \reg{margin_p4_b_no_y_zero_gain}\\
$k-\rho(2a_T-1)m\Delta_1$ without $Y$ & \reg{margin_p4_b_no_y_full_loss}\\
\bottomrule
\end{longtable}

\begin{table}[htbp]\centering\small
\caption{Strict control margins for the selected B pair; computer-assisted.}
\begin{tabular}{@{}lr@{}}\toprule
Quantity & Margin\\\midrule
$k-(2a_T-1)\Delta_0$, investor sees $T$, no $Y$ & \reg{margin_p4_comp_b_no_y_weak}\\
$(2a_T-1)\rho m\Delta_1-k$, investor sees $T$, no $Y$ & \reg{margin_p4_comp_b_no_y_strong}\\
$c_H-B_0(1/2)$, investor sees $T$, no $Y$ & \reg{margin_p4_comp_b_no_y_prior_r0}\\
$c_H-B_1(\mu_+)$, investor sees $T$, no $Y$ & \reg{margin_p4_comp_b_no_y_top_r1}\\
$k-\Delta_0$, investor sees $J$, no $Y$ & \reg{margin_p4_comp_b_class_weak}\\
$\rho m\Delta_1-k$, investor sees $J$, no $Y$ & \reg{margin_p4_comp_b_class_strong}\\
$c_H-B_0(1/2)$, investor sees $J$, no $Y$ & \reg{margin_p4_comp_b_class_prior_r0}\\
$c_H-B_1(1/2)$, investor sees $J$, no $Y$ & \reg{margin_p4_comp_b_class_prior_r1}\\
$B_1(M)-c_H$, investor sees $J$, no $Y$ & \reg{margin_p4_comp_b_class_top_r1}\\\bottomrule
\end{tabular}
\end{table}

\clearpage
\begin{table}[htbp]\centering\small
\caption{Strict margins for the information controls in S(B); computer-assisted.}
\begin{tabular}{@{}lr@{}}\toprule
Quantity & Margin\\\midrule
$k-(2a_T-1)\Delta_0$, investor sees $T$ & \reg{margin_p4_s_no_y_weak}\\
$(2a_T-1)\rho m\Delta_1-k$, investor sees $T$ & \reg{margin_p4_s_no_y_strong}\\
$c_H-B_0(\mu_+)$, investor sees $T$ & \reg{margin_p4_s_no_y_top_r0}\\
$c_H-B_1(\mu_+)$, investor sees $T$ & \reg{margin_p4_s_no_y_top_r1}\\
$k-\Delta_0$, investor sees $J$ & \reg{margin_p4_s_class_weak}\\
$\rho m\Delta_1-k$, investor sees $J$ & \reg{margin_p4_s_class_strong}\\
$c_H-B_0(1/2)$, investor sees $J$ & \reg{margin_p4_s_class_prior_r0}\\
$c_H-B_1(M)$, investor sees $J$ & \reg{margin_p4_s_class_top_r1}\\\bottomrule
\end{tabular}
\end{table}

\begin{table}[htbp]\centering\footnotesize
\caption{The four sufficient margins at the six signal pairs; computer-assisted.}
\begin{tabular}{@{}rrrrrr@{}}\toprule
$a_T$ & $d$ & Private exclusion & Joint entry & Weak order & Strong order\\\midrule
\reg{p4_accuracy_69}&\reg{p4_accuracy_75}&\reg{margin_p4_grid_69_75_private}&\reg{margin_p4_grid_69_75_joint}&\reg{margin_p4_grid_69_75_weak}&\reg{margin_p4_grid_69_75_strong}\\
\reg{p4_accuracy_70}&\reg{p4_accuracy_75}&\reg{margin_p4_grid_70_75_private}&\reg{margin_p4_grid_70_75_joint}&\reg{margin_p4_grid_70_75_weak}&\reg{margin_p4_grid_70_75_strong}\\
\reg{p4_accuracy_71}&\reg{p4_accuracy_74}&\reg{margin_p4_grid_71_74_private}&\reg{margin_p4_grid_71_74_joint}&\reg{margin_p4_grid_71_74_weak}&\reg{margin_p4_grid_71_74_strong}\\
\reg{p4_accuracy_71}&\reg{p4_accuracy_75}&\reg{margin_p4_grid_71_75_private}&\reg{margin_p4_grid_71_75_joint}&\reg{margin_p4_grid_71_75_weak}&\reg{margin_p4_grid_71_75_strong}\\
\reg{p4_accuracy_72}&\reg{p4_accuracy_74}&\reg{margin_p4_grid_72_74_private}&\reg{margin_p4_grid_72_74_joint}&\reg{margin_p4_grid_72_74_weak}&\reg{margin_p4_grid_72_74_strong}\\
\reg{p4_accuracy_72}&\reg{p4_accuracy_75}&\reg{margin_p4_grid_72_75_private}&\reg{margin_p4_grid_72_75_joint}&\reg{margin_p4_grid_72_75_weak}&\reg{margin_p4_grid_72_75_strong}\\\bottomrule
\end{tabular}
\end{table}

\subsection{The arithmetic record}

The certificate uses exact rational intervals.
For $N=60$, it encloses the exponential constant by
\[
 \sum_{n=0}^{N}\frac1{n!}<e<
 \sum_{n=0}^{N}\frac1{n!}+\frac1{N\,N!}.
\]
The remainder bound follows from a geometric series with ratio $1/(N+1)$.
Each arithmetic operation propagates rational endpoints.
A strict sign uses the lower or upper endpoint, without a floating-point tolerance.
Decimal log endpoints round outwards.

The executable derivation is \path{scripts/P4_information/derive.py}.
Its output is \path{scripts/P4_information/run.log}.
The file \texttt{endpoints.json} stores the exact rational endpoints for every registry row.
The file \texttt{receipt.json} records the primitive declaration, method, and output hashes.
Every displayed quantitative output comes from \path{registry/P4_information.csv}.

The derivation prints every order posterior, every price, both class entry rates, and all six investor payoffs for each tuple.
It also prints every entry gap for each private signal at every order.
The accuracy grid records all four sufficient margins and the one-pair margin for every pair.
The finite computations certify the stated signs and values; the proofs above establish the profile bounds and uniqueness.

The B declaration uses $\beta=\reg{beta_B}$ exactly and $\lambda=u\beta/(1+u\beta)$ without rounding.
Its strong ownership display is $\reg{p4_base_r1_own}$.
The earlier exact-entry-match convention used a different $\beta$ and gave a different sixth decimal.
At $r_2$, the current atomless entry display is $\reg{p4_atomless_r2_entry}$.
The signal displays remain unchanged at their declared S(B) vector.

\end{document}
